%%%%%%%%%%%%%%%%%%%%%%%%%%%%%%%%%%%%%%%%%%%%%%%%%%%%%%%%%%%%%%%%%%%%%%%%%%%%
%% Author template for Operations Reseacrh (opre) for articles with no e-companion (EC)
%% Mirko Janc, Ph.D., INFORMS, mirko.janc@informs.org
%% ver. 0.95, December 2010
%%%%%%%%%%%%%%%%%%%%%%%%%%%%%%%%%%%%%%%%%%%%%%%%%%%%%%%%%%%%%%%%%%%%%%%%%%%%
%\documentclass[opre,blindrev]{informs3}
\RequirePackage{rotating} % Need this before \documentclass or sidewaystables aren't centered
\documentclass[ijoc,noblindrev]{informs3} % current current default for manuscript submission

\OneAndAHalfSpacedXI % Made default 4/4/2014 at request
%%\OneAndAHalfSpacedXI % current default line spacing
%%\OneAndAHalfSpacedXII
%%\DoubleSpacedXII

% If hyperref is used, dvi-to-ps driver of choice must be declared as
%   an additional option to the \documentclass. For example
%\documentclass[dvips,opre]{informs3}      % if dvips is used
%\documentclass[dvipsone,opre]{informs3}   % if dvipsone is used, etc.

%%% OPRE uses endnotes. If you do not use them, put a percent sign before
%%% the \theendnotes command. This template does show how to use them.
\usepackage{endnotes}
\let\footnote=\endnote
\let\enotesize=\normalsize
\def\notesname{Endnotes}%
\def\makeenmark{$^{\theenmark}$}
\def\enoteformat{\rightskip0pt\leftskip0pt\parindent=1.75em
  \leavevmode\llap{\theenmark.\enskip}}

% Private macros here (check that there is no clash with the style)

% Natbib setup for author-year style
\usepackage{natbib}
 \bibpunct[, ]{(}{)}{,}{a}{}{,}%
 \def\bibfont{\footnotesize}%
 \def\bibsep{\smallskipamount}%
 \def\bibhang{10pt}%
 \def\newblock{\ }%
 \def\BIBand{and}%


% \renewcommand{\baselinestretch}{1.3}
\usepackage[utf8]{inputenc}

\usepackage[colorinlistoftodos,prependcaption,textsize=small]{todonotes}
\usepackage{rotating}
\usepackage[shortlabels]{enumitem}
\usepackage{placeins}
\usepackage{mathtools}
\usepackage[noend]{algorithm2e}
\usepackage[english]{babel}
\usepackage[margin=1in]{geometry}
\usepackage{bbding}
\usepackage[shortlabels]{enumitem}
%\usepackage{{appendix}}
\usepackage[super]{nth}
\usepackage{bm, bbm}
\usepackage{amsmath, amssymb}
\usepackage{mathtools}
\usepackage{subcaption, placeins}
\usepackage{graphicx}
\usepackage[hidelinks]{hyperref}
\usepackage{pgfplots}
\usepackage{ulem}
% \usepackage{comment}


% \usepackage{algorithm}
% \usepackage{algorithmic}
\usepackage{subcaption}
\usepackage{booktabs}
\usepackage{multirow} 
\usepackage{placeins}
\usepackage{booktabs}
\usepackage{accents}
\usepackage{graphics}
\newcommand{\N}{\mathbb{N}}
\newcommand{\Z}{\mathbb{Z}}
\DeclarePairedDelimiter\ceil{\lceil}{\rceil}
\DeclarePairedDelimiter\floor{\lfloor}{\rfloor}
\usepackage{tikz}
\usepackage{placeins}
\usepackage{float}
\usepackage{graphicx}
\usepackage{accents}
\usepackage{rotating}
\usepackage{graphics}
%\usepackage{todonotes}

\usepackage{duckuments}

\usepackage{color}
\definecolor{blue}{rgb}{0,0,1}
% \usepackage{natbib}
% \setcitestyle{square, numbers}
% \bibliographystyle{abbrvnat}

% \newtheorem{theorem}{Theorem}
% \newtheorem{definition}{Definition}
% \newtheorem{lemma}{Lemma}
\newtheorem{appendixlemma}{Lemma}[section]
% \newtheorem{proposition}{Proposition}
% \newtheorem{corollary}{Corollary}
% \newtheorem{example}{Example}
% \newtheorem{remark}{Remark}
% \newtheorem{assumption}{Assumption}
\setlist[itemize]{leftmargin=*, nosep}

\newcommand{\argminE}{\mathop{\mathrm{argmin}}}  

\newcommand{\ag}[1]{{\color{black}#1}}


\usepackage{etoolbox}
\newcommand{\zerodisplayskips}{%
  \setlength{\abovedisplayskip}{7pt}%
  \setlength{\belowdisplayskip}{5pt}%
  \setlength{\abovedisplayshortskip}{0pt}%
  \setlength{\belowdisplayshortskip}{0pt}}
\appto{\normalsize}{\zerodisplayskips}
\appto{\small}{\zerodisplayskips}
\appto{\footnotesize}{\zerodisplayskips}
\usepackage{xfrac}

% \DeclareMathOperator*{\argmax}{arg\,max}
% \DeclareMathOperator*{\argmin}{arg\,min}

%% Setup of theorem styles. Outcomment only one.
%% Preferred default is the first option.
\TheoremsNumberedThrough     % Preferred (Theorem 1, Lemma 1, Theorem 2)
%\TheoremsNumberedByChapter  % (Theorem 1.1, Lema 1.1, Theorem 1.2)
\ECRepeatTheorems


%% Setup of the equation numbering system. Outcomment only one.
%% Preferred default is the first option.
\EquationsNumberedThrough    % Default: (1), (2), ...
%\EquationsNumberedBySection % (1.1), (1.2), ...

% For new submissions, leave this number blank.
% For revisions, input the manuscript number assigned by the on-line
% system along with a suffix ".Rx" where x is the revision number.
% \MANUSCRIPTNO{MS-0001-1922.65}

%%%%%%%%%%%%%%%%
\begin{document}
%%%%%%%%%%%%%%%%

% Outcomment only when entries are known. Otherwise leave as is and
%   default values will be used.
%\setcounter{page}{1}
%\VOLUME{00}%
%\NO{0}%
%\MONTH{Xxxxx}% (month or a similar seasonal id)
%\YEAR{0000}% e.g., 2005
%\FIRSTPAGE{000}%
%\LASTPAGE{000}%
%\SHORTYEAR{00}% shortened year (two-digit)
%\ISSUE{0000} %
%\LONGFIRSTPAGE{0001} %
%\DOI{10.1287/xxxx.0000.0000}%

% Author's names for the running heads
% Sample depending on the number of authors;
% \RUNAUTHOR{Jones}
% \RUNAUTHOR{Jones and Wilson}
% \RUNAUTHOR{Jones, Miller, and Wilson}
% \RUNAUTHOR{Jones et al.} % for four or more authors
% Enter authors following the given pattern:
\RUNAUTHOR{Ryan Cory-Wright and Andr\'es G\'omez} % Putting alphabetical ordering, happy to discuss if you feel strongly

% Title or shortened title suitable for running heads. Sample:
% Enter the (shortened) title:
\RUNTITLE{Optimal Cross-Validation for Sparse Linear Regression}

% I'm happy to change this title; I think it's an improvement on what we had previously though
\TITLE{Optimal Cross-Validation\\ for Sparse Linear Regression}

% Block of authors and their affiliations starts here:
% NOTE: Authors with same affiliation, if the order of authors allows,
%   should be entered in ONE field, separated by a comma.
%   \EMAIL field can be repeated if more than one author
\ARTICLEAUTHORS{%
\AUTHOR{Ryan Cory-Wright}
\AFF{Department of Analytics, Marketing and Operations, Imperial College Business School, London, UK\\
% IBM Thomas J. Watson Research Center, USA\\
ORCID: \href{https://orcid.org/0000-0002-4485-0619}{$0000$-$0002$-$4485$-$0619$}\\ \EMAIL{r.cory-wright@imperial.ac.uk}} 
\AUTHOR{Andr{\'e}s G{\'o}mez}
\AFF{Department of Industrial and Systems Engineering, Viterbi School of Engineering, University of Southern California, CA\\
ORCID: \href{https://orcid.org/0000-0003-3668-0653}{$0000$-$0003$-$3668$-$0653$}\\
\EMAIL{gomezand@usc.edu}}
% Enter all authors
} % end of the block

\ABSTRACT{%
Given a high-dimensional covariate matrix and a response vector, ridge-regularized sparse linear regression selects a subset of features that explains the relationship between covariates and the response in an interpretable manner. To select the sparsity and robustness of linear regressors, techniques like $k$-fold cross-validation are commonly used for hyperparameter tuning. However, cross-validation substantially increases the computational cost of sparse regression as it requires solving many mixed-integer optimization problems (MIOs) for each hyperparameter combination. To improve upon this state of affairs, we obtain computationally tractable relaxations of $k$-fold cross-validation metrics, facilitating hyperparameter selection after solving $50$--$80\%$ fewer MIOs in practice. These relaxations result in an efficient cyclic coordinate descent scheme, achieving $10\%$--$30\%$ lower validation errors than via traditional methods such as grid search with MCP or GLMNet across a suite of $13$ real-world datasets. 
}%


\KEYWORDS{Cross-validation; perspective formulation; sparse regression; bi-level convex relaxation} 
%\HISTORY{}

\maketitle
\vspace{-10mm}

\section{Introduction}\label{sec:intro}
Over the past fifteen years, Moore's law has spurred an %exponential 
explosion of high-dimensional datasets for scientific discovery across multiple fields \citep[][]{mcafee2012big}. These datasets often consist of a design matrix $\bm{X} \in \mathbb{R}^{n \times p}$ of explanatory variables and an output vector $\bm{y} \in \mathbb{R}^n$ of response variables. Accordingly, practitioners often aim to explain the response variables linearly via the equation $\bm{y}=\bm{X}\bm{\beta}+\bm{\epsilon}$ for a vector of regression coefficients $\bm{\beta} \in \mathbb{R}^p$, which is to be inferred, and a vector of error $\bm{\epsilon}$, typically kept small by minimizing the least squares (LS) error of the regression. 

%A common method to determine the regressors $\bm{\beta}$ is to compute a least squares (LS) estimator by minimizing the sum of squares error. %, $\Vert \bm{\epsilon}\Vert_2^2$. 
Despite its computational efficiency, LS regression exhibits two practical limitations. First, when $p \gg n$, there is not enough data to accurately infer $\bm{\beta}$ via LS, and LS regression generates estimators which perform poorly out-of-sample due to a data curse of dimensionality \citep[][]{buhlmann2011statistics, gamarnik2022sparse}.
Second, LS regression generically selects every feature, including irrelevant ones. This is a significant challenge when regression coefficients are used for high-stakes decision-making tasks and non-zero coefficients inform the decision making.

To tackle the challenges of dimensionality and false discovery, sparse learning has emerged as a popular methodology for explaining the relationship between inputs {\color{black}$\bm{X}$} and outputs {\color{black}$\bm{y}$}. % in an interpretable way. 
%Perhaps the most 
A popular sparse learning model is ridge-regularized sparse regression, which admits the formulation \citep[][]{bertsimas2020sparse,xie2020scalable, hastie2020best,  atamturk2020safe, kenney2021mip, hazimeh2021sparse, liu2023okridge}:
\begin{align}\label{eqn:l0l2_training}
    \min_{\bm{\beta} \in \mathbb{R}^p} \quad &  \frac{\gamma}{2} \Vert \bm{\beta}\Vert_2^2+\Vert \bm{y}-\bm{X}\bm{\beta}\Vert_2^2  \quad \text{s.t.} \quad \Vert \bm{\beta}\Vert_0 \leq \tau,
\end{align} 
where $\tau \in \{1, \ldots, p\}$ and $\gamma >0$ are hyperparameters that respectively model the sparsity and robustness of the linear model $\bm{\beta}$ \citep[cf.][]{xu2008robust, bertsimas2018characterization}, and we assume %here and throughout the paper 
that $\bm{X}, \bm{y}$ have undergone standard preprocessing so that $\bm{y}$ is a zero-mean vector and $\bm{X}$ has zero-mean unit-variance columns, meaning $\gamma$ penalizes each feature equally.


Problem \eqref{eqn:l0l2_training} is \ag{computationally} challenging {\color{black}(indeed, NP-hard \cite{natarajan1995sparse})} and initial formulations could not scale to problems with thousands of features \citep[][]{hastie2020best}. In a more positive direction, by developing and exploiting tight conic relaxations of appropriate substructures of \eqref{eqn:l0l2_training}, e.g., the perspective relaxation \citep{ceria1999convex, stubbs1999branch, gunluk2010perspective}%(see also \cite{dong2015regularization, atamturk2019rank, han20202x2})
, more recent mixed-integer optimization techniques such as branch-and-bound \citep{hazimeh2021sparse} scale to larger instances with thousands of %datapoints and 
features. {\color{black}We refer to \cite{bertsimas2021unified, atamturk2019rank} for reviews of perspective and related relaxations.}

To be sure, the aforementioned works solve \eqref{eqn:l0l2_training} rapidly. Unfortunately, they do not address arguably the most significant difficulty {\color{black}in performing sparse regression}. The hyperparameters $(\tau, \gamma)$ 
{\color{black}are not known to the decision-maker ahead of time, as is often assumed in the literature for convenience.}
%cannot be elicited from an oracle, as is often assumed in the literature. 
Rather, they must be selected by {\color{black}the decision-maker}, which is potentially much more challenging than solving \eqref{eqn:l0l2_training} for a single value of $(\tau, \gamma)$ \citep[]{hansen1992new}. Indeed, selecting $(\tau, \gamma)$ typically involves minimizing a validation metric over a grid of values, which is computationally expensive \citep[]{larochelle2007empirical}. 

{\color{black} Perhaps the most popular validation metric is hold-out \citep{hastie2009elements}, where one omits a portion of the data when training the model and then evaluates performance on this hold-out set as a proxy for the model's test set performance. However, hold-out validation is sometimes called a high-variance approach \citep{hastie2009elements}, because the validation score can vary significantly depending on the hold-out set selected. 

To reduce the variance in this procedure, a number of authors have proposed:
}
\paragraph{The Cross-Validation Paradigm: }
To obtain accurate models that generalize well to unseen data, cross-validation has emerged as a popular model selection paradigm. {\color{black}Early iterations of this paradigm, as reviewed by} \citet{stone1978cross}, {\color{black}suggest solving \eqref{eqn:l0l2_training} with the $i$th data point removed for each $i \in [n]$, and estimating the out-of-sample performance of a solution to Problem \eqref{eqn:l0l2_training} via the average performance of the $n$ estimators with the $i$th training data point removed, on the $i$th data point. This approach is known as leave-one-out cross-validation (LOOCV). 

A popular variant of LOOCV, known as $k$-fold cross-validation, comprises removing subsets of $n/k$ data points at a time and breaking the data into $k$ folds in total, which significantly reduces the computational burden of cross-validation while having less variance than a hold-out approach \citep{burman1989comparative, arlot2010survey}. However, even $k$-fold cross-validation may be prohibitive in the case of MIOs such as \eqref{eqn:l0l2_training}. Indeed, as identified by \cite{hastie2020best}, with a time limit of 3 minutes per MIO, using 10-fold cross-validation to choose between subset sizes $\tau=0,\dots,50$ in an instance of Problem \eqref{eqn:l0l2_training} with $p=100$ and $n=500$ requires 25 hours of computational time.} 


For sparse regression, {\color{black}given a partition $\mathcal{N}_1,\dots, \mathcal{N}_k$ of $[n]$,}  performing {\color{black}$k$-fold cross-validation} corresponds to selecting hyperparameters $\gamma, \tau$ which minimize the function:
{\color{black}
\begin{equation}\label{prob:upperlevelofv_validation}
	h(\gamma,\tau) = \frac{1}{k}\sum_{j=1}^k\sum_{i\in \mathcal{N}_j} (y_i -\bm{x_i}^\top \bm{\beta}^{(\mathcal{N}_j)}(\gamma,\tau))^2
\end{equation}{\color{black}
where $\bm{\beta}^{(\mathcal{N}_j)}(\gamma,\tau)$ denotes an optimal solution to the following lower-level problem for any $\mathcal{N}_j$:
\begin{equation}
	\bm{\beta}^{(\mathcal{N}_j)}(\gamma,\tau) \in \argmin_{\bm{\beta} \in \mathbb{R}^p} \ \frac{\gamma}{2}\Vert \bm{\beta}\Vert_2^2 +\Vert \bm{y}^{(\mathcal{N}_j)}-\bm{X}^{(\mathcal{N}_j)}\bm{\beta}\Vert_2^2\quad \text{s.t.}\quad \Vert \bm{\beta}\Vert_0 \leq \tau, \label{eqn:lowerlevel1}
\end{equation}}
$\gamma>0$ is a hyperparameter, $\tau$ is a sparsity budget, $\bm{X}^{(\mathcal{N}_j)}, \bm{y}^{(\mathcal{N}_j)}$ denotes the dataset with the data in $\mathcal{N}_j$ removed, and we take $\bm{\beta}^{(\mathcal{N}_j)}(\gamma, \tau)$ to be unique for a given $\tau, \gamma$ for convenience}\footnote{This assumption seems plausible, as the training objective is strongly convex for a fixed binary support vector, and therefore for each binary support vector there is indeed a unique solution. One could relax this assumption by defining $h(\gamma, \tau)$ to be the minimum {\color{black}cross-validation} error over all training-optimal solutions $\bm{\beta}^{(i)}$, as is commonly done in the bilevel optimization literature, giving what is called an optimistic formulation of a bilevel problem \citep[see][for a review]{beck2021gentle}. However, this would make the {\color{black}cross-validation} error less tractable.}. In words, $h(\gamma, \tau)$ denotes the average prediction error on each left-out fold for a sparse regressor with hyperparameters $(\gamma, \tau)$ trained on the remaining folds.


{\color{black}We remark that if all sets $\mathcal{N}_j$ are taken to be singletons and $k=n$, minimizing $h$ corresponds to LOOCV. Moreover, if $k=2$ and the term with $j=2$ is removed from $h$, optimizing $h$ reduces to minimizing the hold-out error.}
After selecting $(\gamma, \tau)$, practitioners usually train a final model on the entire dataset, by solving Problem \eqref{eqn:l0l2_training} with the selected hyperparameter combination. %To ensure that $\gamma$ has the same impact in the final model as the cross-validated models, they sometimes first multiply $\gamma$ by the bias correction term $n/(n-n/k)$ \citep[see][for a justification]{liu2019ridge}\footnote{We remark that applying this bias correction term is equivalent to normalizing the least squares error $\Vert \bm{X}\bm{\beta}-\bm{Y}\Vert_2^2$ in the training problem, by dividing this term by the number of data points $n$ (or $n-n/k$).}. 




\paragraph{Our Approach:} We propose techniques for obtaining strong bounds on validation metrics in polynomial time and leverage these bounds to design algorithms for minimizing the {\color{black}cross-validation} error in Sections \ref{sec:lowerbounds}-\ref{sec:corddescent}. By performing a perturbation analysis of perspective relaxations of sparse regression problems, we construct convex relaxations of the $k$-fold cross-validation error, which allows us to minimize it without explicitly solving MIOs at each data fold and for each hyperparameter combination. This results in a branch-and-bound algorithm for hyperparameter selection that is substantially more efficient than state-of-the-art techniques like grid search. {\color{black}As an aside, we remark that as cross-validation is more general than hold-out validation, our convex relaxations can be generalized immediately to the hold-out case.}

In numerical experiments (Section \ref{sec:numres}), we assess the impact of our contributions. We observe on real datasets that our branch-and-bound scheme reduces the number of MIOs that need to be solved by an average of $50\%$--$70\%$. Further, we leverage our branch-and-bound scheme to design a cyclic coordinate scheme that iteratively minimizes $\tau$ and $\gamma$. We observe that on real datasets, our scheme reduces the $k$-fold cross-validation error by an average of $10\%$--$40\%$ compared to state-of-the-art methods like MCP and GLMNet.




\subsection{Literature Review}
Our work falls at the intersection of three areas of the optimization literature. %typically considered in isolation
First, hyperparameter selection techniques for optimizing the performance of a machine learning model by selecting hyperparameters that perform well on a validation set. Second, bilevel approaches that reformulate and solve hyperparameter selection problems as bilevel problems. Finally, perspective reformulation techniques for mixed-integer problems with logical constraints{\color{black}, as discussed above.} To put our contributions into context, we now review {\color{black}the two remaining} areas of the literature.

\paragraph{Hyperparameter Selection Techniques for Machine Learning Problems:} 
A wide variety of hyperparameter selection techniques have been proposed for machine learning problems such as sparse regression, 
{\color{black}including grid search \citep[][]{larochelle2007empirical} as reviewed in Section \ref{sec:intro},}
%or deep learning, 
%which follow the same meta-approach: given a set of hyperparameters $\mathcal{L}:=\{(\gamma, \tau)\}$, we select a combination $(\gamma^\star, \tau^\star)$ which approximately minimize a validation metric $h(\gamma, \tau)$. Perhaps the most popular approach within this paradigm is grid search \citep[][]{larochelle2007empirical}, wherein we let $\mathcal{L}:=\{\gamma_1, \gamma_2, \ldots, \gamma_L\} \times \{\tau_1, \ldots, \tau_L\}$ and minimize a validation metric by evaluating it for each $(\gamma, \tau) \in \mathcal{L}$ separately. Unfortunately, this involves solving a number of training problems exponential in the number of hyperparameters to obtain a solution, i.e., suffers from a curse of dimensionality.% \citep[][]{bellman1961adaptive}. 
%\ag{[I am thinking that this paragraph can be deleted, we have already mentioned all of this.]}
%Originally proposed to break this curse of dimensionality, 
{\color{black}and random search \citep[cf.][]{bergstra2012random}.} %has since emerged as a viable competitor to grid search. 
In random search, we let $\mathcal{L}$ be a random sample from a space of valid hyperparameters, e.g., a uniform distribution over $[10^{-3}, 10^3] \times [p]$ for sparse regression. Remarkably, in settings with many hyperparameters, random search usually outperforms grid search for a given budget on the number of training problems that can be solved, because validation functions often have a lower effective dimension than the number of hyperparameters present in the model \citep[][]{bergstra2012random}. %, and grid search fails to recognize and exploit this structure. 
However, grid search remains competitive for problems with a small number of hyperparameters, such as sparse regression. %Correspondingly, both methods are now implemented in popular hyperparameter selection toolkits such as \verb|Keras| and \verb|Scikit.learn|. 

The modern era of hyperparameter selection strategies was ushered in by the increasing prominence of deep learning methods in applications from voice recognition to drug discovery \citep[see][for a review]{lecun2015deep}. 
The volume of data available and number of hyperparameters needing selection challenged the aforementioned methods and led to new techniques, including evolutionary strategies%\citep{beyer2002evolution}
, Bayesian optimization techniques \citep{ frazier2018tutorial} and bandit methods \cite{falkner2018bohb}. 
% %\citep[see][for a review]{feurer2019hyperparameter}
However, in sparse regression problems where we aim to optimize two hyperparameters, these methods are isomorphic to grid or random search. Further, none of these approaches provide locally optimal hyperparameter combinations with respect to a LOOCV metric, which suggests there is room for improvement upon the state-of-the-art in sparse regression. 

{We point out that current approaches for hyperparameter selection are similar to existing methods for multi-objective mixed-integer optimization. While there has been recent progress in improving multi-objective algorithms for mixed-integer linear programs \citep{lokman2013finding,stidsen2014branch}, %soylu2018search, belotti2016fathoming
a direct application of these methods might be unnecessarily expensive. Indeed, these approaches seek to compute the efficient frontier \citep{boland2015criterion,boland2015criterion2} (i.e., solve problems for all possible values of the regularization parameter), whereas we are interested in only the combination of parameters that optimize a well-defined metric (e.g., \ag{\color{black}the cross-validation error}). %As we show in this work, the leave-one-out information can be exploited to speed up the algorithms substantially.
}




\paragraph{Bilevel Optimization for Hyperparameter Selection:} In a complementary direction, several authors have proposed selecting hyperparameters via bilevel optimization \citep[see][for a general theory]{beck2021gentle}, since \cite{bennett2006model} %kunapuli2008classification
recognized that cross-validation is a special case of bilevel optimization. 
%Indeed, cross-validation is the problem of minimizing a validation metric (an upper-level objective in the parlance of bilevel optimization) with respect to a model which is the optimal solution of a training problem such as Problem \eqref{eqn:l0l2_training} (the solution to a lower-level problem in the parlance of bilevel optimization). 
Therefore, {\color{black}in principle, we could} minimize the {\color{black}cross-validation} error in sparse regression by invoking bilevel techniques. Unfortunately, this approach seems intractable in both theory and practice \citep[][]{ben1990computational, hansen1992new}. Indeed, standard bilevel approaches such as dualizing the lower-level problem are challenging to apply in our context because our lower-level problems are non-convex and cannot easily be dualized.

Although slow in its original implementation, several authors have proposed making hyperparameter optimization more tractable by combining bilevel optimization with tractable modeling paradigms to obtain locally optimal sets of hyperparameters. Among others, \cite{sinha2020gradient} recommends taking a gradient-based approximation of the lower-level problem and thereby reducing the bilevel problem to a single-level problem, \cite{okuno2021lp} advocates selecting hyperparameters by solving the KKT conditions of a bilevel problem, and \cite{ye2022difference} proposes solving bilevel hyperparameter problems via difference-of-convex methods to obtain a stationary point.% of the upper-level objective. 

Specializing our review to regression, three works aim to optimize the performance of regression models on a validation metric. First, \cite{takano2020best} proposes optimizing the $k$-fold validation loss, assuming all folds share the same support. Unfortunately, although their assumption improves their method's tractability, it may lead to subpar statistical performance. Second, \citet{stephenson2021can} proposes first-order methods for minimizing the leave-one-out error in ridge regression problems (without sparsity constraints). However, it is unclear how to generalize their approach to settings with sparsity constraints. Finally, perhaps closest to our work, \cite{kenney2021mip} proposes a bisection algorithm for selecting the optimal sparsity parameter in a sparse regression problem by approximately minimizing the $k$-fold cross-validation error. It is, however, worth noting that this approach is not guaranteed to converge to an optimal sparsity parameter with respect to the $k$-fold error, because it does not develop lower bounds on the $k$-fold error. 

\subsection{Structure}

The rest of the paper is laid out as follows:
\begin{itemize}\color{black}
\item In Section \ref{sec:lowerbounds}, we observe that validation metrics are potentially expensive to evaluate, because they involve solving up to {\color{black}$k+1$} MIOs (in the $k$-fold case), and accordingly develop tractable lower and upper bounds that can be computed without solving any MIOs.
\item In Section \ref{sec:corddescent}, we propose an efficient coordinate descent scheme for identifying locally optimal hyperparameters with respect to the validation error. Specifically, in Section \ref{ssec:parametric}, we develop an efficient scheme for minimizing the {\color{black}cross-validation} error with respect to $\tau$, and in Section \ref{ssec:parametric2}, we propose a scheme for optimizing with respect to $\gamma$.
\item In Section \ref{sec:numres}, we benchmark our proposed approaches on both synthetic and real datasets. On synthetic datasets, we find that optimization significantly reduces the validation error compared to grid search approaches, especially in underdetermined settings. Moreover, the proposed approach leads to a 50-80\% improvement in the number of MIOs solved compared to standard grid search techniques, without sacrificing solution quality.
\end{itemize}

\subsection*{Notation} We let non-boldface characters such as $b$ denote scalars, lowercase bold-faced characters such as $\bm{x}$ denote vectors, uppercase bold-faced characters such as $\bm{A}$ denote matrices, and calligraphic uppercase characters such as $\mathcal{Z}$ denote sets. We let $[n]$ denote the running set of indices $\{1, \dots, n\}$, and $\Vert \bm{x}\Vert_0:=|\{j: x_j \neq 0\}|$ denote the $\ell_0$ pseudo-norm, i.e., the number of non-zero entries in $\bm{x}$. Finally, we let $\bm{e}$ denote the vector of ones, and $\bm{0}$ denote the vector of all zeros. 

%We also invoke matrix operators. We let $\langle \cdot,\cdot \rangle$ denote the Frobenius inner product between two matrices of the same dimension, and $\bm{X}^\dag$ denote the Moore-Penrose pseudoinverse of a matrix $\bm{X}$. If $\bm{x} \neq \bm{0}$ is a vector, then $\bm{x}^\dag=\frac{\bm{x}^\top}{\bm{x}^\top \bm{x}}$; see \citet{johnson1985matrix} for an overview of matrix operators. 

Further, we repeatedly use notation commonplace in the supervised learning literature. We consider a setting where we observe covariates $\bm{X}:=(\bm{x}_1, \ldots, \bm{x}_n) \in \mathbb{R}^{n \times p}$ and response data $\bm{y}:=(y_1, \ldots y_n) \in \mathbb{R}^n$. %, and assume the data $(\bm{x}_i, y_i)$ are drawn from some underlying stochastic process. 
We say that $(\bm{X}, \bm{y})$ is a training set, and let $\bm{\beta}$ denote a regressor fitted on this training set. In cross-validation, we are also interested in the behavior of $\bm{\beta}$ after {\color{black}leaving out portions of} the training set. We let $(\bm{X}^{(i)}, \bm{y}^{(i)})$ denote the training set with the $i$th data point left out, and denote by $\bm{\beta}^{(i)}$ the regressor obtained after leaving out the $i$th point. {\color{black}Similarly, given a partition $\mathcal{N}_1,\dots,\mathcal{N}_k$  of $[n]$ and $j\in [k]$, we let $(\bm{X}^{(\mathcal{N}_j)}, \bm{y}^{(\mathcal{N}_j)})$ denote the training set with the $j$th fold left out, and $\bm{\beta}^{(\mathcal{N}_j)}$ be the associated regressor.} 



\section{Convex Relaxations of $k$-fold Cross-Validation Error}\label{sec:lowerbounds}
% {\color{black}Section \ref{subsec:lowerbounds} proposes selecting the hyperparameters $(\tau, \gamma)$ in  \eqref{eqn:l0l2_training} by minimizing the function $g$ defined in Problem \eqref{eq:param_opt} as a weighted sum of the $k$-fold cross-validation error $h$ and the output stability of a regressor $\mu_h$. From an optimization perspective, this might appear challenging, because {\color{black}each evaluation of $h_j$ requires solving a MIO, thus,}
% evaluating $g$ involves solving $k+1$ MIOs. %Indeed, in modern leave-one-out cross-validation settings, evaluating $g$ may easily involve solving hundreds or thousands of MIOs. 

In this section, we develop tractable upper and lower approximations of the $k$-fold cross-validation error of a sparse regression model which can be evaluated at a given $(\gamma, \tau)$ without solving any MIOs. From a theoretical perspective, one of our main contributions is that, given $\bm{x}\in \mathbb{R}^p$, we show how to construct bounds $\underline \xi,\overline\xi$ such that $\underline \xi\leq \bm{x^\top\beta}^{(\mathcal{N}_j)}\leq \overline \xi$, which we can use to infer out-of-sample predictions. In particular, we leverage this insight to bound from above and below the function:
\begin{align}
	&h(\gamma, \tau)=1/n \sum_{j=1}^k h_j(\gamma, \tau)=1/n\sum_{j=1}^k\sum_{i \in \mathcal{N}_j}\left(y_i -\bm{x_i}^\top \bm{\beta}^{(\mathcal{N}_j)}(\gamma,\tau)\right)^2,\label{eqn:crossValError}
\end{align} 
which, in turn, bounds the $k$-fold cross-validation error.
\subsection{Bounds on the Prediction Spread}\label{sec:predictionBound}

Given any $0< \gamma$, it is well-known that Problem~\eqref{eqn:l0l2_training} admits the conic quadratic relaxation:
\begin{align}\label{eq:persp}
\zeta_{\text{persp}}=\min_{\bm{\beta} \in \mathbb{R}^p,\bm{z}\in [0,1]^p} \quad  \Vert \bm{y}-\bm{X}\bm{\beta}\Vert_2^2 +\frac{\gamma}{2} \sum_{i=1}^p \frac{\beta_i^2}{z_i} \quad \text{s.t.} \quad &\sum_{i=1}^pz_i \leq \tau,
\end{align}
which is also known as the {perspective relaxation} \citep[][]{ceria1999convex, xie2020scalable}.
If integrality constraints $\bm{z}\in \{0,1\}^p$ are added to \eqref{eq:persp}, then the resulting mixed-integer optimization problem (MIO) is a reformulation of \eqref{eqn:l0l2_training}, where the logical constraints $z_i=0$ if $\beta_i= 0 \ \forall i \in [p]$ are implicitly imposed via the domain of the perspective function $\beta_i^2/z_i$. Moreover, the optimal objective $\zeta_{\text{persp}}$ of \eqref{eq:persp} often provides tight lower bounds on the objective value of \eqref{eqn:l0l2_training} \citep[]{pilanci2015sparse,bertsimas2020sparse, askari2022approximation}, and the optimal solution $\bm{\beta}_{\text{persp}}^*$ is often a good estimator in its own right. As we establish in our main results, the perspective relaxation can also be used to obtain accurate approximations of and bounds on the $k$-fold cross-validation error.

\iffalse Our approach to derive such bounds relies on exploiting a restriction of \eqref{eq:persp}. Given any matrix $\bm{W}\in \mathbb{R}^{q\times p}$ with $q\geq 1$, consider the following perspective reformulation parametrized by $\bm{\xi}\in \mathbb{R}^q$ :
\begin{align}\label{eq:perspRestricted}
	\phi(\bm{\xi})=\min_{\bm{\beta}\in \mathbb{R}^p,\bm{z}\in [0,1]^p}\;& \|\bm{X}\bm{\beta}-\bm{y}\|_2^2+\epsilon\|\bm{\beta}\|_2^2+\frac{\gamma-2\epsilon}{2}\sum_{j=1}^p \frac{\beta_j^2}{z_j}\ \text{s.t.}\ \sum_{j=1}^p z_j\leq \tau,\ \bm{W\beta}=\bm{\xi}.
\end{align}
Moreover, suppose we have access to an upper bound of \eqref{eqn:l0l2_training}, i.e., 
\begin{equation}\label{eq:MIPII}
		\bar u\geq \min_{\bm{\beta}\in \mathbb{R}^p}\; \|\bm{X}\bm{\beta}-\bm{y}\|_2^2+\frac{\gamma}{2}\|\bm{\beta}\|_2^2
		\text{ s.t. }\; \|\bm{\beta}\|_0\leq \tau,
	\end{equation} 
 which can easily be obtained using heuristics, and denote by $\bm{\beta}^\star_{\text{MIO}}$ the optimal solution of \eqref{eq:MIPII}. Then observe that if $\phi(\bm{\xi})>\bar u$, we can immediately conclude that $\bm{W}\bm{\beta}^\star_{\text{MIO}}\neq \bm{\xi}$. We say that values in the set $\left\{\bm{\xi}\in \mathbb{R}^q:\phi(\bm{\xi})\leq \bar u\right\}$ are \emph{admissible} for $\bm{W\bm{\beta}^\star_{\text{MIO}}}$ (and those outside the set are non-admissible), and use the set of admissible values to bound the prediction spread. 
 
 To efficiently construct such bounds, the next proposition states how to construct lower bounds on $\phi(\bm{\xi})$, immediately leading to relaxations of the set of admissible values.

 \begin{proposition}\label{prop:socpdual}
Given any matrix $\bm{W}\in \mathbb{R}^{q\times p}$ and $\bm{\xi}\in \mathbb{R}^q$, we find that 
\begin{align*}
	\phi(\bm{\xi})=\max_{\bm{v}, \bm{w}\in \mathbb{R}^p,\bm{\lambda}\in \mathbb{R}^q}\;& \Vert \bm{y}\Vert_2^2-\bm{v}^\top \left(\bm{X}^\top \bm{X}+\varepsilon\mathbb{I}\right)\bm{v} - \frac{\gamma-2\varepsilon}{8} \sum_{j=1}^\tau w_{[j]}^2 \\
 & -\bm{\lambda}^\top \bm{\xi} + \bm{v}^\top \bm{W}^\top \bm{\lambda}-\frac{1}{4}\bm{\lambda}^\top \bm{W}\left(\bm{X}^\top \bm{X}+\varepsilon\mathbb{I}\right)^{-1}\bm{W}^\top\bm{\lambda}\\
 \text{s.t} \quad & \left(\bm{X}^\top \bm{X}+\varepsilon\mathbb{I}\right) \bm{v}=\bm{X}^\top \bm{y}-\frac{\gamma-2\varepsilon}{4}\bm{w}.
\end{align*}
 \end{proposition}
 Observe that Proposition~\ref{prop:socpdual} expresses $\phi$ as a SOCP-representable unconstrained optimization problem -- variables $\bm{v}$ and the equality constraints can be trivially projected out and removed. Moreover, a lower bound on $\phi$ can be easily obtained by fixing the variables in Proposition~\ref{prop:socpdual} to any value.

\fi

\ag{Our first main result (Theorem~\ref{prop:kfoldspread}) reveals that any optimal solution of \eqref{eqn:l0l2_training} lies in an ellipsoid centered at its continuous (perspective) relaxation, and whose radius depends on the duality gap: %This is significant because while there are a priori bounds on the duality gap in the literature in terms of the rank of $\bm{X}$ \citep{askari2022approximation}, these bounds do not give information about the distance between solutions to the MIO and its perspective relaxation, which is uninformative for hyperparameter selection. On the other hand, the following result gives a posteriori guarantees on the distance between optimal solutions to a MIO and its perspective relaxation, which gives rise to an efficient hyperparameter selection strategy:
}


\begin{theorem}\label{prop:kfoldspread}
	Given any $0 < \gamma$ and any bound
	\begin{equation}\label{eq:MIPII}
		\bar u\geq \min_{\bm{\beta}\in \mathbb{R}^p}\; \|\bm{X}\bm{\beta}-\bm{y}\|_2^2+\frac{\gamma}{2}\|\bm{\beta}\|_2^2
		\text{\rm\ s.t. }\; \|\bm{\beta}\|_0\leq \tau,
	\end{equation} 
	the inequality
	\begin{align}\label{ineq:perspbound}
		(\bm{\beta^\star_{\text{persp}}}- \bm{\beta}^\star_{\text{MIO}})^\top \left(\bm{X}^\top \bm{X}+\frac{\gamma}{2} \bm{\mathbb{I}}\right) (\bm{\beta^\star_{\text{persp}}}- \bm{\beta}^\star_{\text{MIO}}) \leq (\bar{u}-\zeta_{\text{persp}})
	\end{align}
	holds, where $\bm{\beta}_{MIO}^*$ is an optimal solution of \eqref{eq:MIPII} and $\bm{\beta}_{persp}^*$ is optimal to \eqref{eq:persp}.
\end{theorem}



\proof{Proof of Theorem \ref{prop:kfoldspread}}{\color{black}
Let $\epsilon>0$ be a small positive constant and let
\begin{align}\label{persp:peturbed}
f_\epsilon(\bm{\beta})&:= \min_{\bm{z} \in [0, 1]^p: \bm{e}^\top \bm{z}\leq k} \quad \|\bm{X}\bm{\beta}-\bm{y}\|_2^2+\frac{\gamma}{2}\sum_{i \in [p]}\frac{\beta_i^2+\epsilon}{z_i}-\frac{\epsilon\gamma p}{2},
\end{align}
denote the objective value of the perspective relaxation at a given $\bm{\beta}$, where we apply the small perturbation $\epsilon$ so that $z_i^\star >0$. Note that $f_\epsilon$ is non-decreasing in $\epsilon$. The function $f(\bm{\beta})$ is twice differentiable with respect to $\bm{\beta}$, and admits the following integral Taylor series expansion about $\bm{\beta}_{persp}^*$, an optimal solution to \eqref{persp:peturbed} \citep[e.g.,][Lemma 3.5.3]{sidford2024notes}
\begin{align*}
f_\epsilon(\bm{\beta})=& f_\epsilon(\bm{\beta}_{persp}^*)+\langle f_\epsilon(\bm{\beta}_{persp}^*), \bm{\beta}-\bm{\beta}_{persp}^*\rangle\\
& +\int_{0}^1 (1-\alpha)(\bm{\beta}-\bm{\beta}_{persp}^*)^\top \nabla^2 f_\epsilon(\bm{\beta}_{persp}^*+\alpha (\bm{\beta}-\bm{\beta}_{persp}^*))(\bm{\beta}-\bm{\beta}_{persp}^*)d\alpha
\end{align*}

Moreover, the Hessian at a given $\bm{\beta}$ is $\nabla^2 f_\epsilon(\bm{\beta})=2 \bm{X}^\top \bm{X}+\gamma\mathrm{Diag}(\bm{z}^*)^{-1}$, where $\bm{z}^*>\bm{0}$ because of the perturbation term in the objective. Since $\bm{z}^*\leq \bm{e}$, the Hessian is such that $\nabla^2 f_\epsilon(\bm{\beta}) \succeq 2 \bm{X}^\top \bm{X}+\gamma\mathbb{I}$. Moreover, replacing $\nabla^2 f_\epsilon(\bm{\beta})$ with a valid lower bound with respect to the Loewener partial order gives a lower bound on $f(\bm{\beta})$. Thus, integrating with respect to $\alpha$ yields the bound
\begin{align*}
f_\epsilon(\bm{\beta})\geq & f_\epsilon(\bm{\beta}_{persp}^*)+(\bm{\beta}-\bm{\beta}_{persp}^*)^\top (\bm{X}^\top \bm{X}+\frac{\gamma}{2} \mathbb{I})(\bm{\beta}-\bm{\beta}_{persp}^*),
\end{align*}
where we omit the first-order term $\langle \nabla f(\bm{\beta}_{persp}^*), \bm{\beta}-\bm{\beta}_{persp}\rangle$ because it is non-negative for an optimal $\bm{\beta}_{persp}^*$ \citep[cf.][Chap. 1]{bertsekas2016nonlinear}. 

The result then follows by inserting $\bm{\beta}_{MIO}$ into the bound, taking limits as $\epsilon \rightarrow 0$ to avoid including perturbation terms, and noting that $f(\bm{\beta}_{MIO})$ does not require that $\bm{z}$ is integral, and thus is a lower bound on $\bar{u}$. We remark that taking limits is justified by, e.g., the monotone convergence theorem \citep{grimmett2020probability}. Indeed, the objective value of $f_\epsilon(\bm{\beta}_{persp}^*)$ is non-increasing as we decrease $\epsilon$, bounded from below by $\zeta_{persp}$, and attains this bound in the limit. 
%Consider the perspective reformulation of \eqref{eq:MIPII} parametrized by $\bm{\xi}\in \mathbb{R}^p$ :
% \begin{align*}
% 	\phi(\bm{\xi})=\min_{\bm{\beta}\in \mathbb{R}^p,\bm{z}\in [0,1]^p}\;& \|\bm{X}\bm{\beta}-\bm{y}\|_2^2+\epsilon\|\bm{\beta}\|_2^2+\frac{\gamma-2\epsilon}{2}\sum_{j=1}^p \frac{\beta_j^2}{z_j}\ \text{s.t.}\ \sum_{j=1}^p z_j\leq \tau,\ \bm{\beta}=\bm{\xi}.
% \end{align*}
% \ag{As noted in \cite{atamturk2020safe}, we have $\beta_j^2/z_j=\max_{w\in \mathbb{R}}w\beta_j-(w^2/4)z_j$ and therefore
% \begin{align*}
% 	\phi(\bm{\xi})=\max_{\bm{w}\in \mathbb{R}^p}\min_{\bm{\beta}\in \mathbb{R}^p,\bm{z}\in [0,1]^p}\;& \|\bm{X}\bm{\beta}-\bm{y}\|_2^2+\epsilon\|\bm{\beta}\|_2^2+\frac{\gamma-2\epsilon}{2}\sum_{j=1}^p \left(w_j\beta_j-\frac{w_j^2}{4}z_j\right)\ \text{s.t.}\; \sum_{j=1}^p z_j\leq \tau,\; \bm{\beta}=\bm{\xi}.
% \end{align*}
% Moreover, letting $\bm{\lambda}\in \mathbb{R}^p$ denote the dual variables associated with the equality constraints $\bm{\beta}=\bm{\xi}$ and taking a Lagrangian relaxation, we find the equivalent expression
% }
% \begin{equation}\label{eq:lagrangian}
% \begin{aligned}
% 	\phi(\bm{\xi})=\max_{\bm{w}\in \mathbb{R}^p,\lambda\in \mathbb{R}^q}\min_{\bm{\beta}\in \mathbb{R}^p,\bm{z}\in [0,1]^p}\;& \|\bm{X}\bm{\beta}-\bm{y}\|_2^2+\epsilon\|\bm{\beta}\|_2^2+\frac{\gamma-2\epsilon}{2}\sum_{j=1}^p \left(w_j\beta_j-\frac{w_j^2}{4}z_j\right)-\bm{\lambda}^\top\xi+\bm{\lambda}^\top \bm{\beta}\\
% 	\text{s.t.}\;& \sum_{j=1}^p z_j\leq \tau.
% \end{aligned}
% \end{equation}
% \ag{Taking derivatives with respect to variables $\bm{\beta}$ and setting to $\bm{0}$, we find that }in any optimal solution of the inner minimization problem\begin{align*}
% \bm{0}=2(\bm{X}^\top \bm{X}+\epsilon\mathbb{I})\bm{\beta}-2\bm{X}^\top \bm{y}+\frac{\gamma-2\epsilon}{2}\bm{w}+ \bm{\lambda}\ 
% \Leftrightarrow\bm{\beta^*}=\left(\bm{X}^\top \bm{X}+\epsilon\mathbb{I}\right)^{-1}\left(\bm{X}^\top \bm{y}-\frac{\gamma-2\epsilon}{4}\bm{w}-\frac{1}{2} \bm{\lambda}\right).
% \end{align*}
% Moreover, we also find that\ag{, in optimal solutions,} $z_j^*=1$ for the largest $\tau$ values of $w_j^2$ and $z_j^*=0$ otherwise -- we use $\sum_{j=1}^\tau w_{[j]}^2$ to denote $\sum_{j=1}^p w_j^2z_j^*$. Finally, the optimal value of $\bm{\lambda}$ is such that 
% \begin{align}
% &\bm{\xi}=\bm{\beta^*}\ \Leftrightarrow\; \bm{\xi}=\left(\bm{X}^\top \bm{X}+\epsilon\mathbb{I}\right)^{-1}\left(\bm{X}^\top \bm{y}-\frac{\gamma-2\epsilon}{4}\bm{w}-\frac{1}{2} \bm{\lambda}\right)\notag\\
% \Leftrightarrow\;& \frac{1}{2} \left(\bm{X}^\top \bm{X}+\epsilon\mathbb{I}\right)^{-1} \bm{\lambda} = \left(\bm{X}^\top \bm{X}+\epsilon\mathbb{I}\right)^{-1}\left(\bm{X}^\top \bm{y}-\frac{\gamma-2\epsilon}{4}\bm{w}\right)-\bm{\xi}\label{eq:valLambda}\\
% \Leftrightarrow\;& \bm{\lambda} = 2\left(\bm{X}^\top \bm{y}-\frac{\gamma-2\epsilon}{4}\bm{w}\right)-2\left(\bm{X}^\top \bm{X}+\epsilon\mathbb{I}\right)\bm{\xi}.\label{eq:valLambda2}
% \end{align}
% First, we set $\bm{\beta}$ to \ag{$\bm{\beta^*}$ in \eqref{eq:lagrangian}} and subtract $\left(\bm{\beta^*}\right)^\top \left(2(\bm{X}^\top \bm{X}+\epsilon\mathbb{I})\bm{\beta^*}-2\bm{X}^\top \bm{y}+\frac{\gamma-2\epsilon}{2}\bm{w}+ \bm{\lambda}\right)$ from the \ag{objective}, since by the KKT conditions this expression equals zero at optimality. This gives:
% \begin{align*}
% 	\phi(\bm{\xi})=\max_{\bm{w}\in \mathbb{R}^p,\lambda\in \mathbb{R}^q}\;& \Vert \bm{y}\Vert_2^2-\bm{\beta}^{*\top} \left(\bm{X}^\top \bm{X}+\varepsilon\mathbb{I}\right)\bm{\beta^*} - \frac{\gamma-2\varepsilon}{8} \sum_{j=1}^\tau w_{[j]}^2 -\bm{\lambda}^\top \bm{\xi}.
% \end{align*}
% Next, we introduce a dummy variable $\bm{v}$ for simplicity, where $\bm{v}$ is such that
% \begin{align*}
%     \left(\bm{X}^\top \bm{X}+\varepsilon\mathbb{I}\right)\bm{v}=\bm{X}^\top \bm{y}-\frac{\gamma-2\varepsilon}{4}\bm{w}.
% \end{align*}
% This allows us to rewrite the above expressions for $\bm{\beta^*}$ in terms of $\bm{\lambda}$, and the optimal value of $\bm{\lambda}$ as 
% \begin{align*}
%     &\bm{\beta}^\star=\bm{v}-\frac{1}{2}\left(\bm{X}^\top \bm{X}+\varepsilon\mathbb{I}\right)^{-1}  \bm{\lambda} \text{, and}\ \bm{\lambda}=2\left(\bm{X}^\top \bm{X}+\varepsilon \mathbb{I}\right)\left(\bm{v}-\bm{\xi}\right).
% \end{align*}
% Rearranging our expression for $\phi$ by substituting for $\bm{\beta^*}$ then gives the expression
% \begin{align*}
% 	\phi(\bm{\xi})=\max_{\bm{v}, \bm{w}\in \mathbb{R}^p,\bm{\lambda}\in \mathbb{R}^q}\;& \Vert \bm{y}\Vert_2^2-\bm{v}^\top \left(\bm{X}^\top \bm{X}+\varepsilon\mathbb{I}\right)\bm{v} - \frac{\gamma-2\varepsilon}{8} \sum_{j=1}^\tau w_{[j]}^2 -\bm{\lambda}^\top \bm{\xi} + \bm{v}^\top  \bm{\lambda}-\frac{1}{4}\bm{\lambda}^\top \left(\bm{X}^\top \bm{X}+\varepsilon\mathbb{I}\right)^{-1}\bm{\lambda}\\
%  \text{s.t} \quad & \left(\bm{X}^\top \bm{X}+\varepsilon\mathbb{I}\right) \bm{v}=\bm{X}^\top \bm{y}-\frac{\gamma-2\varepsilon}{4}\bm{w}.
% \end{align*}
% Moreover, substituting $\bm{\lambda}^\star$
% into $\phi$ and adding the term $\bm{0}=\bm{\lambda}^\top \left(\frac{1}{2}\left(\bm{X}^\top \bm{X}+\varepsilon\mathbb{I}\right)^{-1} \bm{\lambda}-\bm{v}+\bm{\xi}\right)$ allows us to rewrite this expression as
% a maximization problem in $\bm{v,w}$ only, namely:
% \begin{align*}
% 	\phi(\bm{\xi})=\max_{\bm{v}, \bm{w}\in \mathbb{R}^p}\;& \Vert \bm{y}\Vert_2^2-\bm{v}^\top \left(\bm{X}^\top \bm{X}+\varepsilon\mathbb{I}\right)\bm{v} - \frac{\gamma-2\varepsilon}{8} \sum_{j=1}^\tau w_{[j]}^2 +\left(\bm{v}- \bm{\xi}\right)^\top \left(\bm{X}^\top \bm{X}+\varepsilon\mathbb{I}\right)\left(\bm{v}- \bm{\xi}\right)\\
%  \text{s.t} \quad & \left(\bm{X}^\top \bm{X}+\varepsilon\mathbb{I}\right) \bm{v}=\bm{X}^\top \bm{y}-\frac{\gamma-2\varepsilon}{4}\bm{w}.
% \end{align*}
% Therefore, to construct an interval of admissible values, it suffices to fix $\bm{w},\bm{v}$ to any feasible solution of \eqref{eq:perspbound}. Therefore, letting $\zeta(\bm{w},\bm{v})=\|\bm{y}\|_2^2-\bm{v}^\top \left(\bm{X}^\top \bm{X}+\epsilon\mathbb{I}\right)\bm{v}-\frac{\gamma-2\epsilon}{8}\sum_{j=1}^\tau w_{[j]}^2$, it follows that non admissible values of $\xi$ do not satisfy
% \begin{align*}&\zeta(\bm{w},\bm{v})+\left(\bm{v}- \bm{\xi}\right)^\top \left(\bm{X}^\top \bm{X}+\varepsilon\mathbb{I}\right)\left(\bm{v}- \bm{\xi}\right)\leq \bar u\\
% \Leftrightarrow\;&\left(\bm{v}- \bm{\xi}\right)^\top \left(\bm{X}^\top \bm{X}+\varepsilon\mathbb{I}\right)\left(\bm{v}- \bm{\xi}\right)\leq \left(\bar u-\zeta(\bm{w},\bm{v})\right)\ \Leftrightarrow\; \Vert  \left(\bm{X}^\top \bm{X}+\varepsilon\mathbb{I}\right)^{1/2} \left(\bm{v}- \bm{\xi}\right)\Vert_2^2 \leq (\bar{u}-\zeta(\bm{w},\bm{v}))
% \begin{pmatrix} 
% \bar u-\zeta(\bm{w},\bm{v}) & \left(\bm{v}-\bm{W}^\dag \bm{\xi}\right)^\top \\
% \left(\bm{v}-\bm{W}^\dag \bm{\xi}\right) & \left(\bm{X}^\top \bm{X}+\varepsilon\mathbb{I}\right)^{-1}
% \end{pmatrix} \succeq \bm{0},
% \end{align*}
% \ag{To conclude the proof, we fix $\bm{v},\bm{w}$ to the optimal dual variables of \eqref{eq:persp}. These values correspond to setting $\bm{v^\star}= \beta_{\text{persp}}^\star$ (and setting $\bm{w^\star}$ accordingly), in which case $\zeta(\bm{w^\star},\bm{v^\star})=\zeta_{\text{persp}}$, see \cite{atamturk2020safe} for the proof. Moreover, since the result holds for any $\bm{\xi}$, we let $\bm{\xi}=\bm{\beta}^\star_{\text{MIO}}$, concluding the proof}. 
}
\Halmos
\endproof

% \begin{remark}
% The proof of Theorem \ref{prop:kfoldspread} does not directly exploit the fact that $\bm{\beta}_{persp}$ is an optimal solution to a perspective relaxation, rather than a different relaxation. Thus, one could replace the perspective relaxation with a tighter relaxation here, as explored in \cite{atamturk2019rank}. We do not explore this here, as the perspective relaxation is particularly tractable, and we are interested in applying the results developed here to large-scale settings.
% \end{remark}


\ag{Using Theorem~\ref{prop:kfoldspread}, we can compute bounds on $h_j(\gamma,\tau)$ in \eqref{eqn:crossValError} by solving problems of the form 
\begin{subequations}\label{eq:boundComputation}
\begin{align}
\min/\max\;&\sum_{i\in \mathcal{N}_j} \left(y_i -\bm{x_i}^\top \bm{\beta}\right)^2\\
\text{s.t.}\;&(\bm{\beta_{\text{persp}}}^{(\mathcal{N}_j)}- \bm{\beta})^\top \left((\bm{X}^{(\mathcal{N}_j)})^\top \bm{X}^{(\mathcal{N}_j)}+\frac{\gamma}{2} \bm{\mathbb{I}}\right) (\bm{\beta_{\text{persp}}}^{(\mathcal{N}_j)}- \bm{\beta}) \leq (\bar u^{(\mathcal{N}_j)}-\zeta_{\text{persp}}^{(\mathcal{N}_j)}),
\end{align}
\end{subequations}
where $\bm{\beta_{\text{persp}}}^{(\mathcal{N}_j)}$ and $\zeta_{\text{persp}}^{(\mathcal{N}_j)}$ are the optimal solution and objective value of the perspective relaxation with fold $\mathcal{N}_j$ removed, and $\bar u^{(\mathcal{N}_j)}$ is an associated upper bound. Bounds for the function $h(\gamma,\tau)$ then immediately follow by simply adding the bounds associated with $h_j(\gamma,\tau)$ for all $j\in [k]$. %Bounds for $\upsilon_{i,j}(\gamma,\tau)$ in \eqref{eqn:epsilondefn} could be obtained similarly by simply updating the objective. However, we provide a simpler method in the next section.
}

\begin{remark}[Computability of \ag{the} bounds]\label{rem:computability}
	Observe that a lower bound on the $k$-fold error can easily be computed by solving a convex quadratically constrained quadratic problem, while an upper bound can be computed by noticing that the maximization problem \eqref{eq:boundComputation} is a trust region problem in $\bm{\beta}$, which can be reformulated as a semidefinite problem \citep{hazan2016linear}. One could further tighten these bounds by imposing a sparsity constraint on $\bm{\beta}$, but this may not be practically tractable.
\end{remark}

% \begin{remark}[Choice of parameter $\epsilon$]
% 	The role of the parameter $\epsilon$ is to ensure that $\bm{x}$ is in the range of the matrix $\bm{X^\top X}+\epsilon \mathbb{I}$. If the matrix $\bm{x}$ is already in the range of $\bm{X}$, as occurs, for instance, after solving a convex relaxation of the training problem without any observations omitted, then one should set $\epsilon=0$ and replace the inverse in Theorem \ref{prop:kfoldspread}'s bound with a pseudoinverse, as lower values of $\epsilon$ result in stronger perspective relaxations. However, this is not possible in general; for instance, $\bm{x}_i$ need not be in the range of $\bm{X}^{(\mathcal{N}_j)}$ in highly underdetermined settings.
% \end{remark}



%Problems \eqref{eq:boundComputation} can in theory be computed efficiently. In practice, however, solving the required semidefinite problems to optimality may be difficult if $p$ is in the thousands. Moreover, while bounding $k$-fold error can indeed be solved via convex optimization, computing the required stability bounds may require solving non-convex problems. In the next section we derive alternative bounds that, while weaker, can be obtained in closed form.

\subsection{Closed-form Bounds on the Prediction Spread}
\ag{While solving the perspective relaxation \eqref{eq:persp} is necessary to solve the MIO \eqref{eq:MIPII} via branch-and-bound (in particular, the perspective relaxation is the root node in a branch-and-bound scheme \citep{mazumder2023subset}), the additional two optimization problems \eqref{eq:boundComputation} are not. Moreover, solving trust-region problems can be expensive in large-scale problems. Accordingly, in this section, we present alternative bounds that may be weaker, but can be obtained in closed form. In numerical experiments (Section \ref{sec:numres}), these closed-form bounds already reduce the number of MIOs that need to be solved by up to 
80\% when compared to grid search.}

\begin{theorem}\label{prop:l10UB}
	Given any vector $\bm{x}\in \mathbb{R}^p$ and any bound
	\begin{equation}\label{eq:MIPI}
			\bar u\geq \min_{\bm{\beta}\in \mathbb{R}^p}\; \|\bm{X}\bm{\beta}-\bm{y}\|_2^2+\frac{\gamma}{2}\|\bm{\beta}\|_2^2
			\text{\rm\ s.t. }\; \|\bm{\beta}\|_0\leq \tau,
	\end{equation} 
	the inequalities
	\small
	$$\bm{x^\top\beta}_{persp}^*-\sqrt{\left(\bar u-\zeta_{\text{persp}}\right)\bm{x^\top}\left(\bm{X^\top X}+\frac{\gamma}{2}\mathbb{I}\right)^{-1}\bm{x}}\leq \bm{x}^\top \bm{\beta}_{MIO}^*\leq  \bm{x^\top\beta}_{persp}^*+\sqrt{\left(\bar u-\zeta_{\text{persp}}\right)\bm{x^\top}\left(\bm{X^\top X}+\frac{\gamma}{2}\mathbb{I}\right)^{-1}\bm{x}}$$\normalsize
	hold, where $\bm{\beta}_{MIO}^*$ is an optimal solution of \eqref{eq:MIPI} and $\bm{\beta}_{persp}^*$ is optimal to \eqref{eq:persp}.
\end{theorem}

\proof{Proof of Theorem~\ref{prop:l10UB}}
From Theorem~\ref{prop:kfoldspread}, we have the inequality
\begin{align}
(\bm{\beta^\star_{\text{persp}}}- \bm{\beta}^\star_{\text{MIO}})^\top \left(\bm{X}^\top \bm{X}+\frac{\gamma}{2} \bm{\mathbb{I}}\right) (\bm{\beta^\star_{\text{persp}}}- \bm{\beta}^\star_{\text{MIO}}) \leq (\bar{u}-\zeta_{\text{persp}}).
	\end{align}
 By the Schur Complement Lemma \citep[see, e.g.,][]{boyd1994linear}, this is equivalent to
 \begin{align*}
     (\bar{u}-\zeta_{\text{persp}}) \left(\bm{X}^\top \bm{X}+\frac{\gamma}{2} \bm{\mathbb{I}}\right)^{-1}\succeq (\bm{\beta^\star_{\text{persp}}}- \bm{\beta}^\star_{\text{MIO}})(\bm{\beta^\star_{\text{persp}}}- \bm{\beta}^\star_{\text{MIO}})^\top
 \end{align*}
 Next, we can left/right multiply this expression by an arbitrary matrix $\bm{W} \in \mathbb{R}^{m \times p}$. This gives:
 \begin{align*}
     (\bar{u}-\zeta_{\text{persp}}) \bm{W}\left(\bm{X}^\top \bm{X}+\frac{\gamma}{2} \bm{\mathbb{I}}\right)^{-1}\bm{W}^\top\succeq (\bm{W}\bm{\beta^\star_{\text{persp}}}- \bm{W}\bm{\beta}^\star_{\text{MIO}})(\bm{W}\bm{\beta^\star_{\text{persp}}}- \bm{W}\bm{\beta}^\star_{\text{MIO}})^\top.
 \end{align*}
In particular, setting $\bm{W}=\bm{x}^\top$ for a vector $\bm{x} \in \mathbb{R}^{p}$ gives the inequality
 \begin{align*}
     (\bar{u}-\zeta_{\text{persp}}) \bm{x}^\top \left(\bm{X}^\top \bm{X}+\frac{\gamma}{2} \bm{\mathbb{I}}\right)^{-1} \bm{x}\geq (\bm{x}^\top(\bm{\beta^\star_{\text{persp}}}- \bm{\beta}^\star_{\text{MIO}}))^2,
 \end{align*}
 which we rearrange to obtain the result.
	%\ag{The main idea of the proof is similar to Theorem~\ref{prop:kfoldspread}, but instead of imposing the multivariate constraint $\bm{\beta}=\bm{\xi}$ we impose the simpler constraint $\bm{x^\top\beta}=\xi$. In other words, instead of controlling for the value of the regression coefficients, we control for the prediction of an observation with features $\bm{x}$. The arguments until Equation \eqref{eq:valLambda2} are essentially identical. At that point, due to having a one-dimensional vector $\xi$, it is possible to derive the expression of the ball with more direct arguments. We defer the details to Appendix~\ref{ssec:propub}. }
\Halmos
	
\endproof

\begin{corollary}
    For any $\bm{W} \in \mathbb{R}^{m \times p}$ we have that
 \begin{align*}
     (\bar{u}-\zeta_{\text{persp}}) \mathrm{tr}\left(\bm{W} \left(\bm{X}^\top \bm{X}+\frac{\gamma}{2} \bm{\mathbb{I}}\right)^{-1} \bm{W}^\top\right)\geq \Vert 
     \bm{W}(\bm{\beta^\star_{\text{persp}}}- \bm{\beta}^\star_{\text{MIO}})\Vert_2^2,
 \end{align*}
\end{corollary}


Applying Theorem~\ref{prop:l10UB} to the problem 
$$	\bar u^{(\mathcal{N}_j)}\geq \min_{\bm{\beta}\in \mathbb{R}^p}\; \|\bm{X}^{(\mathcal{N}_j)}\bm{\beta}-\bm{y}^{(\mathcal{N}_j)}\|_2^2+\frac{\gamma}{2}\|\bm{\beta}\|_2^2
\text{ s.t. }\; \|\bm{\beta}\|_0\leq \tau,$$  we have the bounds
\begin{align*}
	\underline{\xi}_{i,j}:=\bm{x_i^\top}\bm{\beta}_{persp}^*-\sqrt{\bm{x_i^\top}\left(\bm{X}^{(\mathcal{N}_j)^\top} \bm{X}^{(i)}+\frac{\gamma}{2}\mathbb{I}\right)^{-1}\bm{x_i}\left(\bar u^{(\mathcal{N}_j)}-\zeta^{(\mathcal{N}_j)}\right)},\\
	\bar{\xi}_{i,j}:=\bm{x_i^\top}\bm{\beta}_{persp}^*+\sqrt{\bm{x_i^\top}\left(\bm{X}^{(\mathcal{N}_j)^\top} \bm{X}^{(\mathcal{N}_j)}+\frac{\gamma}{2}\mathbb{I}\right)^{-1}\bm{x_i}\left(\bar u^{(\mathcal{N}_j)}-\zeta^{(\mathcal{N}_j)}\right)}
\end{align*}
where $0 < 2\epsilon < \gamma$ and $\underline{\xi}\leq \bm{x_i^\top}\bm{\beta}_{MIO}^*\leq \bar\xi$. %We can use these bounds to bound \ag{terms $\upsilon_{i,j}(\gamma,\tau)$ in \eqref{eqn:epsilondefn}}.

\begin{corollary}\label{corr:perspUB}
	We have the following bounds on the $i$th prediction error \ag{associated with fold $j$}
	\begin{align}\label{eqn:perspbounds}
		\max\left((y_i-\underline{\xi}_{i,j})^2, (y_i-\bar{\xi}_{i,j})^2\right) \geq \ag{\nu_{i,j}(\gamma, \tau)} \geq \begin{cases} (y_i-\underline{\xi}_{i,j})^2 \quad & \text{if} \ y_i < \underline{\xi}_{i,j}\\
			0 \quad & \text{if} \ y_i \in [\underline{\xi}_{i,j}, \bar{\xi}_{i,j}],\\
			(\bar{\xi}_{i,j}-y_i)^2 \quad & \text{if} \ y_i > \bar{\xi}_{i,j}.
		\end{cases}
	\end{align}
\end{corollary}
Moreover, since $h(\gamma,\tau)=\frac{1}{n}\sum_{j=1}^k\sum_{i\in \mathcal{N}_j}\nu_{i,j}(\gamma,\tau)$, we can compute lower and upper bounds on the $k$-th fold cross-validation error \ag{by adding the individual bounds. Observe that the bounds computed by summing disaggregated bounds could be substantially worse than those obtained by letting $\bm{W}$ be a matrix with all omitted columns in the $j$th fold of $\bm{X}$ in the proof of Theorem \ref{prop:l10UB}. Nonetheless, the approach outlined here might be the only one feasible in large scale instances, as they are obtained directly from the perspective relaxation without solving additional optimization problems, while an aggregated approach would involve solving an auxiliary semidefinite optimization problem. Despite the loss in quality, we show in our computational sections that (combined with the methods discussed in \S\ref{sec:corddescent}), the disaggregated bounds are sufficient to lead to a 50\%-80\% reduction in the number of MIO solved with respect to grid search.}

We conclude this subsection with two remarks.

\begin{remark}[Relaxation Tightness]
	If the perspective relaxation is tight, as occurs 
 when $n$ is sufficiently large under certain assumptions on the data generation process \citep{pilanci2015sparse, reeves2019all}
 then 	$\underline{\xi}=\bar{\xi}=\bm{x_i^\top}\bm{\beta}_{persp}^*$, and Corollary \ref{corr:perspUB}'s bounds on the cross-validation error are definitionally tight. Otherwise, as pointed out in Remark~\ref{rem:intuition}, \eqref{eqn:perspbounds}'s bound quality depends on the tightness of the relaxation and on how close the features $\bm{x}_i$ are to the rest of the data. 
\end{remark}


\begin{remark}[Intuition]\label{rem:intuition}
	Theorem~\ref{prop:l10UB} states that $\bm{x}^\top \bm{\beta}_{MIO}^*\approx\bm{x^\top\beta}_{persp}^*$, where the approximation error is determined by two components. The quantity $\sqrt{\bar u-\zeta_{\text{persp}}}$ is related to the strength of the perspective relaxation, with a stronger relaxation resulting in a better approximation. The quantity $\sqrt{\bm{x^\top}\left(\bm{X^\top X}+\frac{\gamma}{2}\mathbb{I}\right)^{-1}\bm{x}}$ is related to the likelihood that $\bm{x}$ is generated from the same distribution as the rows of $\bm{X}$, with larger likelihoods resulting in better approximations. Indeed, if $n>p$, each column of $\bm{X}$ has $0$ mean but has not been standardized, and each row of $\bm{X}$ is generated iid from a multivariate Gaussian distribution, then $\frac{n(n-1)}{n+1}\bm{x^\top}\left(\bm{X^\top X}\right)^{-1}\bm{x}\sim T^2(p, n-1)$ is Hotelling's two-sample T-square test statistic \citep{hotelling1931generalization}, used to test whether $\bm{x}$ is generated from the same Gaussian distribution. Note that if $\bm{x}$ is drawn from the same distribution as the rows of $\bm{X}$ (as may be the case in cross-validation), then $\mathbb{E}\left[\bm{x^\top}\left(\bm{X^\top X}\right)^{-1}\bm{x}\right]=\frac{p(n+1)}{n(n-p-2)}$.
\end{remark}

\iffalse
\subsection{Approximating the Hypothesis Stability} \label{ssec:genbounds2}

We now leverage Theorem \ref{prop:l10UB} to develop bounds on the hypothesis stability \eqref{eq:hypothesisStability}\ag{, which we repeat for convenience}
\begin{align*}
	\mu_h:=\max_{j\in [k]}\frac{1}{n}\sum_{i=1}^n \left\vert (y_i-\bm{x_i^\top\bm{\beta}^\star} )^2-(y_i-\bm{x_i^\top \beta}^{(\mathcal{N}_j)} )^2\right\vert.,
\end{align*}
\ag{where $\bm{\beta}^\star$ is the regressor obtained by solving the MIO over the entire dataset and $\bm{\beta}^{(\mathcal{N}_j)}$ is the regressor obtained by leaving out fold $j$. From Theorem~\ref{prop:l10UB}, we obtain the following inequalities.}


%To develop these bounds, let us denote by $u$ the optimal value of Problem \eqref{eq:MIPOut0}, i.e., the optimal value of a sparse regression problem run on the entire dataset, let $\zeta$ denote the optimal value of Problem \eqref{eq:MIPOut0}'s perspective relaxation, i.e.,
%\begin{align*}
%	\zeta=\min_{\bm{\beta}\in \mathbb{R}^p,\bm{z}\in [0,1]^p}\;& \|\bm{X}\bm{\beta}-\bm{y}\|_2^2+\frac{\gamma}{2}\sum_{j=1}^p \frac{\beta_j^2}{z_j}\ \text{s.t.}\ \sum_{j=1}^p z_j\leq \tau,
%\end{align*} 
%let $\bm{\beta}_{MIO}$ denote an optimal solution to Problem \eqref{eq:MIPOut0} and $\bm{\beta}_{persp}$ denote an optimal solution to its perspective relaxation. Further, let $u^{(i)}, \zeta^{(i)}$, $\bm{\beta}^{(i)}_{MIO}$, $\bm{\beta}^{(i)}_{persp}$ be equivalent objects with the $i$th data point left out of the training problem. Then, we have the following corollary to Theorem \ref{prop:l10UB}:

\begin{corollary}\label{corr:confidenceadjustment}
    Let $0 < 2\epsilon < \gamma$. Then, we have the following bounds:
    \begin{align}
       \vert \bm{x}_i^\top \bm{\beta}^\star - \bm{x}_i^\top \bm{\beta}_{persp}^\star\vert \leq & \sqrt{\bm{x}_i^\top \left(\bm{X}^\top \bm{X}+\epsilon\mathbb{I}\right)^{-1} \bm{x}_i(u-\zeta)},\\
        \vert \bm{x}_i^\top \bm{\beta}^{(\mathcal{N}_j)} -\bm{x}_j^\top \bm{\beta}_{persp}^{(\mathcal{N}_j)}\vert \leq & \sqrt{\bm{x}_i^\top \left(\bm{X}^{(\mathcal{N}_j)\top} \bm{X}^{(\mathcal{N}_j)}+\epsilon\mathbb{I}\right)^{-1} \bm{x}_i(u^{(\mathcal{N}_j)}-\zeta^{(\mathcal{N}_j)})},
    \end{align}
    \ag{where $\bm{\beta}_{persp}^{\star}$ and $\bm{\beta}_{persp}^{(\mathcal{N}_j)}$ and solutions of the perspective relaxations with all data and with fold $j$ removed, respectively, and $u$, $u^{(\mathcal{N}_j)},\zeta, \zeta^{(\mathcal{N}_j)}$ are upper bounds on the MIO objective values and optimal values for the respective perspective relaxations.}
\end{corollary}
Given these bounds, we can produce lower and upper bounds on the quantities $(y_i-\bm{x}_i^\top{\bm{\beta}^\star}^\top )^2$ and $(y_i-\bm{x_i}^\top\bm{\beta}^{(\mathcal{N}_j)})^2$ in essentially the same manner as Corollary \ref{corr:perspUB}. 

\ag{[I would probably remove a lot about the discussion below.]}
% \begin{remark}
% Similarly to Corollary \ref{corr:lowerbound2}, one can also show that 
% \begin{align}
%     u-\zeta^{(j)} \geq (y_j-\bm{x}_j^\top \bm{\beta}_{\text{MIO}})^2,%,\\
%     %u^{(i)}-\zeta^{(i,j)}\geq (y_j-\bm{x}_j^\top \bm{\beta}^{(i)}_{MIO})^2,
% \end{align}
% where we can replace $u$ with an upper bound on $u$ without loss of correctness.
% % where $\zeta^{(i,j)}$ denotes the optimal value of a perspective relaxation after leaving out the $i$th and $j$th data points.
% Therefore, in practice, we use this result to improve our upper bound on $(y_j-\bm{x}_j^\top \bm{\beta}_{MIO})^2$ as part of our overall approach.
% \end{remark}

From Corollary \ref{corr:confidenceadjustment}, we observe that the hypothesis stability $\mu_h$ is very similar for $\bm{\beta}^\star_{\text{MIO}}$ and $\bm{\beta}^\star_{\text{persp}}$, especially when the perspective relaxation is nearly tight. Moreover, as we will observe empirically in Section \ref{sec:numres}, the cross-validated regressors $\bm{\beta}^\star_{\text{MIO}}$ and $\bm{\beta}^\star_{\text{persp}}$ behave extremely similarly whenever $n$ is sufficiently large. Therefore, since solving the perspective relaxation is much cheaper than solving a MIO, we use $\bm{\beta}^\star_{\text{persp}}$ in lieu of $\bm{\beta}^\star_{\text{MIO}}$ for our experiments involving confidence adjustment. However, we can easily construct a lower bound on $\bm{\beta}^\star_{\text{MIO}}$ by solving the following linear optimization problem{\color{black}
\begin{align*}
    \mu_h \geq \min_{\bm{v} \in \mathbb{R}^n, \bm{W} \in \mathbb{R}^{n \times n}}\max_{i \in [n]} \frac{1}{n} \sum_{j \in [n]}\vert v_i-W_{i,j}\vert \ \text{s.t.} \ \bm{l} \leq \bm{v} \leq \bm{u}, W_{i,j} \in [l_j^{(i)}, u_j^{(i)}] \ \forall i,j \in [n].
\end{align*}}
Moreover, we can construct an upper bound in much the same way after replacing the outer minimization operator with a maximization operator, by solving the following MIO{\color{black}
\begin{align*}
    \mu_h \leq \max_{\substack{\bm{v} \in \mathbb{R}^n, \bm{z} \in \{0, 1\}^{n\times n},\\ \bm{\theta}, \bm{W} \in \mathbb{R}^{n \times n}}} \quad & \max_{i \in [n]} \quad \frac{1}{n} \sum_{j \in [n]}\theta_{i,j} \\
    \text{s.t.} \ \ \ \quad\quad & \theta_{i,j} \geq \vert v_i-W_{i,j}\vert, \quad \forall i,j \in [n]\\
    & \theta_{i,j} \leq v_i-W_{i,j}+M (1-z_{i,j}), \theta_{i,j} \leq -v_i+W_{i,j}+M z_{i,j}\\
    & \bm{l} \leq \bm{v} \leq \bm{u}, W_{i,j} \in [l_j^{(i)}, u_j^{(i)}] \ \forall i,j \in [n],
\end{align*}}
or its Boolean relaxation. Finally, we remark that in some circumstances, we may be willing to evaluate $\bm{\beta}_{MIO}$ exactly, which involves solving one MIO, but not $\bm{\beta}^{(i)}_{MIO}$, which involves solving $n$ MIOs. In this situation, we bound $\mu_h$ almost identically, except the bounds on {\color{black}$v_i$}, which in spirit models $(y_i-\bm{x}_i^\top \bm{\beta}_{MIO})^2$, are now tight, and we bound $\mu_h$ from above by setting {\color{black}$W_{i,j}=u_j^{(i)}$} if {\color{black}$\vert u_j^{(i)}-v_i\vert \geq \vert l_j^{(i)}-v_i\vert$} and $l_j^{(i)}$ otherwise, rather than solving a MIO.
\fi


%In the rest of this section, we discuss how to leverage Theorem~\ref{prop:l10UB} to approximate and bound the confidence-adjusted LOOCV error $g$ without incurring the cost of solving $n$ MIOs. We treat the LOOCV error and the hypothesis stability terms in $g$ separately, as it is straightforward to bound $g$ given bounds on both these terms. We also propose some refinements to Theorem~\ref{prop:l10UB}, which sometimes allow us to obtain tighter bounds.


\subsection{\ag{Further Improvements for Lower Bounds}}\label{sec:genBounds}

Corollary~\ref{corr:perspUB} implies we may obtain a valid upper and lower bound on $h$ at a given hyperparameter combination $\gamma, \tau$ after solving $k$ perspective relaxations and computing $n$ terms of the form $$\sqrt{\bm{x_i^\top}\ag{\left(\bm{X}^{(\mathcal{N}_j)^\top} \bm{X}^{(\mathcal{N}_j)}+\frac{\gamma}{2}\mathbb{I}\right)^{-1}}\bm{x_i}}.$$ %Notably, the latter terms do not depend on $\gamma$ and $\tau$ and thus can be computed a priori before solving \eqref{eq:param_opt}.

A drawback of Corollary~\ref{corr:perspUB} is that if $\bm{x_i^\top}\bm{\beta}_{persp}^*\approx y_i$ for each $i \in \mathcal{N}_j$, i.e., the prediction of the perspective relaxation (without \ag{the $j$th fold}) is close to the response associated with point $i$, then Corollary~\ref{corr:perspUB}'s lower bound is $0$. \ag{A similar situation can happen with the stronger bounds for $h_j(\gamma,\tau)$ obtained from Theorem~\ref{prop:kfoldspread} and Problem \eqref{eq:boundComputation}.} We now propose a different bound on \ag{$h_j(\gamma,\tau)$}, which is sometimes effective in this circumstance.

First, define the function $f(\gamma,\tau)$ to be the in-sample training error \ag{without removing any folds and }with parameters $(\gamma, \tau)$,
\begin{equation*}
f(\gamma,\tau) := \frac{1}{n} \sum_{i=1}^n (y_i -\bm{x_i}^\top \bm{\beta}(\gamma, \tau))^2\quad
\text{s.t.} \quad  \bm{\beta}(\gamma, \tau) \in \argmin_{\bm{\beta} \in \mathbb{R}^p: \ \Vert \bm{\beta}\Vert_0 \leq \tau} \ \frac{\gamma}{2}\Vert \bm{\beta}\Vert_2^2 +\Vert \bm{X}\bm{\beta}-\bm{y}\Vert_2^2,
\end{equation*}
and let \ag{$f_{\mathcal{N}_j}(\gamma, \tau):=\sum_{i\in \mathcal{N}_j}\left(y_i -\bm{x_i}^\top \bm{\beta}(\gamma,\tau)\right)^2$} denote the training error \ag{associated with the $j$th fold}, with $\ag{1/n \sum_{j=1}^k f_{\mathcal{N}_j}(\gamma, \tau)}=f(\gamma, \tau)$. Observe that evaluating $h_{\mathcal{N}_j}(\gamma,\tau)$ involves solving \ag{$k$} MIOs, while evaluating $f$ requires solving one. 

\begin{proposition}\label{prop:LB}
For any $\gamma\geq 0$, any $\tau \in[p]$ \ag{and any $j\in [k]$}, $f_j(\gamma,\tau)\leq h_j(\gamma,\tau)$. Moreover, we have that $f(\gamma, \tau) \leq h(\gamma, \tau)$.
\end{proposition}
\proof{Proof of Proposition \ref{prop:LB}}
Given ${\color{black}j}\in [k]$, consider the following two optimization problems
	\begin{align}
		\min_{\bm{\beta}\in \mathbb{R}^p: \Vert \bm{\beta}\Vert_0 \leq \tau}\sum_{i=1}^n(y_i -\bm{x_i}^\top \bm{\beta})^2+\frac{\gamma}{2}\Vert \bm{\beta}\Vert_2^2\label{eq:prob_full}\\
	\min_{\bm{\beta}\in \mathbb{R}^p: \Vert \bm{\beta}\Vert_0 \leq \tau}\sum_{\ag{i \not\in \mathcal{N}_j}}(y_i -\bm{x_i}^\top \bm{\beta})^2+\frac{\gamma}{2}\Vert \bm{\beta}\Vert_2^2,\label{eq:prob_noK}
	\end{align}
let $\bm{\beta^*}$ be an optimal solution of \eqref{eq:prob_full}, and let $\bm{\beta^{\color{black}j}}$ be an optimal solution of \eqref{eq:prob_noK}. Since \begin{align*}&\sum_{\ag{i \not\in \mathcal{N}_j}}(y_i -\bm{x_i}^\top \bm{\beta^{\color{black}j}})^2+\frac{\gamma}{2}\Vert \bm{\beta^{\color{black}j}}\Vert_2^2\leq \sum_{\ag{i \not\in \mathcal{N}_j}}(y_i -\bm{x}^\top \bm{\beta^*})^2+\frac{\gamma}{2}\Vert \bm{\beta^*}\Vert_2^2,\quad \text{ and }\\
&\sum_{\ag{i \not\in \mathcal{N}_j}}(y_i -\bm{x}^\top \bm{\beta^{\color{black}j}})^2+\ag{\sum_{i\in \mathcal{N}_j}}(y_{\color{black}i} -\bm{x_{\color{black}i}}^\top \bm{\beta^{\color{black}j}})^2+\frac{\gamma}{2}\Vert \bm{\beta^{\color{black}j}}\Vert_2^2\geq \sum_{\ag{i \not\in \mathcal{N}_j}}(y_i -\bm{x_i}^\top \bm{\beta^*})^2+\ag{\sum_{i\in \mathcal{N}_j}}(y_{\color{black}i} -\bm{x_{\color{black}i}}^\top \bm{\beta^*})^2+\frac{\gamma}{2}\Vert \bm{\beta^*}\Vert_2^2,
\end{align*}
we conclude that $\ag{\sum_{i\in \mathcal{N}_j}}(y_{\color{black}i} -\bm{x_{\color{black}i}}^\top \bm{\beta^*})^2\leq  \ag{\sum_{i\in \mathcal{N}_j}}(y_{\color{black}i} -\bm{x_{\color{black}i}}^\top \bm{\beta^{\color{black}j}})^2$. The result immediately follows. \Halmos
\endproof


% \begin{corollary}\label{corr1}
% 	Given any $\gamma_1\leq \gamma_2$ and $\tau_1\geq \tau_2$, $f(\gamma_1,\tau_1)\leq f(\gamma_2,\tau_2)$ and thus $f(\gamma_1,\tau_1)\leq h(\gamma_2,\tau_2)$. 
% \end{corollary}
% % \proof{Proof of Corollary \ref{corr1}}
% % 	One can show that $f$ is non-decreasing in $\gamma$, non-increasing in $\tau$. \Halmos
% % \endproof 

Next, we develop a stronger bound on the $k$-fold error, by observing that our original proof technique relies on interpreting the optimal solution when training on the entire dataset as a feasible solution when leaving out the \ag{$j$}th fold, and that this feasible solution can be improved to obtain a tighter lower bound. Therefore, {given \ag{any} $\bm{z}\in \{0,1\}^p$,} let us define the function:{\color{black}
\begin{align*}
    f^{(\ag{\mathcal{N}_j})}(\bm{z}):=\min_{\bm{\beta} \in \mathbb{R}^p} \quad \frac{\gamma}{2}\sum_{j \in [p]}\beta_j^2+\Vert \bm{X}^{(\ag{\mathcal{N}_j})}\bm{\beta}-\bm{y}^{(\ag{\mathcal{N}_j})}\Vert_2^2 \ \text{s.t.} \ \beta_j=0 \ \text{if} \ z_j=0 \ \forall j \in [p],
\end{align*}}
to be the optimal training loss (including regularization) when we leave out the \ag{$j$th fold} and have the binary support vector $\bm{z}$. Then, fixing $\gamma, \tau$ and letting $u^*$ denote the optimal objective value of \eqref{eq:prob_full}, i.e., the optimal training loss on the entire dataset (including regularization) and $\bm{\beta}^{(\ag{\mathcal{N}_j})}(\bm{z})$ denote an optimal choice of $\bm{\beta}$ for this $\bm{z}$, we have the following result:
\begin{proposition}\label{prop2}
For any $\tau$-sparse binary vector $\bm{z}$, the following inequality holds:
\begin{align}
    u^*\leq f^{(\ag{\mathcal{N}_j})}(\bm{z})+\ag{\sum_{i\in \mathcal{N}_j}}\left(y_i-\bm{x}_i^\top \bm{\beta}^{(\ag{\mathcal{N}_j})}(\bm{z})\right)^2
\end{align}
\end{proposition}
\proof{Proof of Proposition \ref{prop2}}
The right-hand side of this inequality corresponds to the objective value of a feasible solution to \eqref{eq:prob_full}, while {\color{black}$u^\star$} is the optimal objective value of \eqref{eq:prob_full}. \Halmos
\endproof
\begin{corollary}\label{corr:lowerbound2}
Let $\bm{z}$ denote a $\tau$-sparse binary vector. Then, we have the following bound on the $j$th partial \ag{cross-validation} error:
\begin{align}\label{eq:LOOLower}
    h_j(\gamma, \tau) \geq u^*-f^{(\ag{\mathcal{N}_j})}(\bm{z}).
\end{align}
\end{corollary}
\proof{Proof of Corollary \ref{corr:lowerbound2}}
The right-hand side of this bound is maximized by setting $\bm{z}$ to be a binary vector which minimizes $f^{(\ag{\mathcal{N}_j})}(\bm{z})$, and therefore this bound is valid for any $\bm{z}$. \Halmos
\endproof
We close this section with two remarks:
\begin{remark}[Bound quality]
Observe that bound~\eqref{eq:LOOLower} is at least as strong as $f_j(\gamma, \tau)$ with $\bm{z}$ encoding an optimal choice of support in \eqref{eq:prob_full}. Indeed, if $\bm{\beta}^{(\ag{\mathcal{N}_j})}(\bm{z})$ solves \eqref{eq:prob_full}, then both bounds agree and equal $h_j(\gamma, \tau)$ but otherwise \eqref{eq:LOOLower} is strictly stronger. Moreover, since $f_j(\gamma, \tau)$ is typically nonzero, then the bound \eqref{eq:LOOLower} is positive as well and can improve upon the lower bound in \eqref{eqn:perspbounds}. Finally, it is easy to construct an example where the lower bound in \eqref{eqn:perspbounds} is stronger than \eqref{eq:LOOLower}, thus neither lower bound dominates the other.
\end{remark}
\begin{remark}[Computational efficiency]
	Computing lower bound \eqref{eq:LOOLower} for each \ag{$j\in [k]$} requires solving at least one MIO, corresponding to \eqref{eq:prob_full}, which is a substantial improvement over the \ag{$k$} MIOs required to compute $h$ but may still be an expensive computation. However, using any lower bound on $u^*$, for example, corresponding to the optimal solution of a perspective relaxation, gives valid lower bounds. Therefore, in practice, we suggest using a heuristic instead to bound $h_j$ from below, e.g., rounding a perspective relaxation.
\end{remark}

% {\color{black}
% \begin{remark}[Generality of Approach]
% It is noting that the bounding approach developed here is quite general: it leverages the strong convexity of the lower-level problem and results from Fenchel duality to construct the bounds $\underline{\xi}, \overline{\xi}$, rather than any properties specific to regression. Thus, it is straightforward to extend our approach to develop upper/lower bounds on validation metrics in other contexts with strongly convex lower-level objectives, e.g., $\ell_0-\ell_2^2$ logistic regression \citep{deza2022safe}. It would be interesting to investigate whether our approach can be applied in the absence of a strongly convex objective, e.g., optimal decision trees or neural networks.
% \end{remark}
% }


\section{Optimizing the Cross-Validation Loss}\label{sec:corddescent}
In this section, we present an efficient coordinate descent scheme that identifies (approximately) optimal hyperparameters $(\gamma, \tau)$ with respect to the metric: 
\begin{equation}\label{prob:cyclic1}\color{black}
	h(\gamma, \tau):= \frac{1}{n} \sum_{j \in [k]} h_j(\gamma, \tau)
\end{equation}
by iteratively minimizing $\tau$ and $\gamma$. In the tradition of coordinate descent schemes, with initialization $\tau_0, \gamma_0$, we repeatedly solve the following two optimization problems:
\begin{align}\label{prob:cyclic2a}
    \tau_t \in \argmin_{\tau \in[p]} \quad h(\gamma_t, \tau),\\
    \gamma_{t+1} \in \arg \min_{\gamma >0 } \quad h(\gamma, \tau_t),\label{prob:cyclic2b}
\end{align}
until we either detect a cycle or converge to a locally optimal solution. To develop this scheme, in Section \ref{ssec:parametric} we propose an efficient technique for solving Problem {\color{black}\eqref{prob:cyclic2a}}, %i.e., optimizing $\tau$ with a fixed $\gamma$ by leveraging techniques from perspective relaxations and sensitivity analysis, 
and in Section \ref{ssec:parametric2} we propose an efficient technique for (approximately) solving Problem {\color{black}\eqref{prob:cyclic2b}}. Accordingly, our scheme could also be used to identify an optimal choice of $\gamma$ if $\tau$ is already known, e.g., in a context where regulatory constraints specify the number of features that may be included in a model. 

Our overall approach is motivated by three key observations. First, we design a method that obtains local, rather than global, minima, because $h$ is a highly non-convex function and even evaluating $h$ requires solving $n$ MIOs, which suggests that global minima of $h$ may not be attainable in a practical amount of time at scale. Second, we use coordinate descent to seek local minima because if either $\tau$ or $\gamma$ is fixed, it is possible to efficiently optimize the remaining hyperparameter with respect to $h$ by leveraging the convex relaxations developed in the previous section. Third, we should expect our coordinate descent scheme to perform well in practice, because similar schemes are highly effective in other machine learning contexts, e.g., solving certain matrix completion problems in polynomial time \citep{mazumder2011sparsenet, cifuentes2022polynomial}.

%We remark that in a subset of our numerical experiments, we require that the optimal values of $\gamma$ and $\tau$ should be contained within the intervals $[\gamma_{\min}, \gamma_{\max}]$ and $[\tau_{\min}, \tau_{\max}]$. For instance, we require that $\tau_{\max}$ is such that $\tau_{\max}\log \tau_{\max} \leq n$ in our experiments in Section \ref{ssec:realworldata}, because \citet[Theorem 2.5]{gamarnik2022sparse} demonstrated that, under certain assumptions on the data generation process, on the order of $\tau\log \tau$ observations are needed to recover a model with sparsity $\tau_{\max}$.


\subsection{Parametric Optimization of $k$-fold With Respect to Sparsity}\label{ssec:parametric}

Consider the following optimization problem, where $\gamma$ is fixed here and throughout this subsection:
{\color{black}
\begin{equation}\label{eq:paramK}
\begin{aligned}
\min_{\tau\in [p]} \quad h({\gamma}, \tau):=\min_{} & \sum_{j \in [k]}\sum_{i \in \mathcal{N}_j} (y_i-\bm{x}_i^\top \bm{\beta}^{(\mathcal{N}_j)})^2,\\
\text{s.t.} \quad  & \bm{\beta}^{(\mathcal{N}_j)} \in \argmin_{\bm{\beta} \in \mathbb{R}^p: \ \Vert\bm{\beta}\Vert_0\leq \tau} \ \frac{\gamma}{2}\Vert \bm{\beta}\Vert_2^2 +\Vert \bm{X}^{(\mathcal{N}_j)}\bm{\beta}-\bm{y}^{(\mathcal{N}_j)}\Vert_2^2\quad \forall i \in [n].\\[2em]
\end{aligned}
\end{equation}}
% , and $\mu_h$ denotes the hypothesis stability.
This problem can be solved by complete enumeration, i.e., for each $\tau \in[p]$, we compute an optimal $\beta^{(\mathcal{N}_j)}$ for each $j\in [k]$ by solving an MIO. This involves solving $(k+1)p$ MIOs, which is extremely expensive at scale. We now propose a technique for minimizing $h$ without solving all these MIOs: 
% \begin{align}
% \text{Let} \quad h_{\mathcal{N}_j}(\gamma, \tau):= \sum_{i \in \mathcal{N}_j}(y_i -\bm{x_i}^\top \bm{\beta})^2\quad
% 	\text{s.t.} \quad  \bm{\beta} \in \argmin_{\bm{\beta} \in \mathbb{R}^p: \ \Vert\bm{\beta}\Vert_0\leq \tau} \ \frac{\gamma}{2}\Vert \bm{\beta}\Vert_2^2 +\Vert \bm{X}^{(\mathcal{N}_j)}\bm{\beta}-\bm{y}^{(\mathcal{N}_j)}\Vert_2^2,
% \end{align}
% and note that $h(\gamma, \tau)=1/n \sum_{j \in [k]} h_{\mathcal{N}_j}(\gamma, \tau)$. In the special case of leave-one-out cross-validation, we denote $h_{\mathcal{N}_j}$ by $h_i$. From Proposition~\ref{prop:LB}, we find that $f_{\mathcal{N}_j}(\gamma, \tau)\leq h_{\mathcal{N}_j}(\gamma, \tau)$ with 
% \begin{align}
% f_{\mathcal{N}_j}(\gamma, \tau) = \sum_{i \in \mathcal{N}_j}(y_i -\bm{x_i}^\top \bm{\beta})^2\quad
% \text{s.t.} \quad  \bm{\beta} \in \argmin_{\bm{\beta} \in \mathbb{R}^p: \ \Vert\bm{\beta}\Vert_0\leq \tau} \ \frac{\gamma}{2}\Vert \bm{\beta}\Vert_2^2 +\Vert \bm{X}\bm{\beta}-\bm{y}\Vert_2^2.
% \end{align}

% Combining these definitions and observations allows us to propose a method that minimizes $g(\gamma, \tau)$ with respect to $\tau$ without solving $(k+1)p$ MIOs, which we formalize\footnote{We omit some details around bounding the hypothesis stability for the sake of brevity; these bounds can be obtained in much the same way as the bounds on the LOOCV error, with a minor implementation detail around whether to compute $\bm{\beta}(\tau)$, an optimal regressor with no data points left out and a sparsity constraint of $\tau$, before or after running the method; the method performs well in either case.} in Algorithm \ref{alg:parametricK}. 

{Algorithm \ref{alg:parametricK} has two main phases, which both run in a loop. In the first phase, we construct valid lower and upper bounds on $h_{\mathcal{N}_j}(\tau)$ for each $\mathcal{N}_j$ and each $\tau$ without solving any MIOs. We begin by solving, for each potential sparsity budget $\tau\in [p]$, the perspective relaxation with all datapoints included. Call this relaxation's objective value $\bar v_{\tau}$. We then solve each perspective relaxation that arises after omitting one data fold $\mathcal{N}_j: j \in [k]$, with objective values $v_{\tau,\mathcal{N}_j}$ and solutions $\bm{\beta}_{\tau,\mathcal{N}_j}$. %The later solutions produce estimates $h_{\mathcal{N}_j}(\tau)$ of the true $k$-fold error associated with fold $\mathcal{N}_j$ and cardinality $\tau$. 
Next, we compute lower and upper bounds on the $k$-fold error $h_{\mathcal{N}_j}(\tau)$ using the methods derived in Section~\ref{sec:lowerbounds}, which are summarized in the routine \texttt{compute\_bounds} described in Algorithm~\ref{alg:bounds}. By solving $\mathcal{O}(kp)$ relaxations (and no MIOs), we have upper and lower estimates on the $k$-fold error that are often accurate in practice, as described by Theorem~\ref{prop:l10UB}. 

After completing the first loop in Algorithm~\ref{alg:parametricK}, one may already terminate the algorithm. Indeed, according to our numerical experiments in Section \ref{sec:numres}, this already provides high-quality solutions. Alternatively, one may proceed with the second phase of Algorithm~\ref{alg:parametricK} and solve \eqref{prob:cyclic2a} to optimality, at the expense of solving (a potentially large number of) MIOs.

In the second phase, Algorithm~\ref{alg:parametricK} identifies the cardinality $\tau^*$ with the best lower bound (and thus, in an optimistic scenario, the best potential value). Then, it identifies the fold $\mathcal{N}_j^*$ with the largest uncertainty around the $k$-fold estimate $h_{\mathcal{N}_j^*}(\tau^*)$, and solves an MIO to compute the exact partial $k$-fold error. %Observe that the choice of $\mathcal{N}_j^*$ always results in the largest reduction of gap from one iteration to the next. 
This process is repeated until \eqref{eq:paramK} is solved to provable optimality, or a suitable termination condition (e.g., a limit on computational time) %or the number of MIOs solved)
is met. %This is analogous to the pure-explore strategy in multi-armed bandit optimization \cite{bubeck2009pure}, and the best-first branching strategy in branch-and-bound \citep{wolsey2014integer}.
}

%By leveraging the lower and upper bounds on each $h_i(\gamma, \tau)$ and on $\mu_h$ from perspective relaxations we developed in Section \ref{sec:genBounds}--\ref{ssec:genbounds2}, as summarized in Algorithm \ref{alg:bounds}, we immediately have bounds on the LOOCV error for each $\tau$. The idea is then to combine these bounds with our bounds on the hypothesis stability $\mu_h$ and tighten the bounds sequentially, which can be done by selecting the cardinality $\tau$ with the smallest lower bound and solving a MIO corresponding to the index $i$ such that the gap between the lower and upper bounds on $h_i(\gamma_t, \tau)$ from Algorithm \ref{alg:bounds} is as large as possible. If, at each step, one selects the index $\tau$ with the best lower bound, one can significantly reduce the number of MIOs that need to be solved. 

\normalem


\RestyleAlgo{ruled}
\begin{algorithm}[h!] 
	\caption{Computing optimal sparsity parameter for $k$-fold error}
	\label{alg:parametricK}
	\KwData{$\gamma$: $\ell_2^2$ regularization parameter; $\epsilon>0$: desired optimality tolerance; $r$: budget on number of MIOs}
	\KwResult{Cardinality with best estimated $k$-fold error} % Up to here.
        \For{$\tau \in[p]$}{
        $\bar v_{\tau} \gets \min_{\bm{\beta}\in \mathbb{R}^p,\bm{z}\in [0,1]^p}\;\|\bm{X\beta}-\bm{y}\|_2^2+\frac{\gamma}{2}\sum_{i=1}^p\beta_i^2/z_i \text{  s.t.  }\bm{e}^\top \bm{z}\leq \tau$\\
        %$\bar{u}_k, \bm{z}^\tau\gets$ objective value with sparsity $\tau$, binary support of $\bm{\beta}$\\
			\For{$j\in [k]$}{
                $v_{\tau,\mathcal{N}_j} \gets\min_{\bm{\beta}\in \mathbb{R}^p,\bm{z}\in [0,1]^p}\;\|\bm{X}^{(\mathcal{N}_j)}\bm{\beta}-\bm{y}^{(\mathcal{N}_j)}\|_2^2+\frac{\gamma}{2}\sum_{i=1}^p\beta_i^2/z_i \text{  s.t.  }\bm{e}^\top \bm{z}\leq \tau$\\
                $\bm{\beta}_{\tau,\mathcal{N}_j} \in \argmin_{\bm{\beta}\in \mathbb{R}^p,\bm{z}\in [0,1]^p}\;\|\bm{X}^{(\mathcal{N}_j)}\bm{\beta}-\bm{y}^{(\mathcal{N}_j)}\|_2^2+\frac{\gamma}{2}\sum_{i=1}^p\beta_i^2/z_i \text{  s.t.  }\bm{e}^\top \bm{z}\leq \tau$\\
                $h_{\mathcal{N}_j}(\tau) \gets \sum_{i \in \mathcal{N}_j}(y_i-\bm{x_i^\top \beta_{\tau,i}})^2$ \tcp*{Perspective sol. estimates $k$-fold for $\mathcal{N}_j$}
                $u_{\tau,\mathcal{N}_j} \gets \texttt{round}(\bm{\beta}_{\tau,\mathcal{N}_j})$ \tcp*{Any heuristic can be used}
                $\zeta_{\mathcal{N}_j}^L(\tau),\zeta_{\mathcal{N}_j}^U(\tau)\gets \texttt{compute\_bounds}(\mathcal{N}_j,\bm{\beta}_{\tau,\mathcal{N}_j},\bar v_{\tau},v_{\tau,\mathcal{N}_j},u_{\tau,\mathcal{N}_j})$
		}
  %$\bm{l}, \bm{u}, \bm{l}^{(\mathcal{N}_j)}, \bm{u}^{(\mathcal{N}_j)} \gets \text{compute\_bounds}$ \tcp*{Similarly to Algorithm \ref{alg:bounds}}\
   % $\underline{\mu_h}(\tau) \gets \min_{\substack{\bm{x} \in \mathbb{R}^n, \\\bm{Y} \in \mathbb{R}^{n \times n}}}\max_{i \in [n]} \frac{1}{n} \sum_{j \in [n]}\vert x_i-Y_{i,j}\vert$\; 
  %  $\overline{\mu_h}(\tau) \gets \max_{\substack{\bm{x} \in \mathbb{R}^n, \\\bm{Y} \in \mathbb{R}^{n \times n}}}\max_{i \in [n]} \frac{1}{n} \sum_{j \in [n]}\vert x_i-Y_{i,j}\vert \ $\ 
 %   $\text{s.t.} \ \bm{l} \leq \bm{x} \leq \bm{u}, Y_{i,j} \in [l_i^{(\mathcal{N}_j)}, u_i^{(\mathcal{N}_j)}] \ \forall i \in [n], j \in [k]: \ i \in \mathcal{N}_j \ \text{for optimization problems } \overline{\mu_h}, \underline{\mu_h}$\tcp*{For max problem, solve relaxation of MIO in practice}
	}
$LB\gets \min_{\tau \in[p]}\sum_{j \in [k]} \zeta_{\mathcal{N}_j}^L(\tau)$; 
$UB\gets \min_{\tau \in[p]}\sum_{j \in [k]} \zeta_{\mathcal{N}_j}^U(\tau)$ \tcp*{Bounds on $k$-fold} 
        $num\_mip\gets 0$\\
	\Repeat{$(UB-LB)/UB\geq \epsilon$ or $num\_mip>r$}{
		$\tau^*\gets{\arg \min}_{\tau\in [p]} \sum_{i=1}^n \zeta_{\mathcal{N}_j}^L(\tau)$ \tcp*{Cardinality with best bound}
		$\mathcal{N}_j^*\gets \arg \max_{j\in [k]} \{\zeta_{\mathcal{N}_j}^U(\tau)-\zeta_{\mathcal{N}_j}^L(\tau)\}$ \tcp*{Fold with largest $k$-fold uncertainty}
            $h_{\mathcal{N}_j^*}(\tau^*)\gets \min_{\bm{\beta}\in \mathbb{R}^p,\bm{z}\in \{0,1\}^p}\;\|\bm{X}^{(\mathcal{N}_j^*)}\bm{\beta}-\bm{y}^{(\mathcal{N}_j^*)}\|_2^2+\frac{\gamma}{2}\|\bm{\beta}\|_2^2 \text{  s.t.  }\bm{e}^\top \bm{z}\leq \tau^*$\tcp*{Solve MIO}
            $\zeta_{\mathcal{N}_j^*}^L(\tau^*)\gets h_{\mathcal{N}_j^*}(\tau^*),\; \zeta_{\mathcal{N}_j^*}^U(\tau^*)\gets h_{\mathcal{N}_j^*}(\tau^*)$\\	
            $LB\gets \min_{\tau \in[p]}\sum_{j \in [k]} \zeta_i^L(\tau)$\\
            $UB\gets \min_{\tau \in[p]}\sum_{j \in [k]} \zeta_{\mathcal{N}_j}^U(\tau)$\\
            $num\_mip\gets num\_mip+1$\\[0.5em]
	}
        \Return{${\arg \min}_{\tau\in [p]} \sum_{j \in [k]} h_{\mathcal{N}_j}(\tau)$}\tcp*{Cardinality with best error}
	%\Return{${\arg \min}_{\tau\in [p]} \sum_{i=1}^n \zeta_i^U(\tau)$}\tcp*{Cardinality with best guaranteed error}
\end{algorithm}

\begin{algorithm}[h!] 
	\caption{$\texttt{compute\_bounds}(\mathcal{N}_j,\bm{\beta},\bar v,v,u)$}
	\label{alg:bounds}
	\KwData{$\mathcal{N}_j$: fold left out; $\bm{\beta}$: optimal solution of perspective relaxation with $\mathcal{N}_j$ left out; $\bar v$: lower bound of obj val of MIO with all data; $v$: optimal obj value of perspective relaxation with $\mathcal{N}_j$ left out; $u$: upper bound of obj val of MIO with $\mathcal{N}_j$ left out }
	\KwResult{Lower and upper bounds on the $k$-fold error attributable to datapoint $i$}
 $\underline{\xi}\gets\bm{x_i^\top\beta}-\sqrt{\bm{x_i^\top}\left(\bm{X}^{(\mathcal{N}_j)^\top} \bm{X}^{(\mathcal{N}_j)}+\frac{\gamma}{2}\mathbb{I}\right)^{-1}\bm{x_i}\left(u-v\right)}$\\
$\overline{\xi}\gets\bm{x_i^\top\beta}+\sqrt{\bm{x_i^\top}\left(\bm{X}^{(\mathcal{N}_j)^\top} \bm{X}^{(\mathcal{N}_j)}+\frac{\gamma}{2}\mathbb{I}\right)^{-1}\bm{x_i}\left(u-v\right)}$\\
  $\zeta^L\gets \bar v-u,\; \zeta^U\gets \max\{(y_i-\underline{\xi})^2, (\overline{\xi}-y_i)^2\}$\\
  \If{$\underline{\xi}> y_i$}{
   $\zeta^L\gets \max\{\zeta^L, (\underline{\xi}-y_i)^2\}$
  }
  \If{$\overline{\xi}< y_i$}{
   $\zeta^L\gets \max\{\zeta^L, (y_i-\overline{\xi})^2\}$
  }
	\Return{$\left(\zeta^L, \zeta^U\right)$}
\end{algorithm}

To solve each MIO in Algorithm \ref{alg:parametricK}, we invoke a Generalized Benders Decomposition scheme \citep{geoffrion1972generalized}, which was specialized to sparse regression problems by \cite{bertsimas2020sparse}, enhanced with some ideas from the optimization literature summarized in the works \cite{bertsimas2020sparse2, hazimeh2020fast}. For the sake of conciseness, we defer these implementation details to Appendix \ref{sec:append.implementationdetails}.

%One could also store screening rules for each $\tau$ and each fold of the data in memory, although this is more memory intensive.

\paragraph{Algorithm \ref{alg:parametricK} in Action: } Figure \ref{fig:algvarykinaction} depicts visually the lower and upper bounds on $g$ from Algorithm \ref{alg:bounds} (left) and after running Algorithm \ref{alg:parametricK} to completion (right) on a synthetic sparse regression instance generated in the fashion described in our numerical experiments, with $k=n$, $n=200, p=20$, $\gamma=1/\sqrt{n}$, $\tau_{\text{true}}=10$, $\rho=0.7$, $\nu=1$, where $\tau \in\{2, \ldots, 19\}$, and using the outer-approximation method of \cite{bertsimas2020sparse} as our solver for each MIO with a time limit of $60$s. We observe that Algorithm \ref{alg:parametricK} solved $1694$ MIOs to identify the optimal $\tau$, which is a 53\% improvement on complete enumeration. Interestingly, when $\tau=19$, the perspective relaxation is tight after omitting any fold of the data and we have tight bounds on the LOOCV error without solving any MIOs. In our computational experiments, see Section \ref{sec:compExactLOOCV}, we test Algorithm~\ref{alg:parametricK} on real datasets and find that it reduces the number of MIOs that need to be solved by 50-80\% with respect to complete enumeration. For more information on how the bounds evolve over time, we provide a GIF with one frame each time a MIO is solved at the link \url{https://drive.google.com/file/d/1EZdNwlV9sEEnludGGM7v2nGpB7tzZvz4/view?usp=sharing}.

\begin{figure}[h!]\centering
\begin{subfigure}[t]{.45\linewidth}
\includegraphics[scale=0.35]{fig2_mspaper_iter1.pdf}
\end{subfigure}
\begin{subfigure}[t]{.45\linewidth}
\includegraphics[scale=0.35]{fig2_mspaper_iterx.pdf}
\end{subfigure}
\caption{Comparison of initial bounds on LOOCV ($k$-fold with $k=n$) from Algorithm \ref{alg:bounds} (left) and bounds after running Algorithm \ref{alg:parametricK} (right) for a synthetic sparse regression instance where $p=20, n=200, \tau_{\text{true}}=10$, for varying $\tau$. The black number in the top middle depicts the iteration number of the method.}
\label{fig:algvarykinaction}
\end{figure}

% \paragraph{Minimizing the Confidence-Adjusted $k$-fold Error at Scale:} When minimizing $g$ with respect to $\tau$, Algorithm \ref{alg:parametricK} performs well in settings where $p < 1000$, often reducing the number of MIOs that need solving by $70\%$ or more compared to complete enumeration, as we demonstrated in Table \ref{fig:reductionMioNodes}. Unfortunately, in high-dimensional settings, this is not enough of an improvement over complete enumeration to be numerically tractable, due to the number of MIOs that need solving. Accordingly, in certain settings, we eschew Algorithm \ref{alg:parametricK} and instead use the more tractable saddle-point method of \cite{bertsimas2020sparse2} to estimate $\bm{\beta}_{\text{persp}}^{(\mathcal{N}_j)}$. Notably, this saddle-point method recovers the same regressors as exact methods when $n \gg p$ \citep[see][for precise conditions]{bertsimas2020sparse2}. %Moreover, as we demonstrate in Sections \ref{sec:genBounds}-\ref{ssec:genbounds2}, this estimator provides a similar, although slightly different, confidence-adjusted LOOCV error to the MIO, and can therefore be considered a heuristic with the same limiting behavior as Algorithm \ref{alg:parametricK}.

\subsection{Parametric Optimization of $k$-fold Error With Respect to $\gamma$} \label{ssec:parametric2}
In this section, we propose a technique for approximately minimizing the $k$-fold error with respect to the regularization hyperparameter $\gamma$. 

We begin with two observations from the literature. First, as observed by \cite{stephenson2021can}, the LOOCV error $h(\gamma,\tau)$ is often quasi-convex with respect to $\gamma$ when $\tau=p$. Second, \cite{bertsimas2021unified, bertsimas2022scalable} reports that, for sparsity-constrained problems, the optimal support does not often change as we vary $\gamma$. Combining these observations suggests that, after optimizing $\tau$ with $\gamma$ fixed, a good strategy for minimizing $g$ with respect to $\gamma$ is to fix the optimal support $\bm{z}^{(\mathcal{N}_j)}$ with respect to each fold $i$ and invoke a root-finding method to find a $\gamma$ which locally minimizes %\footnote{We remark that one could also generate a grid of allowable values of $\gamma$ and proceed in the same manner as in the previous section, although this would be less numerically tractable and may not even generate locally optimal $\gamma$'s.} 
$g$.


Accordingly, we now use the fact that $\gamma$ and $\bm{z}^{(\mathcal{N}_j)}$ fully determine $\bm{\beta}^{(\mathcal{N}_j)}$ to rewrite
\begin{align*}
    \min_{\bm{\beta} \in \mathbb{R}^p} \quad & \frac{\gamma}{2}\Vert \bm{\beta}\Vert_2^2+\Vert \bm{X}\bm{\beta}-\bm{y}\Vert_2^2 \ \text{s.t.} \ \beta_i =0 \ \text{if} \ \hat{z}_i=0,\\
\text{as}\qquad &  \bm{\beta}^\star=\left(\frac{\gamma}{2}\mathbb{I}+\bm{X}^\top \mathrm{Diag}(\hat{\bm{z}})\bm{X}\right)^{-1}\mathrm{Diag}(\hat{\bm{z}})\bm{X}^\top \bm{y}.
\end{align*}
Therefore, we fix each $\bm{z}^{(\mathcal{N}_j)}$ and substitute the resulting expressions for each $\bm{\beta}^{(\mathcal{N}_j)}$ into the $k$-fold error. This substitution yields the following univariate optimization problem, which can be solved via standard root-finding methods to approximately minimize the $k$-fold loss:
\begin{align}\label{prob:newton1}
            \min_{\gamma >0} \sum_{j \in [k]}\sum_{i \in \mathcal{N}_j} \left(y_i-\bm{X}_i^\top \mathrm{Diag}(\bm{z}^{(\mathcal{N}_j)}) \left(\frac{\gamma}{2} \mathbb{I}+\bm{X}^{(i)\top} \mathrm{Diag}(\bm{z}^{(\mathcal{N}_j)})\bm{X}^{(i)^\top}\right)^{-1}\mathrm{Diag}(\bm{z}^{(\mathcal{N}_j)})\bm{X}^{{(i)}^\top} \bm{y}^{(i)}\right)^2.
\end{align}
% Moreover, if $\delta >0$ and we are interested in minimizing the confidence-adjusted $k$-fold error, rather than the $k$-fold error itself, we assume that the index $j$ at which the expression 
% \begin{align*}
% \mu_h:=\max_{j\in [k]}\frac{1}{n}\sum_{j=1}^n \left\vert (y_i-{\bm{\beta}^\star}^\top \bm{x}_i)^2-(y_i-\bm{\beta}^{(\mathcal{N}_j)} \bm{x}_i)^2\right\vert
% \end{align*}
% attains its maximum\footnote{We pick the first index $j$ which attains this maximum in the rate case of ties.}, does not vary as we vary $\gamma$. Fixing $j$ then allows us to derive a similar approximation for the hypothesis stability, namely:
% \begin{align*}
%     \mu_h(\gamma, \tau)\approx \frac{1}{n}\sum_{i \in [n]} & \bigg\vert \left(y_i-\bm{X}_i^\top \mathrm{Diag}(\bm{z}^{(j)}) \left(\frac{\gamma}{2}  \mathbb{I}+\bm{X}^{(j)\top} \mathrm{Diag}(\bm{z}^{(j)})\bm{X}^{(j)^\top}\right)^{-1}\mathrm{Diag}(\bm{z}^{(j)})\bm{X}^{{(j)}^\top} \bm{y}^{(j)}\right)^2\\
%    &  - \left(y_i-\bm{X}_i^\top \mathrm{Diag}(\bm{z}) \left(\frac{\gamma}{2} \mathbb{I}+\bm{X}^{\top} \mathrm{Diag}(\bm{z}^{})\bm{X}^{\top}\right)^{-1}\mathrm{Diag}(\bm{z}^{})\bm{X}^{{}^\top} \bm{y}^{}\right)^2 \bigg\vert,
% \end{align*}
% where $\bm{z}$ denotes the optimal support when no data observations are omitted. With these expressions, it is straightforward to minimize the confidence-adjusted LOOCV error with respect to $\gamma$. 
Details on minimizing $\gamma$ using \verb|Julia| are provided in Appendix \ref{append.4.2}.

\section{Numerical Experiments}\label{sec:numres} 

We now present numerical experiments testing our proposed methods. First, in Section~\ref{sec:compExactLOOCV}, we study the computational savings of using Algorithm~\ref{alg:parametricK} over a complete grid search when optimizing the $k$-fold error, as a function of the sparsity parameter $\tau$. Then, in Section~\ref{sec:comp_real}, we compare the statistical performance of the coordinate descent algorithm proposed in Section~\ref{sec:corddescent} with alternatives in the literature.
%Finally, in Section \ref{sec:syntStats2}, we combine our proposed approach with a confidence adjustment procedure proposed in a separate work \citep{cory2025stability} and demonstrate that, in conjunction with the confidence adjustment procedure, our approach also leads to lower test set errors than via other methods in the literature.



\subsection{Exact K-fold Optimization}\label{sec:compExactLOOCV}

We first assess whether Algorithm~\ref{alg:parametricK} significantly reduces the number of MIOs that need to be solved to minimize the kCV error with respect to $\tau$, compared to grid search. We set either $k=n$ or $k=10$, corresponding to leave-one-out and ten-fold cross-validation problems \eqref{eq:paramK} respectively.

We compare the performance of two approaches. First, a standard grid search approach (\texttt{Grid}), where we solve the inner MIO in \eqref{eq:paramK} for all combinations of cardinality $\tau\in [p]$ and all folds of the data $j\in [k]$, and select the hyperparameter combination which minimizes the objective. To ensure the quality of the resulting solution, we solve all MIOs to optimality (without any time limit). Second, we consider using Algorithm~\ref{alg:parametricK} with parameter $r=\infty$ (thus solving MIOs to optimality until the desired optimality gap $\epsilon$ for problem~\eqref{eq:paramK} is proven). We test regularization parameter $\gamma\in \{0.01,0.02,0.05,0.10,0.20,0.50,1.00\}$ in Algorithm~\ref{alg:parametricK}, and solve all MIOs via their perspective reformulations, namely
\begin{align*}
	\min_{\bm{\beta}\in \mathbb{R}^p,\bm{z}\in \{0,1\}^p}\;& \|\bm{X}\bm{\beta}-\bm{y}\|_2^2+\frac{\gamma}{2}\sum_{j=1}^p \frac{\beta_j^2}{z_j}\ \text{s.t.}\ \sum_{j=1}^p z_j\leq \tau,
\end{align*} 
using Mosek 10.0. Since the approach \texttt{Grid} involves solving $\mathcal{O}(kp)$ MIOs (without a time limit), we are limited to testing these approaches on small datasets, and accordingly use the Diabetes, Housing, Servo, and AutoMPG datasets for this experiment, as described by \cite{gomez2021mixed}. Moreover, we remark that the specific solution times and the number of nodes expanded by each method are not crucial, as those could vary substantially if relaxations other than the perspective are used, different solvers or solution approaches are used, or if advanced techniques are implemented (but both methods would be affected in the same way). Thus, we focus our analysis on relative performance.

Figures~\ref{fig:reductionMioNodes} and~\ref{fig:reductionMioNodes10} summarize the percentage reduction of the number of MIOs and the number of branch-and-bound nodes achieved by Algorithm~\ref{alg:parametricK} over \texttt{Grid}, computed as 
\small$$ \text{Reduction in MIOs}=\frac{\text{\# MIO}_{\texttt{Grid}}-\text{\# MIO}_{\text{Alg.\ref{alg:parametricK}}}}{\text{\# MIO}_{\texttt{Grid}}},\;\; \text{Reduction in nodes}=\frac{\text{\# nodes}_{\texttt{Grid}}-\text{\# nodes}_{\text{Alg.\ref{alg:parametricK}}}}{\text{\# nodes}_{\texttt{Grid}}},$$\normalsize
where $\text{\# MIO}_{\text{Y}}$ and $\text{\# nodes}_{\text{Y}}$ indicate the number of MIOs or branch-and-bound nodes used by method Y.  Tables \ref{tab:gridVsAlg} and \ref{tab:gridVsAlg10} present the detailed computational results. 

\begin{figure}[h!]\centering
\begin{subfigure}[t]{.45\linewidth}
\includegraphics[width=\textwidth,trim={12cm 6cm 12cm 6cm},clip]{mio_reduction.pdf}
\end{subfigure}
\begin{subfigure}[t]{.45\linewidth}
\includegraphics[width=\textwidth,trim={12cm 6cm 12cm 6cm},clip]{nodes_reduction.pdf}
\end{subfigure}
\caption{Reduction in the number of MIO solved (left) and the total number of branch-and-bound nodes (right) when using Algorithm~\ref{alg:parametricK} for leave-one-out cross-validation, when compared with \texttt{Grid} (i.e., independently solving $\mathcal{O}(pn)$ MIOs) in four real datasets. The distributions shown in the figure correspond to solving the same instance with different values of $\gamma$. All MIOs are solved to optimality, without imposing any time limits. } % Old version-running another version at the moment which will be larger-scale, since this version only halved the number of MIOs that needed solving.
\label{fig:reductionMioNodes}
\end{figure}

\begin{figure}[h!]\centering
\begin{subfigure}[t]{.45\linewidth}
\includegraphics[width=\textwidth,trim={12cm 6cm 12cm 6cm},clip]{mio_reduction10.pdf}
\end{subfigure}
\begin{subfigure}[t]{.45\linewidth}
\includegraphics[width=\textwidth,trim={12cm 6cm 12cm 6cm},clip]{nodes_reduction10.pdf}
\end{subfigure}
\caption{Reduction in the number of MIO solved (left) and the total number of branch-and-bound nodes (right) when using Algorithm~\ref{alg:parametricK} for 10-fold cross-validation, when compared with \texttt{Grid} (i.e., independently solving $\mathcal{O}(pk)$ MIOs) in four real datasets. The distributions shown in the figure correspond to solving the same instance with different values of $\gamma$. All MIOs are solved to optimality, without imposing any time limits. } % Old version-running another version at the moment which will be larger-scale, since this version only halved the number of MIOs that needed solving.
\label{fig:reductionMioNodes10}
\end{figure}


	\begin{table}
    [h!]\footnotesize
    \begin{center}
 \caption{Comparison between using Algorithm~\ref{alg:parametricK} and solving $\mathcal{O}(pn)$ MIOs independently (\texttt{Grid}) for leave-one-out cross-validation in four real datasets, for different values of regularization $\gamma$. Times reported are in minutes, and correspond to the time to solve all required mixed-integer optimization problems to optimality. No time limits are imposed on the MIOs. Algorithm~\ref{alg:parametricK} consistently reduces to number of calls to the MIO solver by 50-85\%.}
 \label{tab:gridVsAlg}
	\setlength{\tabcolsep}{2pt}
	\begin{tabular}{l r r r|r r r | r r r  | r r r }
		
		\hline
		\multirow{2}{*}{Dataset}&\multirow{2}{*}{$p$}&\multirow{2}{*}{$n$}&\multirow{2}{*}{$\gamma$}& \multicolumn{3}{c|}{\underline{\texttt{Grid}}}&\multicolumn{3}{c|}{\underline{Algorithm~\ref{alg:parametricK}}}&\multicolumn{3}{c}{\textbf{\underline{Improvement}}}\\
		&&&&Time&\# MIO&Nodes&Time&\# MIO&Nodes&\textbf{Time}&\textbf{\# MIO}&\textbf{Nodes}\\
		\hline
			\multirow{7}{*}{Diabetes}&\multirow{7}{*}{11}&\multirow{7}{*}{442}&
		$0.01$&65&3,978&126,085&37&1,714&58,332&\textbf{45\%}&\textbf{56\%}&\textbf{53\%}\\
		&&&$0.02$&52&3,978&82,523&36&1,699&50,333&\textbf{30\%}&\textbf{56\%}&\textbf{37\%}\\
		&&&$0.05$&42&3,978&42,411&26&1,868&27,342&\textbf{29\%}&\textbf{52\%}&\textbf{35\%}\\
		&&&$0.10$&39&3,978&31,116&25&1,652&15,456&\textbf{34\%}&\textbf{53\%}&\textbf{48\%}\\
		&&&$0.20$&35&3,978&22,165&20&1,316&9,111&\textbf{42\%}&\textbf{67\%}&\textbf{58\%}\\
		&&&$0.50$&32&3,978&11,889&15&1,147&4,444&\textbf{50\%}&\textbf{71\%}&\textbf{59\%}\\
		&&&$1.00$&34&3,978&9,278&14&820&2,416&\textbf{58\%}&\textbf{79\%}&\textbf{73\%}\\
		&&&&&&&&&&&&\\
		\multirow{7}{*}{Housing}&\multirow{7}{*}{13}&\multirow{7}{*}{506}&
		$0.01$&247&6,072&512,723&91&1,867&216,411&\textbf{59\%}&\textbf{69\%}&\textbf{57\%}\\
		&&&$0.02$&187&6,072&324,238&64&1,711&139,293&\textbf{65\%}&\textbf{70\%}&\textbf{56\%}\\
		&&&$0.05$&166&6,072&216,116&87&1,679&91,822&\textbf{45\%}&\textbf{69\%}&\textbf{57\%}\\
		&&&$0.10$&40&6,072&96,387&18&1,814&40,112&\textbf{51\%}&\textbf{69\%}&\textbf{58\%}\\
		&&&$0.20$&82&6,072&68,581&34&1,599&24,899&\textbf{55\%}&\textbf{73\%}&\textbf{63\%}\\
		&&&$0.50$&90&6,072&60,067&34&1,233&20,231&\textbf{62\%}&\textbf{79\%}&\textbf{65\%}\\
		&&&$1.00$&107&6,072&49,770&22&947&13,111&\textbf{77\%}&\textbf{84\%}&\textbf{73\%}\\
		&&&&&&&&&&&&\\
		\multirow{7}{*}{Servo}&\multirow{7}{*}{19}&\multirow{7}{*}{167}&
		$0.01$&466&3,006&1,669,537&259&1,099&938,012&\textbf{41\%}&\textbf{60\%}&\textbf{44\%}\\
		&&&$0.02$&110&3,006&811,432&51&989&399,980&\textbf{52\%}&\textbf{66\%}&\textbf{51\%}\\
		&&&$0.05$&44&3,006&324,877&25&965&160,112&\textbf{77\%}&\textbf{84\%}&\textbf{73\%}\\
		&&&$0.10$&23&3,006&162,223&9&679&58,136&\textbf{59\%}&\textbf{77\%}&\textbf{64\%}\\
		&&&$0.20$&15&3,006&76,739&8&898&33,030&\textbf{48\%}&\textbf{70\%}&\textbf{57\%}\\
		&&&$0.50$&10&3,006&40,197&4&561&10,299&\textbf{56\%}&\textbf{81\%}&\textbf{74\%}\\
		&&&$1.00$&8&3,006&25,683&4&479&6,639&\textbf{52\%}&\textbf{84\%}&\textbf{74\%}\\
			&&&&&&&&&&&&\\
		\multirow{7}{*}{AutoMPG}&\multirow{7}{*}{25}&\multirow{7}{*}{392}&
		$0.01$&1,100&9,408&6,772,986&584&2,999&3,221,031&\textbf{46\%}&\textbf{67\%}&\textbf{48\%}\\
		&&&$0.02$&1,356&9,408&3,900,417&412&2,433&1,698,234&\textbf{67\%}&\textbf{70\%}&\textbf{52\%}\\
		&&&$0.05$&519&9,408&2,286,681&212&2,659&1,012,099&\textbf{56\%}&\textbf{70\%}&\textbf{50\%}\\
		&&&$0.10$&355&9,408&1,548,369&139&2,675&681,344&\textbf{59\%}&\textbf{71\%}&\textbf{56\%}\\
		&&&$0.20$&143&9,408&629,020&64&2,387&281,001&\textbf{54\%}&\textbf{71\%}&\textbf{55\%}\\
		&&&$0.50$&66&9,408&176,950&28&2,101&56,165&\textbf{58\%}&\textbf{76\%}&\textbf{67\%}\\
		&&&$1.00$&68&9,408&116,982&36&1,477&28,112&\textbf{43\%}&\textbf{84\%}&\textbf{74\%}\\
		\hline
	\end{tabular}
    \end{center}
\end{table}



	\begin{table}[h!]\footnotesize
    \begin{center}
 \caption{Comparison between using Algorithm~\ref{alg:parametricK} and solving $\mathcal{O}(kp)$ MIOs independently (\texttt{Grid}) for 10-fold cross validation in four real datasets, for different values of regularization $\gamma$. Times reported are in minutes, and correspond to the time to solve all required mixed-integer optimization problems to optimality. No time limits are imposed on the MIOs.}
 \label{tab:gridVsAlg10}
	\setlength{\tabcolsep}{2pt}
	\begin{tabular}{l r r r|r r r | r r r  | r r r }
		
		\hline
		\multirow{2}{*}{Dataset}&\multirow{2}{*}{$p$}&\multirow{2}{*}{$n$}&\multirow{2}{*}{$\gamma$}& \multicolumn{3}{c|}{\underline{\texttt{Grid}}}&\multicolumn{3}{c|}{\underline{Algorithm~\ref{alg:parametricK}}}&\multicolumn{3}{c}{\textbf{\underline{Improvement}}}\\
		&&&&Time&\# MIO&Nodes&Time&\# MIO&Nodes&\textbf{Time}&\textbf{\# MIO}&\textbf{Nodes}\\
		\hline
			\multirow{7}{*}{Diabetes}&\multirow{7}{*}{11}&\multirow{7}{*}{442}&
		$0.01$&3&396&11,666&2&242&8,224&\textbf{14\%}&\textbf{39\%}&\textbf{30\%}\\
		&&&$0.02$&2&396&8,371&2&235&6,785&\textbf{12\%}&\textbf{41\%}&\textbf{19\%}\\
		&&&$0.05$&2&396&4,436&2&228&3,430&\textbf{10\%}&\textbf{42\%}&\textbf{23\%}\\
		&&&$0.10$&2&396&3,185&2&247&2,277&\textbf{10\%}&\textbf{38\%}&\textbf{29\%}\\
		&&&$0.20$&1&396&2,268&1&206&1,536&\textbf{8\%}&\textbf{48\%}&\textbf{32\%}\\
		&&&$0.50$&1&396&1,233&1&149&643&\textbf{26\%}&\textbf{62\%}&\textbf{48\%}\\
		&&&$1.00$&1&396&872&1&93&287&\textbf{42\%}&\textbf{77\%}&\textbf{67\%}\\
		&&&&&&&&&&&&\\
		\multirow{7}{*}{Housing}&\multirow{7}{*}{13}&\multirow{7}{*}{506}&
		$0.01$&25&600&48,069&19&321&35,227&\textbf{25\%}&\textbf{47\%}&\textbf{27\%}\\
		&&&$0.02$&19&600&34,915&14&310&25,090&\textbf{28\%}&\textbf{48\%}&\textbf{28\%}\\
		&&&$0.05$&14&600&21,350&10&303&14,933&\textbf{29\%}&\textbf{50\%}&\textbf{30\%}\\
		&&&$0.10$&10&600&11,012&7&300&7,308&\textbf{31\%}&\textbf{50\%}&\textbf{34\%}\\
		&&&$0.20$&9&600&7,406&5&230&3,524&\textbf{46\%}&\textbf{62\%}&\textbf{52\%}\\
		&&&$0.50$&9&600&6,168&3&141&1,977&\textbf{62\%}&\textbf{77\%}&\textbf{68\%}\\
		&&&$1.00$&8&600&4,993&2&66&930&\textbf{77\%}&\textbf{89\%}&\textbf{81\%}\\
		&&&&&&&&&&&&\\
		\multirow{7}{*}{Servo}&\multirow{7}{*}{19}&\multirow{7}{*}{167}&
		$0.01$&15&288&148,168&12&191&128,592&\textbf{16\%}&\textbf{34\%}&\textbf{13\%}\\
		&&&$0.02$&8&288&77,457&7&190&67,416&\textbf{10\%}&\textbf{34\%}&\textbf{13\%}\\
		&&&$0.05$&3&288&29,056&3&157&23,653&\textbf{16\%}&\textbf{45\%}&\textbf{19\%}\\
		&&&$0.10$&2&288&15,951&2&146&12,562&\textbf{16\%}&\textbf{49\%}&\textbf{21\%}\\
		&&&$0.20$&1&288&8,117&1&155&6,275&\textbf{12\%}&\textbf{46\%}&\textbf{23\%}\\
		&&&$0.50$&1&288&4,028&1&201&2,922&\textbf{3\%}&\textbf{30\%}&\textbf{27\%}\\
		&&&$1.00$&1&288&2,541&1&206&1,768&\textbf{1\%}&\textbf{28\%}&\textbf{30\%}\\
			&&&&&&&&&&&&\\
		\multirow{7}{*}{AutoMPG}&\multirow{7}{*}{25}&\multirow{7}{*}{392}&
		$0.01$&111&936&691,816&76&389&460,187&\textbf{31\%}&\textbf{58\%}&\textbf{33\%}\\
		&&&$0.02$&68&936&401,905&44&374&264,179&\textbf{35\%}&\textbf{60\%}&\textbf{34\%}\\
		&&&$0.05$&42&936&225,318&30&396&161,639&\textbf{28\%}&\textbf{58\%}&\textbf{28\%}\\
		&&&$0.10$&30&936&149,243&20&389&98,261&\textbf{35\%}&\textbf{58\%}&\textbf{34\%}\\
		&&&$0.20$&14&936&61,534&10&389&41,323&\textbf{32\%}&\textbf{58\%}&\textbf{33\%}\\
		&&&$0.50$&7&936&17,865&4&318&8,550&\textbf{43\%}&\textbf{66\%}&\textbf{52\%}\\
		&&&$1.00$&6&936&10,848&3&251&4,480&\textbf{48\%}&\textbf{73\%}&\textbf{59\%}\\
		\hline
	\end{tabular}
    \end{center}
\end{table}

We observe that across these four datasets, Algorithm~\ref{alg:parametricK} reduces the number of MIO that need to be solved by an average of 70\% for leave-one-out cross-validation and by 52\% for 10-fold cross-validation. The overall number of branch-and-bound nodes is reduced by an average of 57\% for leave-one-out cross-validation and 35\% for 10-fold cross-validation (the reduction in computational times is similar to the reduction of nodes). 
We observe that solution times for both methods decrease on a given dataset as $\gamma$ increases (as expected, since the perspective reformulation is stronger). Interestingly, while the improvements of Algorithm~\ref{alg:parametricK} over \texttt{Grid} (in terms of time, MIOs solved and nodes) are more pronounced in regimes with large regularization $\gamma$, this effect on $\gamma$ is slight: Algorithm~\ref{alg:parametricK} consistently results in improvements over 40\% (and often more) even for the smallest values of $\gamma$ tested. These results indicate that the relaxations of the bilevel optimization \eqref{eq:paramK} derived in \S\ref{sec:lowerbounds} are sufficiently strong to avoid solving most of the MIOs that traditional methods such as \texttt{Grid} would solve, without sacrificing solution quality. The proposed methods are especially beneficial for settings where $k$ is large, that is, in the settings that would require more MIOs and are more computationally expensive using standard approaches. 

The resulting approach still requires solving several MIOs, but, as we show throughout the rest of this section, approximating each MIO with its perspective relaxation yields similarly high-quality statistical estimators at a fraction of the computational cost.




\subsection{Statistical results with real data}\label{sec:comp_real}
%Nonetheless, these metrics interest us since they allow us to measure the impact of accounting for stability on the number of features selected. 


%We consider two problem settings. First, a smaller-scale setting $(p=50, \tau_{\text{true}}=10)$ that allows us to benchmark two implementations of our coordinate descent approach: \begin{enumerate}
    %\item An exact implementation of our approach, where we optimize $\tau$ according to Algorithm \ref{alg:parametricK}, using \verb|Gurobi| version $9.5.1$ to solve all MIOs with a time limit of $120$s, and warm-start \verb|Gurobi| with a greedily rounded solution to each MIO's perspective relaxation (computed using \verb|Mosek| version $10.0$) before running \verb|Gurobi|. We denote this approach by ``EX'' (stands for EXact).
    %\item An approximate implementation of our approach, where we optimize $\tau$ by greedily rounding the perspective relaxation of each MIO we encounter (computed using \verb|Mosek| version $10.0$), and using these greedily rounded solutions, rather than optimal solutions to MIOs, to optimize the leave-one-out error with respect to $\tau$. We denote this approach by ``GD'' (stands for GreeDy).
%\end{enumerate}

In experiments with real data, solving MIOs to optimality or using grid search to identify a globally optimal solution with respect to the kCV loss is not practical due to excessive computational requirements. Accordingly, instead of solving each MIO subproblem to optimality, we use the saddle-point method of \cite{bertsimas2020sparse2} with default parameters to obtain high-quality solutions quickly. We term this implementation of our coordinate descent approach ``SP'' (stands for Saddle Point), and set $\tau_{\min}=2, \tau_{\max}$ such that $\tau_{\max}\log \tau_{\max}\leq \min(n,p)$ when optimizing $\tau$ in this experiment\footnote{We previously tried setting $\tau_{\max}=p$. We found that this yielded the same optimal hyperparameters, but increased the total runtime of the method substantially.} as in \cite{gamarnik2022sparse}

We compare against the following state-of-the-art methods, using in-built functions to approximately minimize the cross-validation loss with respect to the method's hyperparameters via grid search, and subsequently fit a regression model on the entire dataset with these cross-validated parameters (see also \cite{bertsimas2020sparse2} for a detailed discussion of these approaches):
\begin{itemize}
    \item The \verb|ElasticNet| method in the ubiquitous \verb|GLMNet| package, with grid search on their parameter $\alpha \in \{0, 0.1, 0.2, \ldots, 1\}$% using $100$-fold cross-validation as in \citep{bertsimas2020sparse2}. 
    \item The Minimax Concave Penalty (MCP) %and Smoothly Clipped Absolute Deviation Penalty (SCAD) 
    as implemented in the \verb|R| package \verb|ncvreg|, using the \verb|cv.ncvreg| function with %$100$ folds and 
    default parameters to (approximately) minimize the cross-validation error.
    %\item The \verb|L0Learn.cvfit| method implemented in the \verb|L0Learn| \verb|R| package \citep[cf.][]{hazimeh2020fast}, with $n$ folds, a grid of $10$ different values of $\gamma$ and default parameters otherwise. 
\end{itemize}
We compare performance in terms of the Mean Square Error, namely
\begin{align*}
    MSE(\bm{\beta}):=\frac{1}{n}\sum_{i=1}^n (y_{i}-\bm{x}_{i}^\top \bm{\beta})^2,
\end{align*}
which can either be taken over the validation set (CV)-- that is, the objective \eqref{prob:upperlevelofv_validation} we attempt to minimize-- or over an unseen test set (MSE), acting as a proxy for generalization error. 

To measure the validation and test set errors, we repeat the following procedure five times and report the average result: we randomly shuffle the data into $80\%$ training/validation and $20\%$ test data, perform five-fold cross-validation on the $80\%$ training/validation data, fit a model with the cross-validated $(\tau, \gamma)$ on the combined $80\%$ train/validation data, and evaluate the model's test-set performance on the remaining $20\%$ test data. We also report the average value of $\tau$, the cross-validated sparsity, for each method.

We observe in Table \ref{tab:comparison_ourmethods5} that SP obtains a lower five-fold cross-validation error (CV) for all datasets, being on average $9.5\%$ lower than GLMNet and $37\%$ lower than MCP. This result indicates that the proposed approach is significantly better at obtaining high-quality solutions in terms of the cross-validation error, even without solving the MIO subproblems to optimality. 

However, this improved validation error does not necessarily translate into an improved test set error, due to the optimizer's curse \citep[cf.][]{smith2006optimizer, cawley2010over}, where the cross-validation error consistently underestimates the out-of-sample error of the estimators. In overdetermined regimes, this error is relatively small, and achieving an improved cross-validation error indeed leads to better estimators out-of-sample: SP performs $2.0\%$ better out-of-sample than GLMNet and $3.3\%$ better out-of-sample than MCP. However, in underdetermined settings, the kCV error of SP underestimates the out-of-sample MSE by a factor of two. Thus, direct minimization of the cross-validation error leads to overfitting and worse estimators out-of-sample, and GLMNet performs $24\%$ better out-of-sample MSE despite a greater cross-validation error (although SP remains $1.2\%$ better out-of-sample than MCP). Finally, we note that the average runtime across all datasets was $1334$s for SP, $0.213$s for MCP, and $1.79$s for GLMNet.

\begin{table}[h]
\centering\footnotesize
\begin{tabular}{@{}l r r r r r r r r r r r @{}} \toprule
Dataset & n & p & \multicolumn{3}{c@{\hspace{0mm}}}{SP} &  \multicolumn{3}{c@{\hspace{0mm}}}{MCP} & \multicolumn{3}{c@{\hspace{0mm}}}{GLMNet}\\
\cmidrule(l){4-6} \cmidrule(l){7-9} \cmidrule(l){10-12}  &   &   & $\tau$ & CV & MSE&$\tau$ & CV & MSE &$\tau$ & CV & MSE \\\midrule
Wine & 6497 & 11 & 9.8 & 0.544 & 0.543  & 10.8 & 0.543 & 0.543 & 11 & 0.542 & 0.543\\
Auto-MPG & 392 & 25 & 17.4 & 	8.550 & 	8.871 & 	 	16 & 	9.154 & 	8.993 & 	20.8 & 	8.600 & 	8.979 \\ 
Hitters & 263 & 19 & 10.2 & 	0.076 & 	0.077 & 	13 & 	0.081 & 	0.081 & 	16 & 	0.076 & 	0.079 \\ 
Prostate & 97 & 8 & 4.4 & 	0.522 & 	0.567 & 	5.8 & 	0.572 & 	0.574 & 	6.4 & 	0.524 & 	0.561 \\ 
Servo & 167 & 19 & 12 & 	0.732 & 	0.725 & 13.8 & 	0.752 & 	0.722 & 	16.4 & 	0.685 & 	0.715 \\ 
Housing2 & 506 & 91 & 66 & 	11.76 & 	11.36 & 	36.8 & 	16.03 & 	14.89 & 	66 & 	12.68 & 	13.16 \\ 
Toxicity & 38 & 9 & 3.2 & 	0.037 & 	0.060 & 2.6 & 	0.049 & 	0.057 & 	5 & 	0.037 & 	0.061 \\ 
Steam & 25 & 8 & 2.2 & 	0.401 & 	0.479 & 2.2 & 	0.583 & 	0.559 & 	4.6 & 	0.502 & 	0.495 \\ 
Alcohol2 & 44 & 21 & 2.6 & 	0.229 & 	0.256 & 	2 & 	0.228 & 	0.232 & 	3.8 & 	0.237 & 	0.255 \\ \midrule
TopGear & 242 & 373 & 26.2 & 	0.037 & 	0.061 & 		7.4 & 	0.054 & 	0.073 & 	29.8 & 	0.044 & 	0.056 \\ 
Bardet & 120 & 200 & 21.8 & 	0.007 & 	0.011 & 	6 & 	0.008 & 	0.010 & 	30.2 & 	0.007 & 	0.009 \\ 
Vessel & 180 & 486 & 23.2 & 	0.014 & 	0.031 & 	2.8 & 	0.033 & 	0.033 & 	49.6 & 	0.018 & 	0.023 \\ 
Riboflavin & 71 & 4088 & 9.6 & 	0.149 & 	0.364 & 	8 & 	0.349 & 	0.352 & 	105.6 & 	0.212 & 	0.285 \\ 
\bottomrule
\end{tabular}
\caption{Average performance of five-fold version of methods across a suite of real-world datasets where the ground truth is unknown (and may not be sparse), sorted by how overdetermined the dataset is ($n/p$), and separated into the underdetermined and overdetermined cases. In both cases, the coordinate descent approach yields better five-fold cross-validation error than alternative approaches. In overdetermined settings, cyclic coordinate descent returns sparser solutions than MCP or GLMNet and achieves a better out-of-sample MSE on average. In underdetermined settings, GLMNet results in the better out-of-sample MSE, suggesting that optimization of the cross-validation error may lead to overfitting in this regime.
}
\label{tab:comparison_ourmethods5}
\end{table}

In a follow-up work \citep{cory2025stability}, we investigate techniques for improving the test-set performance of cross-validation more broadly and SP specifically, by optimizing both the stability of the regressors and their cross-validation error. This is because stable regressors are less likely to disappoint significantly out-of-sample, as identified by \cite{bousquet2002stability}. This involves separate ideas from the generalization theory literature, and thus is beyond the scope of the present work.

\section{Conclusion}
In this paper, we propose a new optimization-based approach 
for cross-validating hyperparameters in ridge-regularized sparse regression problems, by leveraging perspective relaxations and bounds on the cross-validation error. The proposed approach is effective at substantially decreasing the number of MIOs and branch-and-bound nodes explored required to optimize the cross-validation loss. In practice, when combined with state-of-the-art methods to find high-quality solutions to MIO problems quickly, we show that the proposed method is able to improve upon statistical packages in the literature, particularly in overdetermined settings where the cross-validation error is an accurate proxy of the out-of-sample error.  



 \subsubsection*{Acknowledgments: } Andr\'es G\'omez is supported in part by grantFA9550-24-1-0086 from the Air Force Office of Scientific Research. Ryan Cory-Wright gratefully acknowledges the MIT-IBM Research Lab for hosting him while part of this work was conducted. We are grateful to Jean Pauphilet for valuable discussions on an earlier draft of this manuscript.
{\def\enotesize{\footnotesize}

\theendnotes
}

\FloatBarrier
{
%\bibliographystyle{abbrvnat} 
% \setlength{\bibsep}{0pt plus 0.0ex}
\scriptsize
\bibliographystyle{informs2014}
\begin{thebibliography}{61}
    \providecommand{\natexlab}[1]{#1}
    \providecommand{\url}[1]{\texttt{#1}}
    \providecommand{\urlprefix}{URL }
    
    \bibitem[{Arlot \protect\BIBand{} Celisse(2010)}]{arlot2010survey}
    Arlot S, Celisse A (2010) A survey of cross-validation procedures for model
      selection. \emph{Statistics surveys} 4:40--79.
    
    \bibitem[{Askari et~al.(2022)Askari, d'Aspremont, \protect\BIBand{}
      Ghaoui}]{askari2022approximation}
    Askari A, d'Aspremont A, Ghaoui LE (2022) Approximation bounds for sparse
      programs. \emph{SIAM Journal on Mathematics of Data Science} 4(2):514--530.
    
    \bibitem[{Atamt{\"u}rk \protect\BIBand{} G{\'o}mez(2020)}]{atamturk2020safe}
    Atamt{\"u}rk A, G{\'o}mez A (2020) Safe screening rules for l0-regression from
      perspective relaxations. \emph{ICML}, 421--430.
    
    \bibitem[{Atamt{\"u}rk \protect\BIBand{} G{\'o}mez(2025)}]{atamturk2019rank}
    Atamt{\"u}rk A, G{\'o}mez A (2025) Rank-one convexification for sparse
      regression. \emph{Journal of Machine Learning Research, conditionally
      accepted} .
    
    \bibitem[{Beck \protect\BIBand{} Schmidt(2021)}]{beck2021gentle}
    Beck Y, Schmidt M (2021) A gentle and incomplete introduction to bilevel
      optimization .
    
    \bibitem[{Ben-Ayed \protect\BIBand{} Blair(1990)}]{ben1990computational}
    Ben-Ayed O, Blair CE (1990) Computational difficulties of bilevel linear
      programming. \emph{Operations Research} 38(3):556--560.
    
    \bibitem[{Bennett et~al.(2006)Bennett, Hu, Ji, Kunapuli, \protect\BIBand{}
      Pang}]{bennett2006model}
    Bennett KP, Hu J, Ji X, Kunapuli G, Pang JS (2006) Model selection via bilevel
      optimization. \emph{The 2006 IEEE International Joint Conference on Neural
      Network Proceedings}, 1922--1929 (IEEE).
    
    \bibitem[{Bergstra \protect\BIBand{} Bengio(2012)}]{bergstra2012random}
    Bergstra J, Bengio Y (2012) Random search for hyper-parameter optimization.
      \emph{Journal of Machine Learning Research} 13(2).
    
    \bibitem[{Bertsekas(2016)}]{bertsekas2016nonlinear}
    Bertsekas D (2016) \emph{Nonlinear Programming} (Athena Scientific), 3rd
      edition.
    
    \bibitem[{Bertsimas \protect\BIBand{}
      Copenhaver(2018)}]{bertsimas2018characterization}
    Bertsimas D, Copenhaver MS (2018) Characterization of the equivalence of
      robustification and regularization in linear and matrix regression.
      \emph{European Journal of Operational Research} 270(3):931--942.
    
    \bibitem[{Bertsimas \protect\BIBand{}
      Cory-Wright(2022)}]{bertsimas2022scalable}
    Bertsimas D, Cory-Wright R (2022) A scalable algorithm for sparse portfolio
      selection. \emph{INFORMS Journal on Computing} 34(3):1489--1511.
    
    \bibitem[{Bertsimas et~al.(2021)Bertsimas, Cory-Wright, \protect\BIBand{}
      Pauphilet}]{bertsimas2021unified}
    Bertsimas D, Cory-Wright R, Pauphilet J (2021) A unified approach to
      mixed-integer optimization problems with logical constraints. \emph{SIAM
      Journal on Optimization} 31(3):2340--2367.
    
    \bibitem[{Bertsimas et~al.(2020)Bertsimas, Pauphilet, \protect\BIBand{}
      Van~Parys}]{bertsimas2020sparse2}
    Bertsimas D, Pauphilet J, Van~Parys B (2020) Sparse regression: Scalable
      algorithms and empirical performance. \emph{Statistical Science}
      35(4):555--578.
    
    \bibitem[{Bertsimas \protect\BIBand{} Van~Parys(2020)}]{bertsimas2020sparse}
    Bertsimas D, Van~Parys B (2020) Sparse high-dimensional regression: Exact
      scalable algorithms and phase transitions. \emph{The Annals of Statistics}
      48(1):300--323.
    
    \bibitem[{Boland et~al.(2015{\natexlab{a}})Boland, Charkhgard,
      \protect\BIBand{} Savelsbergh}]{boland2015criterion}
    Boland N, Charkhgard H, Savelsbergh M (2015{\natexlab{a}}) A criterion space
      search algorithm for biobjective integer programming: The balanced box
      method. \emph{INFORMS Journal on Computing} 27(4):735--754.
    
    \bibitem[{Boland et~al.(2015{\natexlab{b}})Boland, Charkhgard,
      \protect\BIBand{} Savelsbergh}]{boland2015criterion2}
    Boland N, Charkhgard H, Savelsbergh M (2015{\natexlab{b}}) A criterion space
      search algorithm for biobjective mixed integer programming: The triangle
      splitting method. \emph{INFORMS Journal on Computing} 27(4):597--618.
    
    \bibitem[{Bousquet \protect\BIBand{} Elisseeff(2002)}]{bousquet2002stability}
    Bousquet O, Elisseeff A (2002) Stability and generalization. \emph{The Journal
      of Machine Learning Research} 2:499--526.
    
    \bibitem[{Boyd et~al.(1994)Boyd, El~Ghaoui, Feron, \protect\BIBand{}
      Balakrishnan}]{boyd1994linear}
    Boyd S, El~Ghaoui L, Feron E, Balakrishnan V (1994) \emph{Linear matrix
      inequalities in system and control theory} (SIAM).
    
    \bibitem[{B{\"u}hlmann \protect\BIBand{} Van
      De~Geer(2011)}]{buhlmann2011statistics}
    B{\"u}hlmann P, Van De~Geer S (2011) \emph{Statistics for high-dimensional
      data: methods, theory and applications} (Springer Science \& Business Media).
    
    \bibitem[{Burman(1989)}]{burman1989comparative}
    Burman P (1989) A comparative study of ordinary cross-validation, v-fold
      cross-validation and the repeated learning-testing methods. \emph{Biometrika}
      76(3):503--514.
    
    \bibitem[{Cawley \protect\BIBand{} Talbot(2010)}]{cawley2010over}
    Cawley GC, Talbot NL (2010) On over-fitting in model selection and subsequent
      selection bias in performance evaluation. \emph{The Journal of Machine
      Learning Research} 11:2079--2107.
    
    \bibitem[{Ceria \protect\BIBand{} Soares(1999)}]{ceria1999convex}
    Ceria S, Soares J (1999) Convex programming for disjunctive convex
      optimization. \emph{Math. Prog.} 86:595--614.
    
    \bibitem[{Cifuentes \protect\BIBand{} Moitra(2022)}]{cifuentes2022polynomial}
    Cifuentes D, Moitra A (2022) Polynomial time guarantees for the burer-monteiro
      method. \emph{Advances in Neural Information Processing Systems}
      35:23923--23935.
    
    \bibitem[{Cory-Wright \protect\BIBand{} G{\'o}mez(2025)}]{cory2025stability}
    Cory-Wright R, G{\'o}mez A (2025) Stability regularization improves the
      performance of cross-validation. \emph{submitted} .
    
    \bibitem[{Falkner et~al.(2018)Falkner, Klein, \protect\BIBand{}
      Hutter}]{falkner2018bohb}
    Falkner S, Klein A, Hutter F (2018) Bohb: Robust and efficient hyperparameter
      optimization at scale. \emph{International Conference on Machine Learning},
      1437--1446 (PMLR).
    
    \bibitem[{Frazier(2018)}]{frazier2018tutorial}
    Frazier PI (2018) A tutorial on bayesian optimization. \emph{arXiv preprint
      arXiv:1807.02811} .
    
    \bibitem[{Gamarnik \protect\BIBand{} Zadik(2022)}]{gamarnik2022sparse}
    Gamarnik D, Zadik I (2022) Sparse high-dimensional linear regression.
      estimating squared error and a phase transition. \emph{The Annals of
      Statistics} 50(2):880--903.
    
    \bibitem[{Geoffrion(1972)}]{geoffrion1972generalized}
    Geoffrion AM (1972) Generalized {B}enders decomposition. \emph{Journal of
      Optimization Theory and Applications} 10(4):237--260.
    
    \bibitem[{G{\'o}mez \protect\BIBand{} Prokopyev(2021)}]{gomez2021mixed}
    G{\'o}mez A, Prokopyev OA (2021) A mixed-integer fractional optimization
      approach to best subset selection. \emph{INFORMS Journal on Computing}
      33(2):551--565.
    
    \bibitem[{Grimmett \protect\BIBand{} Stirzaker(2020)}]{grimmett2020probability}
    Grimmett G, Stirzaker D (2020) \emph{Probability and random processes} (Oxford
      university press).
    
    \bibitem[{G{\"u}nl{\"u}k \protect\BIBand{}
      Linderoth(2010)}]{gunluk2010perspective}
    G{\"u}nl{\"u}k O, Linderoth J (2010) Perspective reformulations of mixed
      integer nonlinear programs with indicator variables. \emph{Math. Prog.}
      124(1):183--205.
    
    \bibitem[{Hansen et~al.(1992)Hansen, Jaumard, \protect\BIBand{}
      Savard}]{hansen1992new}
    Hansen P, Jaumard B, Savard G (1992) New branch-and-bound rules for linear
      bilevel programming. \emph{SIAM Journal on scientific and Statistical
      Computing} 13(5):1194--1217.
    
    \bibitem[{Hastie et~al.(2009)Hastie, Tibshirani, Friedman, \protect\BIBand{}
      Friedman}]{hastie2009elements}
    Hastie T, Tibshirani R, Friedman JH, Friedman JH (2009) \emph{The elements of
      statistical learning: data mining, inference, and prediction}, volume~2
      (Springer).
    
    \bibitem[{Hastie et~al.(2020)Hastie, Tibshirani, \protect\BIBand{}
      Tibshirani}]{hastie2020best}
    Hastie T, Tibshirani R, Tibshirani R (2020) Best subset, forward stepwise or
      {L}asso? analysis and recommendations based on extensive comparisons.
      \emph{Statistical Science} 35(4):579--592.
    
    \bibitem[{Hazan \protect\BIBand{} Koren(2016)}]{hazan2016linear}
    Hazan E, Koren T (2016) A linear-time algorithm for trust region problems.
      \emph{Mathematical Programming} 158(1-2):363--381.
    
    \bibitem[{Hazimeh \protect\BIBand{} Mazumder(2020)}]{hazimeh2020fast}
    Hazimeh H, Mazumder R (2020) Fast best subset selection: Coordinate descent and
      local combinatorial optimization algorithms. \emph{Operations Research}
      68(5):1517--1537.
    
    \bibitem[{Hazimeh et~al.(2022)Hazimeh, Mazumder, \protect\BIBand{}
      Saab}]{hazimeh2021sparse}
    Hazimeh H, Mazumder R, Saab A (2022) Sparse regression at scale:
      Branch-and-bound rooted in first-order optimization. \emph{Mathematical
      Programming} 196(1):347--388.
    
    \bibitem[{Hotelling(1931)}]{hotelling1931generalization}
    Hotelling H (1931) The generalization of student's ratio. \emph{The Annals of
      Mathematical Statistics} 2(3):360--378.
    
    \bibitem[{Kenney et~al.(2021)Kenney, Chiaromonte, \protect\BIBand{}
      Felici}]{kenney2021mip}
    Kenney A, Chiaromonte F, Felici G (2021) {MIP}-boost: Efficient and effective l
      0 feature selection for linear regression. \emph{Journal of Computational and
      Graphical Statistics} 30(3):566--577.
    
    \bibitem[{Larochelle et~al.(2007)Larochelle, Erhan, Courville, Bergstra,
      \protect\BIBand{} Bengio}]{larochelle2007empirical}
    Larochelle H, Erhan D, Courville A, Bergstra J, Bengio Y (2007) An empirical
      evaluation of deep architectures on problems with many factors of variation.
      \emph{Proc. 24th Int. Conf. Mach. Learn.}, 473--480.
    
    \bibitem[{LeCun et~al.(2015)LeCun, Bengio, \protect\BIBand{}
      Hinton}]{lecun2015deep}
    LeCun Y, Bengio Y, Hinton G (2015) Deep learning. \emph{Nature}
      521(7553):436--444.
    
    \bibitem[{Liu et~al.(2023)Liu, Rosen, Zhong, \protect\BIBand{}
      Rudin}]{liu2023okridge}
    Liu J, Rosen S, Zhong C, Rudin C (2023) Okridge: Scalable optimal k-sparse
      ridge regression. \emph{Advances in neural information processing systems}
      36:41076--41258.
    
    \bibitem[{Lokman \protect\BIBand{} K{\"o}ksalan(2013)}]{lokman2013finding}
    Lokman B, K{\"o}ksalan M (2013) Finding all nondominated points of
      multi-objective integer programs. \emph{Journal of Global Optimization}
      57:347--365.
    
    \bibitem[{Mazumder et~al.(2011)Mazumder, Friedman, \protect\BIBand{}
      Hastie}]{mazumder2011sparsenet}
    Mazumder R, Friedman JH, Hastie T (2011) Sparsenet: Coordinate descent with
      nonconvex penalties. \emph{Journal of the American Statistical Association}
      106(495):1125--1138.
    
    \bibitem[{Mazumder et~al.(2023)Mazumder, Radchenko, \protect\BIBand{}
      Dedieu}]{mazumder2023subset}
    Mazumder R, Radchenko P, Dedieu A (2023) Subset selection with shrinkage:
      Sparse linear modeling when the snr is low. \emph{Operations Research}
      71(1):129--147.
    
    \bibitem[{McAfee et~al.(2012)McAfee, Brynjolfsson, Davenport, Patil,
      \protect\BIBand{} Barton}]{mcafee2012big}
    McAfee A, Brynjolfsson E, Davenport TH, Patil D, Barton D (2012) Big data: The
      management revolution. \emph{Harvard Business Review} 90(10):60--68.
    
    \bibitem[{Natarajan(1995)}]{natarajan1995sparse}
    Natarajan BK (1995) Sparse approximate solutions to linear systems. \emph{SIAM
      journal on computing} 24(2):227--234.
    
    \bibitem[{Okuno et~al.(2021)Okuno, Takeda, Kawana, \protect\BIBand{}
      Watanabe}]{okuno2021lp}
    Okuno T, Takeda A, Kawana A, Watanabe M (2021) On lp-hyperparameter learning
      via bilevel nonsmooth optimization. \emph{J. Mach. Learn. Res.} 22:245--1.
    
    \bibitem[{Pilanci et~al.(2015)Pilanci, Wainwright, \protect\BIBand{}
      El~Ghaoui}]{pilanci2015sparse}
    Pilanci M, Wainwright MJ, El~Ghaoui L (2015) Sparse learning via boolean
      relaxations. \emph{Mathematical Programming} 151(1):63--87.
    
    \bibitem[{Reeves et~al.(2019)Reeves, Xu, \protect\BIBand{}
      Zadik}]{reeves2019all}
    Reeves G, Xu J, Zadik I (2019) The all-or-nothing phenomenon in sparse linear
      regression. \emph{Conference on Learning Theory}, 2652--2663 (PMLR).
    
    \bibitem[{Sidford(2024)}]{sidford2024notes}
    Sidford A (2024) Optimization algorithms. \emph{Lecture notes for Introduction
      to Optimization Theory and Optimization Algorithms} Stanford University.
    
    \bibitem[{Sinha et~al.(2020)Sinha, Khandait, \protect\BIBand{}
      Mohanty}]{sinha2020gradient}
    Sinha A, Khandait T, Mohanty R (2020) A gradient-based bilevel optimization
      approach for tuning hyperparameters in machine learning. \emph{arXiv preprint
      arXiv:2007.11022} .
    
    \bibitem[{Smith \protect\BIBand{} Winkler(2006)}]{smith2006optimizer}
    Smith JE, Winkler RL (2006) The optimizer’s curse: Skepticism and
      postdecision surprise in decision analysis. \emph{Management Science}
      52(3):311--322.
    
    \bibitem[{Stephenson et~al.(2021)Stephenson, Frangella, Udell,
      \protect\BIBand{} Broderick}]{stephenson2021can}
    Stephenson W, Frangella Z, Udell M, Broderick T (2021) Can we globally optimize
      cross-validation loss? quasiconvexity in ridge regression. \emph{Advances in
      Neural Information Processing Systems} 34.
    
    \bibitem[{Stidsen et~al.(2014)Stidsen, Andersen, \protect\BIBand{}
      Dammann}]{stidsen2014branch}
    Stidsen T, Andersen KA, Dammann B (2014) A branch and bound algorithm for a
      class of biobjective mixed integer programs. \emph{Management Science}
      60(4):1009--1032.
    
    \bibitem[{Stone(1978)}]{stone1978cross}
    Stone M (1978) Cross-validation: A review. \emph{Statistics: A Journal of
      Theoretical and Applied Statistics} 9(1):127--139.
    
    \bibitem[{Stubbs \protect\BIBand{} Mehrotra(1999)}]{stubbs1999branch}
    Stubbs RA, Mehrotra S (1999) A branch-and-cut method for 0-1 mixed convex
      programming. \emph{Mathematical programming} 86:515--532.
    
    \bibitem[{Takano \protect\BIBand{} Miyashiro(2020)}]{takano2020best}
    Takano Y, Miyashiro R (2020) Best subset selection via cross-validation
      criterion. \emph{Top} 28(2):475--488.
    
    \bibitem[{Xie \protect\BIBand{} Deng(2020)}]{xie2020scalable}
    Xie W, Deng X (2020) Scalable algorithms for the sparse ridge regression.
      \emph{SIAM Journal on Optimization} 30(4):3359--3386.
    
    \bibitem[{Xu et~al.(2008)Xu, Caramanis, \protect\BIBand{}
      Mannor}]{xu2008robust}
    Xu H, Caramanis C, Mannor S (2008) Robust regression and lasso. \emph{Advances
      in neural information processing systems} 21.
    
    \bibitem[{Ye et~al.(2022)Ye, Yuan, Zeng, \protect\BIBand{}
      Zhang}]{ye2022difference}
    Ye JJ, Yuan X, Zeng S, Zhang J (2022) Difference of convex algorithms for
      bilevel programs with applications in hyperparameter selection.
      \emph{Mathematical Programming} 1--34.
    
    \end{thebibliography}
    
}

\newpage
\ECSwitch
\ECHead{Supplementary Material}



% \begin{table}[h]
% \centering\footnotesize
% \begin{tabular}{@{}l r r r r r r r r r r r r r r r r r @{}} \toprule
% Dataset & n & p & \multicolumn{3}{c@{\hspace{0mm}}}{$\delta=10^{-0.5}$} & \multicolumn{3}{c@{\hspace{0mm}}}{$\delta=10^{-1}$} & \multicolumn{3}{c@{\hspace{0mm}}}{$\delta=10^{-1.5}$} & \multicolumn{3}{c@{\hspace{0mm}}}{$\delta=10^{-2}$} \\
% \cmidrule(l){4-6} \cmidrule(l){7-9} \cmidrule(l){10-12} \cmidrule(l){13-15} \cmidrule(l){16-18} &   &   & $\tau$ & CV & MSE & $\tau$ & CV & MSE & $\tau$ & CV & MSE & $\tau$ & CV & MSE \\\midrule
% Wine &	6497 &	11 &	1.8 & 	0.603 & 	0.724 & 	3.4 & 	0.5 & 	0.687 & 	9.6 & 	0.46 & 	0.573 & 	9.6 & 	0.443 & 	0.573 \\ 
% Auto-MPG &	392 &	25 &	13.4 & 	20.471 & 	11.228 & 	17.4 & 	11.228 & 	9.083 & 	17.8 & 	8.231 & 	8.926 & 	17.8 & 	7.259 & 	8.926 \\ 
% Hitters &	263 &	19 &	15.6 & 	0.098 & 	0.072 & 	14.8 & 	0.072 & 	0.079 & 	14.2 & 	0.064 & 	0.079 & 	12 & 	0.061 & 	0.079 \\ 
% Prostate &	97 &	8 &	2.4 & 	0.878 & 	0.572 & 	4.6 & 	0.572 & 	0.615 & 	4.8 & 	0.469 & 	0.613 & 	4.8 & 	0.434 & 	0.613 \\ 
% Servo &	167 &	19 &	11.2 & 	1.833 & 	0.995 & 	15 & 	0.995 & 	0.719 & 	13.6 & 	0.723 & 	0.725 & 	13.6 & 	0.634 & 	0.725 \\ 
% Housing2 &	506 &	91 &	72.4 & 	33.742 & 	17.312 & 	65.2 & 	17.312 & 	11.9 & 	66.4 & 	11.901 & 	11.705 & 	62.8 & 	10.164 & 	11.772 \\ 
% Toxicity &	38 &	9 &	3.4 & 	0.094 & 	0.051 & 	3.4 & 	0.051 & 	0.061 & 	3.4 & 	0.037 & 	0.061 & 	3.4 & 	0.032 & 	0.061 \\ 
% Steam &	25 &	8 &	3.4 & 	1.021 & 	0.554 & 	3.4 & 	0.554 & 	0.591 & 	2.2 & 	0.386 & 	0.487 & 	2.2 & 	0.341 & 	0.487 \\ 
% Alcohol2 &	44 &	21 &	2.6 & 	3.637 & 	1.271 & 	2.4 & 	1.271 & 	0.236 & 	2.6 & 	0.524 & 	0.241 & 	2.6 & 	0.285 & 	0.232 \\ \midrule
% TopGear &	242 &	373 &	26.2 & 	0.161 & 	0.068 & 	33.4 & 	0.068 & 	0.053 & 	35.6 & 	0.036 & 	0.060 & 	39.6 & 	0.026 & 	0.062 \\
% Bardet &	120 &	200 &	20.6 & 	0.030 & 	0.014 & 	21.2 & 	0.014 & 	0.010 & 	23 & 	0.008 & 	0.010 & 	20 & 	0.006 & 	0.011 \\ 
% Vessel &	180 &	486 &	27.6 & 	0.575 & 	0.191 & 	23.2 & 	0.191 & 	0.030 & 	20.6 & 	0.070 & 	0.030 & 	24.2 & 	0.029 & 	0.031 \\ 
% Riboflavin &	71 &	4088 &	9.4 & 	0.538 & 	0.253 & 	14 & 	0.253 & 	0.341 & 	12.4 & 	0.172 & 	0.368 & 	12.4 & 	0.144 & 	0.368 \\ 
% \bottomrule
% \end{tabular}
% \caption{Performance of five-fold version of methods across real-world datasets where the ground truth is unknown (continued).}
% \label{tab:comparison_ourmethods6}
% \end{table}


% \begin{table}[h]
% \centering\footnotesize
% \begin{tabular}{@{}l r r r r r r r r r r r r r r r r r r r r@{}} \toprule
% Dataset & n & p & \multicolumn{3}{c@{\hspace{0mm}}}{$\delta=10^{-2.5}$} & \multicolumn{3}{c@{\hspace{0mm}}}{$\delta=10^{-3}$} & \multicolumn{3}{c@{\hspace{0mm}}}{$\delta=10^{-3.5}$} & \multicolumn{3}{c@{\hspace{0mm}}}{$\delta$ calibrated} \\
% \cmidrule(l){4-6} \cmidrule(l){7-9} \cmidrule(l){10-12} \cmidrule(l){13-15}  &   &   & $\tau$ & CV & MSE & $\tau$ & CV & MSE & $\tau$ & CV & MSE & $\delta$ & $\tau$ & MSE\\\midrule
% Wine &	6497 &	11 &	9.6 & 	0.438 & 	0.573 & 	9.8 & 	0.436 & 	0.568 & 	9.8 & 	0.436 & 	0.568 & $10^{-3.5}$ & 9.8 & 0.568 \\
% Auto-MPG &	392 &	25 &	17.4 & 	6.961 & 	8.910 & 	17.4 & 	6.854 & 	8.939 & 	17.4 & 	6.822 & 	8.939 & $10^{-2.5}$ & 17.4 & 8.910\\
% Hitters &	263 &	19 &	12 & 	0.061 & 	0.079 & 	12 & 	0.060 & 	0.079 & 	11.4 & 	0.061 & 	0.080 &  $10^{-2.0}$ & 12 & 0.079\\
% Prostate &	97 &	8 &	3.8 & 	0.418 & 	0.607 & 	3.8 & 	0.414 & 	0.607 & 	3.8 & 	0.413 & 	0.607 & $10^{-2.5}$ & 3.8 & 0.607 \\
% Servo & 167 &	19 &	13 & 	0.606 & 	0.724 & 	13 & 	0.606 & 	0.724 & 	13 & 	0.597 & 	0.724 &  $10^{-3.5}$ & 13 & 0.724\\
% Housing2 &	506 &	91 &	64.4 & 	9.670 & 	11.821 & 	64.4 & 	9.496 & 	11.821 & 	65.8 & 	9.414 & 	11.687 & $10^{-2.0}$ & 62.8 & 11.772\\
% Toxicity &	38 &	9 &	3.4 & 	0.031 & 	0.061 & 	3.4 & 	0.030 & 	0.061 & 	3.4 & 	0.030 & 	0.061 & $10^{-3.5}$ & 3.4 & 0.061\\
% Steam &	25 &	8 &	2.2 & 	0.327 & 	0.487 & 	2.2 & 	0.323 & 	0.487 & 	2.2 & 	0.321 & 	0.487 & $10^{-1.5}$ & 2.2 & 0.487\\
% Alcohol2 &	44 &	21 &	2.6 & 	0.211 & 	0.232 & 	2.6 & 	0.187 & 	0.232 & 	2.6 & 	0.179 & 	0.232 & $10^{-3.0}$ & 2.6 & 0.232\\ \midrule
% TopGear &	242 &	373 &	43.6 & 	0.021 & 	0.070 & 	41.2 & 	0.021 & 	0.064 & 	41.2 & 	0.021 & 	0.069 & $10^{-1.5}$ & 35.6 & 0.060\\
% Bardet &	120 &	200 &	20 & 	0.006 & 	0.011 & 	20 & 	0.006 & 	0.011 & 	20 & 	0.006 & 	0.011 & $10^{-2.0}$ & 20 & 0.011\\
% Vessel &	180 &	486 &	24.6 & 	0.017 & 	0.029 & 	24.6 & 	0.013 & 	0.030 & 	24.6 & 	0.012 & 	0.030 & $10^{-3.0}$ & 24.6 & 0.030\\
% Riboflavin &	71 &	4088 &	11.6 & 	0.134 & 	0.383 & 	11.6 & 	0.131 & 	0.383 & 	13.4 & 	0.125 & 	0.370 & $10^{-1.5}$ & 12.4 & 0.368\\
% \bottomrule
% \end{tabular}
% \caption{Performance of five-fold version of methods across real-world datasets where the ground truth is unknown (continued).}
% \label{tab:comparison_ourmethods7}
% \end{table}

% \FloatBarrier






\section{Implementation Details}\label{sec:append.implementationdetails}
To solve each MIO in Algorithm \ref{alg:parametricK}, we invoke a Generalized Benders Decomposition scheme \citep{geoffrion1972generalized}, which was specialized to sparse regression problems by \cite{bertsimas2020sparse}. For any fixed $\gamma, \tau$, the method proceeds by minimizing a piecewise linear approximation of 
\begin{align}
    f(\bm{z}, \gamma):=\min_{\bm{\beta} \in \mathbb{R}^p: \ \Vert \bm{\beta}\Vert_0 \leq \tau} \ \frac{\gamma}{2}\sum_{j \in [p]}\frac{\beta_j^2}{z_j}+\Vert \bm{X}^{}\bm{\beta}-\bm{y}^{}\Vert_2^2,\label{prob:persp1}
\end{align}
% which, by strong duality, can be rewritten as
% \begin{align}\label{prob:saddlepoint}
%     h(\bm{z}, \gamma):=\max_{\bm{\alpha} \in \mathbb{R}^{n-1}} \ -\frac{1}{2\gamma}\sum_{j \in [p]}z_j \bm{\alpha}^{\top}\bm{X}_j^{(i)}\bm{X}_j^{(i)\top}\bm{\alpha}^{}-\frac{1}{2}\bm{\alpha}^{\top} \bm{\alpha}^{}+\bm{y}^{(i)\top}\bm{\alpha}^{}+\Vert \bm{y}^{(i)}\Vert_2^2
% \end{align}
until it either converges to an optimal solution or encounters a time limit. %Although this method requires exponentially many iterations to converge in the worst-case, much recent evidence in sparse regression \citep[][]{bertsimas2020sparse, bertsimas2020sparse2} %and more generally \citep[][]{dey2021branch} 
% suggests that this occurs with very low probability in practice. %Rather, this method often converges in a practical amount of time even with $100,000$s of features, although this scalability depends on the problem data and can be questioned in some circumstances, see \cite{atamturk2019rank}. % for a discussion of sparse regression's tractability and how it depends on the diagonal dominance of the sample covariance matrix). 
% Therefore, by running this method on each fold and at each sparsity budget (possibly with a time limit), we can minimize the LOOCV error with respect to $\tau$. 

We now discuss two enhancements that improve this method's performance in practice.

\paragraph{Warm-Starts:} as noted by \cite{ bertsimas2021unified}, a greedily rounded solution to the Boolean relaxation constitutes an excellent warm-start for a Generalized Benders Decomposition scheme. Therefore, when computing the lower and upper bounds on $h_{\mathcal{N}_j}(\gamma, \tau)$ for each $\tau$ by solving a perspective relaxation, we save the greedily rounded solution to the relaxation in memory, and provide the relevant rounding as a high-quality warm-start before solving the corresponding MIO. 

\paragraph{Screening Rules:} as observed by \cite{atamturk2020safe}, if we have an upper bound on the optimal value of $f(\bm{z}, \gamma)$, say $\bar{f}$, 
an optimal solution to the Boolean relaxation of minimizing \eqref{prob:persp1} over $\bm{z} \in [0, 1]^p$, say $(\bm{\beta}, \bm{z})$, and a lower bound on the optimal value of $h(\bm{z}, \gamma)$ from the Boolean relaxation, say $\underaccent{\bar}{f}$
then, letting $\beta_{[\tau]}$ be the $\tau$th largest value of $\beta$ in absolute magnitude, we have the following screening rules:
    \begin{itemize}
        \item If $\beta_i^2 \leq \beta_{[\tau+1]}^2$ and $\underaccent{\bar}{f}-\frac{1}{2\gamma}(\beta^2_{i}-\beta^2_{[\tau]})>\bar{f}$ then $z_i=0$.
        \item If $\beta_i^2 \geq \beta_{[\tau]}^2$ and $\underaccent{\bar}{f}+\frac{1}{2\gamma}(\beta_{i}^2-\beta^2_{[\tau+1]})>\bar{f}$ then $z_i=1$.
    \end{itemize}
Accordingly, to reduce the dimensionality of our problems, we solve a perspective relaxation for each fold of the data with $\tau=\tau_{\max}$ as a preprocessing step, and screen out the features where $z_i=0$ at $\tau=\tau_{\max}$ (for this fold of the data) before running Generalized Benders Decomposition.

\subsection{Implementation Details for Section \ref{ssec:parametric2}}\label{append.4.2}
In our numerical experiments, we find local minimizers of our approximation of $g$ by invoking the \verb|ForwardDiff| function in \verb|Julia| to automatically differentiate our approximation of $g$, and subsequently identify local minima via the \verb|Order0| method in the \verb|Roots.jl| package, which is designed to be a robust root-finding method. To avoid convergence to a low-quality local minimum, we run the search algorithm initialized at the previous iterate $\gamma_{t-1}$ and seven points log-uniformly distributed in $[10^{-3}, 10^{1}]$, and set $\gamma_t$ to be the local minima with the smallest estimated error. Moreover, to ensure numerical robustness, we require that $\gamma_t$ remains within the bounds $[10^{-3}, 10^1]$ and project $\gamma_t$ onto this interval if it exceeds these bounds (this almost never occurs in practice, because the data is preprocessed to be standardized). This approach tends to be very efficient in practice, particularly when the optimal support does not vary significantly as we vary $\gamma$. 

\end{document}
\section{Datasets}\label{sec:datasets}

We now describe the datasets we use to test the methods proposed in this paper, and competing alternatives in the literature. We use both synthetically generated data and real data in our experiments. This is because synthetic data allows us to control the ground truth and measure the accuracy of our methods in statistical settings, while real data allows us to measure the performance of our methods on datasets that arise in practice, and ensure that any performance gains with respect to out-of-sample MSE are not an artifice of the data generation process. 

\subsection{Synthetic datasets}\label{sec:synt_generation} We follow the experimental setup in \cite{bertsimas2020sparse2}. Given a fixed number of features $p$, number of datapoints $n$, true sparsity $1\leq \tau_{\text{true}}\leq p$, autocorrelation parameter $0\leq \rho\leq 1$ and signal to noise parameter $\nu$:
\begin{enumerate}
\item The rows of the model matrix are generated iid from a $p$-dimensional multivariate Gaussian distribution $\mathcal{N}(\bm{0},\bm{\Sigma})$, where $\Sigma_{ij}=\rho^{|i-j|}$ for all $i,j\in [p]$.
\item A ``ground-truth" vector $\bm{\beta}_{\text{true}}$ is sampled with exactly $\tau_{\text{true}}$ non-zero coefficients. The position of the non-zero entries is randomly chosen from a uniform distribution, and the value of the non-zero entries is either $1$ or $-1$ with equal probability.
\item The response vector is generated as $\bm{y}=\bm{X\beta_{\text{true}}}+\bm{\varepsilon}$, where each $\varepsilon_i$ is generated iid from a scaled normal distribution such that $\sqrt{\nu}=\|\bm{X\beta_{\text{true}}}\|_2/\|\bm{\varepsilon}\|_2$. 
\item We standardize $\bm{X}, \bm{y}$ to normalize and center them.
\end{enumerate}










\section{Numerical Experiments}\label{sec:numres} 

We now present numerical experiments testing our proposed methods. First, in Section~\ref{sec:compExactLOOCV}, we study the computational savings of using Algorithm~\ref{alg:parametricK} over a complete grid search when optimizing the $k$-fold error as a function of the sparsity parameter $\tau$. Next, in Sections \ref{sec:syntStats1}, we use synthetic data to benchmark the statistical performance of the proposed method against alternatives in the literature, and observe that our approach leads to significantly lower validation errors, although this sometimes translates to higher test errors. 
%Finally, in Section \ref{sec:syntStats2}, we combine our proposed approach with a confidence adjustment procedure proposed in a separate work \citep{cory2025stability} and demonstrate that, in conjunction with the confidence adjustment procedure, our approach also leads to lower test set errors than via other methods in the literature.

For conciseness, we describe the datasets used throughout the section in Appendix \ref{sec:datasets}.

{\color{black}\paragraph{Evaluation Metrics:} We now discuss the evaluation metrics that we use throughout this section and are standard in subset selection \citep[e.g.,][]{bertsimas2020sparse2, hastie2020best}. 

Suppose that our data observations $(\bm{x}_i, y_i)$ are generated according to some stochastic process via $y_i=\bm{x}_i^\top \bm{\beta}+\epsilon_i$, where $\epsilon_i$ is zero-mean noise and $\bm{\beta}_{\text{true}}$ is some fixed but unknown ground truth regressor. Then, we assess the statistical performance of various methods in terms of their accuracy:
\begin{align*}
    A(\bm{\beta}):=\frac{\Vert \bm{\beta}_{\text{true}} \circ \bm{\beta}\Vert_0}{\Vert \bm{\beta}_{\text{true}}\Vert_0},
\end{align*}
i.e., the proportion of true features that are selected, and their false discovery rate
\begin{align*}
    FDR(\bm{\beta}):=\frac{\vert \{j: \beta_j\neq 0, \beta_{\text{true},j}=0\}\vert}{\Vert \bm{\beta}_{\text{true}}\Vert_0},
\end{align*}
i.e., the proportion of selected features not included in the true support. 

It is worth noting that these metrics rely on knowledge of the ground truth $\bm{\beta}_{\text{true}}$, and thus cannot be applied when the ground truth is unknown. %Nonetheless, these metrics interest us since they allow us to measure the impact of accounting for stability on the number of features selected. 
Accordingly, we also compare performance in terms of the Mean Square Error, namely
\begin{align*}
    MSE(\bm{\beta}):=\frac{1}{n}\sum_{i=1}^n (y_{i}-\bm{x}_{i}^\top \bm{\beta})^2,
\end{align*}
which can either be taken over a training set to obtain the in-sample Mean Square Error, or over a test set to obtain the out-of-sample Mean Square Error.

%The methods developed here require computing $M$, an upper bound on the loss. Accordingly, we approximate $M$ by $\max_{i \in [n]} y_i^2$, the largest $y_i^2$ over the training set, throughout this section.
}



\subsection{Exact K-fold Optimization}\label{sec:compExactLOOCV}

We first assess whether Algorithm~\ref{alg:parametricK} significantly reduces the number of MIOs that need to be solved to minimize the kCV error with respect to $\tau$, compared to grid search. We set either $k=n$ or $k=10$, corresponding to leave-one-out and ten-fold cross-validation problems \eqref{eq:paramK} respectively.

We compare the performance of two approaches. First, a standard grid search approach (\texttt{Grid}), where we solve the inner MIO in \eqref{eq:paramK} for all combinations of cardinality $\tau\in [p]$ and all folds of the data $j\in [k]$, and select the hyperparameter combination which minimizes the objective. To ensure the quality of the resulting solution, we solve all MIOs to optimality (without any time limit). Second, we consider using Algorithm~\ref{alg:parametricK} with parameter $r=\infty$ (thus solving MIOs to optimality until the desired optimality gap $\epsilon$ for problem~\eqref{eq:paramK} is proven). We test regularization parameter $\gamma\in \{0.01,0.02,0.05,0.10,0.20,0.50,1.00\}$ in Algorithm~\ref{alg:parametricK}, and solve all MIOs via their perspective reformulations, namely
\begin{align*}
	\min_{\bm{\beta}\in \mathbb{R}^p,\bm{z}\in \{0,1\}^p}\;& \|\bm{X}\bm{\beta}-\bm{y}\|_2^2+\frac{\gamma}{2}\sum_{j=1}^p \frac{\beta_j^2}{z_j}\ \text{s.t.}\ \sum_{j=1}^p z_j\leq \tau,
\end{align*} 
using Mosek 10.0. Since the approach \texttt{Grid} involves solving $\mathcal{O}(kp)$ MIOs (without a time limit), we are limited to testing these approaches on small datasets, and accordingly use the Diabetes, Housing, Servo, and AutoMPG datasets for this experiment, as described by \cite{gomez2021mixed}. Moreover, we remark that the specific solution times and the number of nodes expanded by each method are not crucial, as those could vary substantially if relaxations other than the perspective are used, different solvers or solution approaches are used, or if advanced techniques are implemented (but both methods would be affected in the same way). Thus, we focus our analysis on relative performance.

We now summarize our experimental results and defer the details to Tables~\ref{tab:gridVsAlg} and \ref{tab:gridVsAlg10} of Appendix~\ref{sec:compTable}. Figures~\ref{fig:reductionMioNodes} and~\ref{fig:reductionMioNodes10} summarize the percentage reduction of the number of MIOs and the number of branch-and-bound nodes achieved by Algorithm~\ref{alg:parametricK} over \texttt{Grid}, computed as 
\small$$ \text{Reduction in MIOs}=\frac{\text{\# MIO}_{\texttt{Grid}}-\text{\# MIO}_{\text{Alg.\ref{alg:parametricK}}}}{\text{\# MIO}_{\texttt{Grid}}},\;\; \text{Reduction in nodes}=\frac{\text{\# nodes}_{\texttt{Grid}}-\text{\# nodes}_{\text{Alg.\ref{alg:parametricK}}}}{\text{\# nodes}_{\texttt{Grid}}},$$\normalsize
where $\text{\# MIO}_{\text{Y}}$ and $\text{\# nodes}_{\text{Y}}$ indicate the number of MIOs or branch-and-bound nodes used by method Y.  

\begin{figure}[h!]\centering
\begin{subfigure}[t]{.45\linewidth}
\includegraphics[width=\textwidth,trim={12cm 6cm 12cm 6cm},clip]{mio_reduction.pdf}
\end{subfigure}
\begin{subfigure}[t]{.45\linewidth}
\includegraphics[width=\textwidth,trim={12cm 6cm 12cm 6cm},clip]{nodes_reduction.pdf}
\end{subfigure}
\caption{Reduction in the number of MIO solved (left) and the total number of branch-and-bound nodes (right) when using Algorithm~\ref{alg:parametricK} for leave-one-out cross-validation, when compared with \texttt{Grid} (i.e., independently solving $\mathcal{O}(pn)$ MIOs) in four real datasets. The distributions shown in the figure correspond to solving the same instance with different values of $\gamma$. All MIOs are solved to optimality, without imposing any time limits. } % Old version-running another version at the moment which will be larger-scale, since this version only halved the number of MIOs that needed solving.
\label{fig:reductionMioNodes}
\end{figure}

\begin{figure}[h!]\centering
\begin{subfigure}[t]{.45\linewidth}
\includegraphics[width=\textwidth,trim={12cm 6cm 12cm 6cm},clip]{mio_reduction10.pdf}
\end{subfigure}
\begin{subfigure}[t]{.45\linewidth}
\includegraphics[width=\textwidth,trim={12cm 6cm 12cm 6cm},clip]{nodes_reduction10.pdf}
\end{subfigure}
\caption{Reduction in the number of MIO solved (left) and the total number of branch-and-bound nodes (right) when using Algorithm~\ref{alg:parametricK} for 10-fold cross-validation, when compared with \texttt{Grid} (i.e., independently solving $\mathcal{O}(pk)$ MIOs) in four real datasets. The distributions shown in the figure correspond to solving the same instance with different values of $\gamma$. All MIOs are solved to optimality, without imposing any time limits. } % Old version-running another version at the moment which will be larger-scale, since this version only halved the number of MIOs that needed solving.
\label{fig:reductionMioNodes10}
\end{figure}

We observe that across these four datasets, Algorithm~\ref{alg:parametricK} reduces the number of MIO that need to be solved by an average of 70\% for leave-one-out cross-validation and by 52\% for 10-fold cross-validation. The overall number of branch-and-bound nodes is reduced by an average of 57\% for leave-one-out cross-validation and 35\% for 10-fold cross-validation (the reduction in computational times is similar to the reduction of nodes). These results indicate that the relaxations of the bilevel optimization \eqref{eq:paramK} derived in \S\ref{sec:lowerbounds} are sufficiently strong to avoid solving most of the MIOs that traditional methods such as \texttt{Grid} would solve, without sacrificing solution quality. The proposed methods are especially beneficial for settings where $k$ is large, that is, in the settings that would require more MIOs and are more computationally expensive using standard approaches. The resulting approach still requires solving several MIOs, but, as we show throughout the rest of this section, approximating each MIO with its perspective relaxation yields similarly high-quality statistical estimators at a fraction of the computational cost.

\subsection{Statistical results with real data}\label{sec:comp_real}

We observe that similarly to Section \ref{ssec:realworldata}, without any confidence adjustment, we obtain a five-fold cross-validation error $9.4\%$ lower than GLMNet and $37\%$ lower than MCP, respectively. Similarly, this does not translate to a lower test set error, due to the optimizer's curse. In particular, the average MSE without (with calibrated) confidence adjustment is $5.7\%$ ($6.2\%$) higher than GLMNet, and the average regressor is $32.8\%$ ($29.8\%$) sparser than the regressor generated by GLMNet. The MSE without (with calibrated) confidence adjustment is also $3\%$ ($2.5\%$) lower than MCP.

Interestingly, the impact of the confidence adjustment procedure on the MSE appears to be negligible in the five-fold case (the MSEs with and without confidence adjustment are statistically indistinguishable), likely because the five-fold error is a more biased and lower variance estimator of the test set error than the leave-one-out error in underdetermined settings \citep{hastie2009elements}. Indeed, the optimized five-fold error when $\delta=0$ for Riboflavin, the most underdetermined dataset, is nearly three times the optimized leave-one-out error when $\delta=0$. This suggests that the benefits of confidence adjustment may be most pronounced in the leave-one-out case where there is more variance, particularly in underdetermined settings with many folds. Indeed, over the four underdetermined datasets, the average MSE of models trained via a calibrated leave-one-out approach is $20\%$ lower than the average MSE of models trained via a calibrated five-fold approach.

\begin{table}[h]
\centering\footnotesize
\begin{tabular}{@{}l r r r r r r r r r r r r r r@{}} \toprule
Dataset & n & p & \multicolumn{3}{c@{\hspace{0mm}}}{$\delta=0.0$} & \multicolumn{3}{c@{\hspace{0mm}}}{$\delta=1.0$} & \multicolumn{3}{c@{\hspace{0mm}}}{MCP} & \multicolumn{3}{c@{\hspace{0mm}}}{GLMNet}\\
\cmidrule(l){4-6} \cmidrule(l){7-9} \cmidrule(l){10-12} \cmidrule(l){13-15}  &   &   & $\tau$ & CV & MSE & $\tau$ & CV & MSE &$\tau$ & CV & MSE &$\tau$ & CV & MSE \\\midrule
Wine & 6497 & 11 & 9.8 & 0.435 & 0.543 & 1.2 & 0.886 & 0.750 & 10.8 & 0.435 & 0.542 & 11 & 0.433 & 0.543\\
Auto-MPG & 392 & 25 & 17.4 & 	6.840 & 	8.871 & 	9 & 	47.855 & 	13.095 & 	16 & 	7.323 & 	8.993 & 	20.8 & 	6.880 & 	8.979 \\ 
Hitters & 263 & 19 & 10.2 & 	0.060 & 	0.077 & 	14.4 & 	0.178 & 	0.081 & 	13 & 	0.065 & 	0.081 & 	16 & 	0.061 & 	0.079 \\ 
Prostate & 97 & 8 & 4.4 & 	0.418 & 	0.567 & 	1.2 & 	1.842 & 	0.659 & 	5.8 & 	0.457 & 	0.574 & 	6.4 & 	0.420 & 	0.561 \\ 
Servo & 167 & 19 & 12 & 	0.585 & 	0.725 & 	6.2 & 	4.457 & 	0.810 & 	13.8 & 	0.601 & 	0.722 & 	16.4 & 	0.548 & 	0.715 \\ 
Housing2 & 506 & 91 & 66 & 	9.415 & 	11.356 & 	60.4 & 	84.76 & 	17.216 & 	36.8 & 	12.823 & 	14.895 & 	66 & 	10.142 & 	13.158 \\ 
Toxicity & 38 & 9 & 3.2 & 	0.029 & 	0.060 & 	3.6 & 	0.234 & 	0.061 & 	2.6 & 	0.039 & 	0.057 & 	5 & 	0.030 & 	0.061 \\ 
Steam & 25 & 8 & 2.2 & 	0.321 & 	0.479 & 	4 & 	2.629 & 	0.814 & 	2.2 & 	0.466 & 	0.559 & 	4.6 & 	0.401 & 	0.495 \\ 
Alcohol2 & 44 & 21 & 2.6 & 	0.183 & 	0.256 & 	5.6 & 	11.17 & 	0.267 & 	2 & 	0.182 & 	0.232 & 	3.8 & 	0.189 & 	0.255 \\ \midrule
TopGear & 242 & 373 & 26.2 & 	0.030 & 	0.061 & 	17.8 & 	0.444 & 	0.062 & 	7.4 & 	0.044 & 	0.073 & 	29.8 & 	0.035 & 	0.056 \\ 
Bardet & 120 & 200 & 21.8 & 	0.005 & 	0.011 & 	20 & 	0.083 & 	0.010 & 	6 & 	0.007 & 	0.010 & 	30.2 & 	0.006 & 	0.009 \\ 
Vessel & 180 & 486 & 23.2 & 	0.011 & 	0.031 & 	24.4 & 	1.798 & 	0.030 & 	2.8 & 	0.026 & 	0.033 & 	49.6 & 	0.014 & 	0.023 \\ 
Riboflavin & 71 & 4088 & 9.6 & 	0.120 & 	0.364 & 	16 & 	1.331 & 	0.408 & 	8 & 	0.280 & 	0.352 & 	105.6 & 	0.170 & 	0.285 \\ 
\bottomrule
\end{tabular}
\caption{Average performance of five-fold version of methods across a suite of real-world datasets where the ground truth is unknown (and may not be sparse), sorted by how overdetermined the dataset is ($n/p$), and separated into the underdetermined and overdetermined cases. In overdetermined settings, cyclic coordinate descent (without confidence) returns sparser solutions than MCP or GLMNet and maintains a comparable average MSE. In underdetermined settings, cyclic coordinate descent with confidence returns significantly sparse solutions than GLMNet with a comparable MSE, and more accurate (although denser) solutions than MCP.
}
\label{tab:comparison_ourmethods5}
\end{table}

\subsection{Sparse Regression on Synthetic Data}\label{sec:syntStats1}
We now benchmark our coordinate descent approach on synthetic sparse regression problems where the ground truth is known to be sparse, but the number of non-zeros is unknown. The goal of this section is to test whether implementing the methods described in this paper leads to improve validation and test error vs alternatives in the literature.

We consider two problem settings. First, a smaller-scale setting $(p=50, \tau_{\text{true}}=10)$ that allows us to benchmark two implementations of our coordinate descent approach: 
\begin{enumerate}
    \item An exact implementation of our approach, where we optimize $\tau$ according to Algorithm \ref{alg:parametricK}, using \verb|Gurobi| version $9.5.1$ to solve all MIOs with a time limit of $120$s, and warm-start \verb|Gurobi| with a greedily rounded solution to each MIO's perspective relaxation (computed using \verb|Mosek| version $10.0$) before running \verb|Gurobi|. We denote this approach by ``EX'' (stands for EXact).
    \item An approximate implementation of our approach, where we optimize $\tau$ by greedily rounding the perspective relaxation of each MIO we encounter (computed using \verb|Mosek| version $10.0$), and using these greedily rounded solutions, rather than optimal solutions to MIOs, to optimize the leave-one-out error with respect to $\tau$. We denote this approach by ``GD'' (stands for GreeDy).
\end{enumerate}
For both approaches, we optimize $\gamma$ as described in Section \ref{ssec:parametric2}, and set $\tau_{\min}=4, \tau_{\max}=20$.

We also consider a large-scale setting $(p=1000, \tau_{\text{true}}=20)$ where grid search is not sufficient to identify a globally optimal solution with respect to the kCV loss. In this setting, the subproblems are too numerically expensive to solve exactly, and accordingly, we optimize $\tau$ using an approach very similar to ``GD'', except to optimize $\tau$ we solve each subproblem using the saddle-point method of \cite{bertsimas2020sparse2} with default parameters, rather than greedily rounding the perspective relaxations of MIOs. This approach generates solutions that are almost identical to those generated by GD, but is more scalable. We term this implementation of our coordinate descent approach ``SP'' (stands for Saddle Point), and set $\tau_{\min}=10, \tau_{\max}=40$ when optimizing $\tau$ in this experiment. 

We compare against the following state-of-the-art methods, using in-built functions to approximately minimizes the cross-validation loss with respect to the method's hyperparameters via grid search, and subsequently fit a regression model on the entire dataset with these cross-validated parameters (see also \cite{bertsimas2020sparse2} for a detailed discussion of these approaches):
\begin{itemize}
    \item The \verb|ElasticNet| method in the ubiquitous \verb|GLMNet| package, with grid search on their parameter $\alpha \in \{0, 0.1, 0.2, \ldots, 1\}$, using $100$-fold cross-validation as in \citep{bertsimas2020sparse2}. 
    \item The Minimax Concave Penalty (MCP) and Smoothly Clipped Absolute Deviation Penalty (SCAD) as implemented in the \verb|R| package \verb|ncvreg|, using the \verb|cv.ncvreg| function with $100$ folds and default parameters to (approximately) minimize the cross-validation error.
    \item The \verb|L0Learn.cvfit| method implemented in the \verb|L0Learn| \verb|R| package \citep[cf.][]{hazimeh2020fast}, with $n$ folds, a grid of $10$ different values of $\gamma$ and default parameters otherwise. 
\end{itemize}

We remark that, in preliminary experiments, we found that using the \verb|cv.GLMNet| and \verb|cv.ncvreg| to minimize the kCV error when $k=n$ was orders-of-magnitude more expensive than other approaches. Accordingly, we settled with minimizing the $100$-fold cross-validation error as a surrogate.
%in preliminary experiments, we considered using the \verb|cv.GLMNet| and \verb|cv.ncvreg| functions with $n$ folds to approximately minimize the leave-one-out error for the \verb|ElasticNet|, \verb|SCAD| and \verb|MCP| penalties. However, these methods take several orders of magnitude more computational time than either \verb|L0Learn| or Algorithm 
%\ref{alg:sbb} {\color{red}[Missing ref]} to cross-validate with $n$ folds, which is computationally prohibitive. Therefore, we follow \cite{bertsimas2020sparse2} in minimizing the $100$-fold cross-validation error as a surrogate of the leave-one-out error for these methods.

%\paragraph{Data Generation Process: } We follow the experimental setup in \cite{bertsimas2020sparse2}, set $\tau_{\text{true}}=10, p=50$ for the first experiment (resp. $\tau_{\text{true}}=20, p=1000$ for the second one) and generate a $k_{\text{true}}$ sparse vector $\bm{\beta} \in \mathbb{R}^p$ with entries $\{0, \pm 1\}$. Next, we draw the columns of $\bm{X}_i \in \mathbb{R}^{p \times n} \sim \mathcal{N}(0,1)^{p \times n}$ from a multivariate normal distribution with zero mean and covariance matrix $\bm{\Sigma}$, where $\Sigma_{i,j}=\rho^{\vert i-j\vert}$ and $\rho=0.7$ (resp. $\rho=0.2$). %Further, we draw another $10,000$ observations $(\bm{X}_i, y_i)$ from the same stochastic process, which we reserve for a test set that evaluates the out-of-sample performance of each method. 
%We set the signal-to-noise ratio to be $\nu=1$ and standardize $\bm{X}, \bm{Y}$ to normalize and center them. Accordingly, our first experiment approximately corresponds to the ``medium noise, high correlation'' setting studied in \cite{bertsimas2020sparse2}, while the second experiment approximately corresponds to their ``medium noise, low correlation'' setting.% studied by \cite{bertsimas2020sparse2}.%, although we take $k_{\text{true}}, p$ to be smaller to compensate for the increase computational cost in approximately minimizing the leave-one-out error rather than using a hold-out set. %For each method, we identify the set of hyperparameters which minimizes the leave-one-out cross-validation error, and fit a model to the provided dataset at these parameter values. 


\paragraph{Experimental Methodology: } We measure each method's ability to recover the ground truth (true positive rate) while avoiding detecting irrelevant features (false discovery rate). We consider two sets of synthetic data, following \cite{bertsimas2020sparse2}: a small (medium noise, high correlation) dataset: $\tau_{\text{true}}=10$, $p=50$, $\rho=0.7$ and $\nu=1$; and a large (medium noise, low correlation) dataset: $\tau_{\text{true}}=20$, $p=1000$, $\rho=0.2$ and $\nu=1$. Figure \ref{fig:sparseregression_synthetic} reports results in small instances with varying number of samples $n \in \{10, 20, \ldots, 200\}$, and Figure \ref{fig:sparseregression_synthetic2} reports results for large datasets with  $n \in \{100, 200, \ldots, 3000\}$. We report the average leave-one-out error and average MSE on a different (out-of-sample) set of $10000$ observations of $\bm{X}, \bm{y}$ drawn from the same distribution. Further, we report the average cross-validated support size and runtime for both experiments in \S \ref{append.section5.2}.


\begin{figure}[h!]
    \centering
    \begin{subfigure}[t]{.45\linewidth}
    \centering
            \includegraphics[scale=0.35]{synthetic_regression_tpr_p50.pdf}
    \end{subfigure}
        \begin{subfigure}[t]{.45\linewidth}
            \includegraphics[scale=0.35]{fpr_synthetic_p50.pdf}
    \end{subfigure}\\
        \begin{subfigure}[t]{.45\linewidth}
            \includegraphics[scale=0.35]{l10_synthetic_p50.pdf}
    \end{subfigure}
    \begin{subfigure}[t]{.45\linewidth}
            \includegraphics[scale=0.35]{mses_synthetic_p50.pdf}
    \end{subfigure}\\
    % \begin{subfigure}[t]{.45\linewidth}
    %         \includegraphics[scale=0.35]{synthetic_regression_sparsity_p50.pdf}
    % \end{subfigure}
    %     \begin{subfigure}[t]{.45\linewidth}
    %         \includegraphics[scale=0.35]{ts_synthetic_p50.pdf}
    % \end{subfigure}\\
   \caption{Average accuracy (top left), false discovery rate (top right), normalized validation error (bottom left), normalized MSE on test set (bottom right), as $n$ increases with $p=50, \tau_{\text{true}}=10$, for coordinate descent with $\tau$ optimized using Algorithm \ref{alg:parametricK} (EX), coordinate descent with $\tau$ optimized by greedily rounding perspective relaxations (GD), GLMNet, MCP, SCAD, and L0Learn. We average results over $25$ datasets.} 
   \label{fig:sparseregression_synthetic}
\end{figure}

\begin{figure}[h!]
    \centering
    \begin{subfigure}[t]{.45\linewidth}
    \centering
            \includegraphics[scale=0.35]{synthetic_regression_tpr_p1000.pdf}
    \end{subfigure}
        \begin{subfigure}[t]{.45\linewidth}
            \includegraphics[scale=0.35]{fpr_synthetic_p1000.pdf}
    \end{subfigure}\\
        \begin{subfigure}[t]{.45\linewidth}
            \includegraphics[scale=0.35]{l10_synthetic_p1000.pdf}
    \end{subfigure}
    \begin{subfigure}[t]{.45\linewidth}
            \includegraphics[scale=0.35]{mses_synthetic_p1000.pdf}
    \end{subfigure}\\
    % \begin{subfigure}[t]{.45\linewidth}
    %         \includegraphics[scale=0.35]{synthetic_regression_sparsity_p1000.pdf}
    % \end{subfigure}
    %     \begin{subfigure}[t]{.45\linewidth}
    %         \includegraphics[scale=0.35]{ts_synthetic_p1000.pdf}
    % \end{subfigure}\\
   \caption{Average accuracy (top left), false discovery rate (top right), normalized validation error (bottom left), and normalized MSE on test set (bottom right) as $n$ increases with $p=1000, \tau_{\text{true}}=20$, for coordinate descent with a saddle-point method to solve each training problem when optimizing $\tau$ (SP), GLMNet, MCP, SCAD, and L0Learn. We average results over $25$ datasets.} 
   \label{fig:sparseregression_synthetic2}
\end{figure}


\paragraph{Accuracy and Performance of Methods: }We observe that our coordinate descent schemes and \verb|L0Learn| consistently provide the best performance in large-sample settings, by returning sparser solutions with a lower false discovery rate and a similar out-of-sample MSE to all other methods when $n>p$. On the other hand, \verb|GLMNet| appears to perform best when $n \ll p$, where it consistently returns solutions with a lower out-of-sample MSE and less out-of-sample disappointment than any other method. Thus, the best-performing methods vary depending on the number of samples, as recently suggested by a number of authors \citep{hastie2020best, bertsimas2020sparse}. In particular, coordinate descent performs worse than \verb|GLMNet| when $p=1000$ and $n<2000$ despite having a much lower kCV error, precisely because of the dangers of minimizing the kCV error without confidence adjustment, as discussed in the introduction. 

\paragraph{Out-of-Sample Disappointment: } We observe that all methods suffer from the optimizer's curse \citep[cf.][]{smith2006optimizer, van2021data}, with the average MSE on the test set being consistently larger than the average leave-one-out error on the validation set, especially when $n$ is smaller. However, out-of-sample disappointment is most pronounced for our coordinate descent schemes and \verb|L0Learn| (bottom two panels of Figures \ref{fig:sparseregression_synthetic}--\ref{fig:sparseregression_synthetic2}), which consistently exhibit the lowest LOOCV error at all sample sizes but the highest test set error in small-data settings. As explained in more detail in \cite{cory2025stability}, this phenomenon can be explained by the fact that \textit{optimizing} the kCV error may generate optimistic estimates of the corresponding test set error. 

%On the other hand, \verb|GLMNet| generates solutions that do not disappoint as significantly out-of-sample at any problem size. This occurs because \verb|GLMNet| is a robust method, not a sparse method \citep[cf.][]{bertsimas2018characterization}, and its increased robustness improves its stability. Unfortunately, it consistently makes more false discoveries than all other methods, with a false discovery rate of $50\%$--$80\%$. Thus, it is a viable choice in small-data settings where accuracy is unimportant but should not be used to elicit the ground truth. {[\color{red} I would probably delete this paragraph]}


\subsection{Confidence-Adjusted Sparse Regression on Synthetic Data}\label{sec:syntStats2}
We now benchmark our coordinate descent schemes with confidence adjustment. In particular, we revisit the problem settings studied in the previous section and consider setting $\delta \in \{1, \sqrt{10}\}$ in our GD and SP implementations of coordinate descent. Specifically, we solve each subproblem using greedy rounding when $p=50, \tau_{\text{true}}=10$ and via a saddle-point method when $p=1000, \tau_{\text{true}}=20$. For ease of comparison, we also report the performance of these inexact methods without any confidence adjustment, as reported in the previous section {\color{black}(here denoted by $\delta =0$)}.

\paragraph{Experimental Methodology: } We implement the same methodology as in the previous section, and vary $n \in \{10, 20, \ldots, 100\}$ for small instances (Figure \ref{fig:sparseregression_synthetic3}) and $n \in \{100, 200, \ldots, 1500\}$ for large instances (Figure \ref{fig:sparseregression_synthetic4}). We report the average accuracy, average false positive rate, average cross-validated support size and average MSE on a different (out-of-sample) set of $10000$ observations of $\bm{X}, \bm{y}$ drawn from the same distribution.%, but unobserved during the optimization process. 

\begin{figure}[h!]
    \centering
    \begin{subfigure}[t]{.45\linewidth}
    \centering
            \includegraphics[scale=0.35]{synthetic_regression_tpr_p50_chebyshev.pdf}
    \end{subfigure}
        \begin{subfigure}[t]{.45\linewidth}
            \includegraphics[scale=0.35]{fpr_synthetic_p50_chebyshev.pdf}
    \end{subfigure}\\
    \begin{subfigure}[t]{.45\linewidth}
            \includegraphics[scale=0.35]{synthetic_regression_sparsity_p50_chebyshev.pdf}
    \end{subfigure}
    \begin{subfigure}[t]{.45\linewidth}
            \includegraphics[scale=0.35]{mses_synthetic_p50_chebyshev.pdf}
    \end{subfigure}\\
   \caption{Average accuracy (top left), false discovery rate (top right), cross-validated support (bottom left), and normalized MSE on test set (bottom right) as $n$ increases with $p=50, \tau_{\text{true}}=10$, with $\delta \in \{1, \sqrt{10}\}$, and without any confidence adjustment ({\color{black}$\delta =0$}). We average results over $25$ datasets.} 
   \label{fig:sparseregression_synthetic3}
\end{figure}


\begin{figure}[h!]
    \centering
    \begin{subfigure}[t]{.45\linewidth}
    \centering
            \includegraphics[scale=0.35]{tps_p1000_chebyshev.pdf}
    \end{subfigure}
        \begin{subfigure}[t]{.45\linewidth}
            \includegraphics[scale=0.35]{fprs_p1000_chebyshev.pdf}
    \end{subfigure}\\
    \begin{subfigure}[t]{.45\linewidth}
            \includegraphics[scale=0.35]{synthetic_regression_sparsity_p1000_chebyshev.pdf}
    \end{subfigure}
        \begin{subfigure}[t]{.45\linewidth}
            \includegraphics[scale=0.35]{mses_synthetic_p1000_chebyshev.pdf}
    \end{subfigure}\\
   \caption{Average accuracy (top left), false discovery rate (top right), cross-validated support (bottom left), and normalized MSE on test set (bottom right) as $n$ increases with $p=1000, \tau_{\text{true}}=20$, for coordinate descent with confidence adjustment and $\delta \in \{1.0, \sqrt{10}\}$, and without any confidence adjustment ({\color{black} $\delta =0$}). We use a saddle-point method to approximately solve all sparse regression subproblems for all methods, and average results over $25$ datasets.} 
   \label{fig:sparseregression_synthetic4}
\end{figure}


\paragraph{Impact of Confidence Adjustment on Test Set Performance: } We observe that for both problem settings, accounting for out-of-sample disappointment via confidence adjustment significantly improves the test set performance of sparse regression methods with hyperparameters selected via cross-validation. On average, we observe a $3.04\%$ (resp. $2.55\%$) average MSE improvement for $\delta=1$ (resp. $\delta=\sqrt{10}$) when $p=50$, and a $5.48\%$ (resp. $2.75\%$ ) average MSE improvement for $\delta=1$ (resp. $\delta=\sqrt{10}$) when $p=1000$, with the most significant out-of-sample performance gains occurring when $n$ is smallest. This performance improvement highlights the benefits of accounting for out-of-sample disappointment by selecting more stable sparse regression models, and suggests that accounting for model stability when cross-validating is a viable alternative to selecting hyperparameters that minimize the leave-one-out error that often yields better test set performance. 

Additionally, when $p=1000$, accounting for model stability via confidence adjustment yields sparser regressors with a significantly lower false discovery rate ($5\%$ vs. $40\%$ when $n=1000$), which suggests that models selected via a confidence adjustment procedure may sometimes be less likely to select irrelevant features. However, we caution that when $p=50$, the models selected via a confidence-adjusted procedure exhibit a similar false discovery rate to models selected by minimizing the LOOCV error, so this effect does not appear to occur uniformly.

All-in-all, the best value of $\delta$ to select for a confidence adjustment procedure appears to depend on the amount of data available, with larger values like $\delta=\sqrt{10}$ performing better when $n$ is small, but being overly conservative when $n$ is large, and smaller values like $\delta=1$ providing a less significant benefit when $n$ is very small but performing more consistently when $n$ is large. 



\section{Conclusion}
In this paper, we propose a new optimization-based approach 
for cross-validating hyperparameters in ridge-regularized sparse regression problems, by leveraging perspective relaxations and branch-and-bound techniques. Our approach obtains locally optimal hyperparameter combinations with $p=1000$ features in a few hours, and thus is a viable hyperparameter selection technique in offline settings where sparse and stable regressors are desirable. 





